% 第18章：自动驾驶场景
\chapter{自动驾驶场景}
\label{ch:18}

\paragraph*{本章阅读锚点}
本章的锚点是自动驾驶场景。原始 3DGS 的静态重建不足以描述道路参与者；驾驶方法要把动态物体、传感器融合、场景编辑和仿真闭环接起来。


自动驾驶系统从原型到量产的跨越需要海量、多样、安全关键的测试数据，而真实路测代价高昂且长尾场景难以复现。基于仿真的闭环验证成为不可或缺的补充手段，但传统仿真器的视觉真实感与真实世界存在显著差距。3DGS 以其照片级真实感渲染和显式场景表示，为构建高保真自动驾驶数字孪生提供了兼具保真度和灵活性的新范式。

本章围绕 136 篇相关论文，从五个维度系统梳理 3DGS 在自动驾驶中的技术进展：第~\ref{sec:18_1}~节讨论街景重建与渲染（48 篇），涵盖从基础表示到前馈重建、从外观解耦到轨迹外推的完整技术栈；第~\ref{sec:18_2}~节介绍仿真与数据生成（38 篇），覆盖闭环仿真、场景编辑和多模态数据合成；第~\ref{sec:18_3}~节分析动态驾驶场景（25 篇），聚焦动静态分解、实例感知和运动建模；第~\ref{sec:18_4}~节探讨传感器融合与定位（13 篇），处理 LiDAR、雷达和相机等多传感器统一仿真；第~\ref{sec:18_5}~节介绍 3D 占用估计与感知（12 篇），将高斯表示引入下游感知任务。第~\ref{sec:18_summary}~节总结全章。

\begin{figure}[htbp]
\centering
\resizebox{\textwidth}{!}{\begin{tikzpicture}[
  route/.style={rectangle, rounded corners=5pt, draw=#1, fill=#1!10,
    text width=3.2cm, minimum height=0.9cm, align=center, font=\small},
  method/.style={rectangle, rounded corners=3pt, draw=#1, fill=#1!6,
    font=\scriptsize, text width=3.0cm, align=center, minimum height=0.6cm},
  arr/.style={-{Stealth[length=4pt]}, techgray!60},
  >=Stealth
]

  \node[route=techgray, text width=4cm, font=\small\bfseries] (root) at (18.2, 5.6) {自动驾驶场景方法分类};
  \node[route=techblue] (b0) at (5.1, 3.1) {街景重建与渲染};
  \node[method=techblue] (b0_l0) at (1.5, 0.6) {街景表示\\{\tiny StreetGaussians, DrivingGaussian, ADM-GS}};
  \node[method=techblue] (b0_l1) at (5.1, 0.6) {动态建模\\{\tiny ReconDrive, StreetForward, PointForward}};
  \node[route=techgreen] (b1) at (12.3, 3.1) {仿真与数据生成};
  \node[method=techgreen] (b1_l0) at (8.7, 0.6) {闭环仿真\\{\tiny SplatSim-Drive, DriveSim-GS}};
  \node[method=techgreen] (b1_l1) at (12.3, 0.6) {数据合成\\{\tiny DrivingForward, XYZCylinder, VGD}};
  \node[route=techorange] (b2) at (18.2, 3.1) {动态驾驶场景};
  \node[method=techorange] (b2_l0) at (15.9, 0.6) {动静分解\\{\tiny Street Gaussians, DynaGS-Drive}};
  \node[method=techorange] (b2_l1) at (19.5, 0.6) {综合动态\\{\tiny 4DGS-Drive, OmniRe, PVG}};
  \node[route=techpurple] (b3) at (23.6, 3.1) {传感器融合};
  \node[method=techpurple] (b3_l0) at (23.6, 0.6) {多传感器融合\\{\tiny LiDAR-GS-Drive, RadarGS}};
  \node[method=techpurple] (b3_l1) at (27.3, 0.6) {定位\\{\tiny GeomGS, LocGS, PoseGS}};
  \node[route=techblue] (b4) at (30.9, 3.1) {3D占用与感知};
  \node[method=techblue] (b4_l0) at (30.9, 0.6) {高斯辅助感知\\{\tiny OccGS, GaussianOcc, GS-Perc}};
  \node[method=techblue] (b4_l1) at (35.0, 0.6) {占用估计\\{\tiny OccSplat, GSOcc}};
  \draw[arr] (root) -- (b0);
  \draw[arr] (root) -- (b1);
  \draw[arr] (root) -- (b2);
  \draw[arr] (root) -- (b3);
  \draw[arr] (root) -- (b4);
  \draw[arr] (b0) -- (b0_l0);
  \draw[arr] (b0) -- (b0_l1);
  \draw[arr] (b1) -- (b1_l0);
  \draw[arr] (b1) -- (b1_l1);
  \draw[arr] (b2) -- (b2_l0);
  \draw[arr] (b2) -- (b2_l1);
  \draw[arr] (b3) -- (b3_l0);
  \draw[arr] (b3) -- (b3_l1);
  \draw[arr] (b4) -- (b4_l0);
  \draw[arr] (b4) -- (b4_l1);
\end{tikzpicture}}
\caption{自动驾驶场景方法分类}
\label{fig:ch18_driving_taxonomy}
\end{figure}

%=====================================================================
\section{街景重建与渲染}
\label{sec:18_1}

街景重建是 3DGS 在自动驾驶中应用最为广泛的子方向，其核心目标是从车载传感器数据构建可自由视角浏览的高保真三维场景。标准 3DGS 假设相机围绕中心物体拍摄，获得充足的多视角覆盖。但自动驾驶场景中，相机安装在移动车辆上，形成前向为主的窄视角分布，大部分场景区域仅在 1--2 个视角中可见。此外，城市街道背景广阔且弱纹理（如路面、天空），标准 Structure from Motion（SfM，运动恢复结构）初始化在这些区域往往失败；前景车辆和行人在训练视角间造成动态遮挡，需要与静态背景分离处理；不同时段和天气条件下采集的数据存在跨视角外观不一致。这些因素使得标准 3DGS 的直接应用效果受限，催生了大量针对性方法。

\begin{figure}[htbp]
\centering
\resizebox{\textwidth}{!}{\begin{tikzpicture}[
  box/.style={rectangle, rounded corners=4pt, draw=#1, fill=#1!10,
    text width=3.2cm, align=center, font=\small, minimum height=1cm},
  sbox/.style={rectangle, rounded corners=3pt, draw=#1!60, fill=#1!5,
    font=\scriptsize, text width=2.4cm, align=center, minimum height=0.5cm},
  arr/.style={-{Stealth[length=6pt]}, thick, techgray},
  darr/.style={-{Stealth[length=6pt]}, thick, techgray, dashed},
  >=Stealth
]

  \node[box=techblue] (d1) at (3.4, 0.6) {数据采集(多相机/LiDAR)};
  \node[box=techgreen] (d2) at (10.9, 0.6) {街景重建(静态背景)};
  \node[box=techorange] (d3) at (17.8, 0.6) {动态建模(车辆/行人)};
  \node[box=techpurple] (d4) at (24.7, 0.6) {仿真引擎(闭环测试)};
  \node[box=techblue] (d5) at (31.5, 0.6) {训练验证(感知模型)};
  \node[font=\scriptsize, text=techgray] at (10.9, -0.6) {天空/地面分离};
  \node[font=\scriptsize, text=techgray] at (17.8, -0.6) {实例级动态高斯};
  \node[font=\scriptsize, text=techgray] at (24.7, -0.6) {传感器模拟};
  \draw[arr] (d1) -- (d2);
  \draw[arr] (d2) -- (d3);
  \draw[arr] (d3) -- (d4);
  \draw[arr] (d4) -- (d5);
\end{tikzpicture}}
\caption{自动驾驶场景 pipeline：从数据采集到仿真引擎}
\label{fig:ch18_driving_pipeline}
\end{figure}

\subsection{基础街景表示与动态建模}

本节讨论5条技术路线：（1）基础街景表示与动态建模；（2）前馈与可泛化重建；（3）外观解耦与多轨迹重建；（4）轨迹外推与车道变换；（5）特殊场景与领域方法。

早期工作聚焦于如何用高斯基元有效表示包含动态元素的城市场景，奠定了该方向的技术基础。核心问题在于：标准 3DGS 假设场景完全静态，而城市街道中充斥着移动的车辆和行人。如何将动态物体与静态背景分离、如何在各自的坐标系中高效建模、以及如何将两者组合渲染成统一的高保真图像，成为这一方向的核心挑战。

\paragraph{动静分离的显式高斯街景表示}

Street Gaussians（2024.01）~\cite{2401_01339} 是该方向的里程碑工作，奠定了动态街景显式高斯表示的基本范式。

\textbf{Baseline 与问题。} 在 Street Gaussians 之前，动态街景重建主要依赖 NeRF 类隐式方法（如 UniSim、Mars），这些方法将场景建模为连续神经辐射场，训练和渲染速度慢，且难以对前景车辆进行单独编辑。3DGS 虽然在静态场景中实现了实时渲染，但无法直接处理包含移动车辆的动态街景。

\textbf{核心改动。} Street Gaussians 将动态城市街景表示为一组带语义 logits 的 3D 高斯点云，分为静态背景和 $N$ 个前景车辆两组独立建模。背景模型使用世界坐标系下的标准 3D 高斯（位置 $\boldsymbol{\mu}_b$、协方差 $\Sigma_b$、不透明度 $o_b$、球谐系数 $z_b$）。前景车辆的每个高斯定义在物体局部坐标系中，通过可优化的追踪位姿 $(R_t, T_t)$ 变换到世界坐标系：

\begin{equation*}
\boldsymbol{\mu}_w = R_t \boldsymbol{\mu}_o + T_t, \quad R_w = R_t R_o, \quad \Sigma_w = R_w S_o S_o^T R_w^T
\end{equation*}

其中下标 $w$ 表示世界坐标，$o$ 表示物体局部坐标。关键创新在于 4D 球谐函数模型：将每个 SH 系数 $z_{m,l}$ 替换为一组傅里叶变换系数 $\{f_i \in \mathbb{R}^k\}$，通过逆离散傅里叶变换恢复时间 $t$ 的外观：$z_{m,l}(t) = \sum_{i=0}^{k-1} f_i \cos(2\pi i t / T)$。直觉上，任何随时间周期性变化的外观（如阴影移动、反光闪烁）都可以用少数几个不同频率的余弦波叠加近似——正如音频信号可以用几个主频分量重建。这样以极少的额外参数编码了时间信息，避免了为每帧单独存储 SH 系数的高存储开销。

\textbf{渲染与初始化。} 渲染时，将背景和变换后的前景点云拼接后统一投影到 2D 图像空间进行 $\alpha$-混合。天空区域用高分辨率立方体贴图单独建模。初始化方面，用聚合 LiDAR 点云替代 SfM 点云以获得更好的几何先验，物体模型用 3D 边界框内的 LiDAR 点或随机采样点初始化。

\textbf{设计动机。} 显式点云表示使得前景与背景可自由组合、支持场景编辑操作；可学习的位姿残差补偿追踪器的噪声；4D 傅里叶 SH 以 $O(k)$ 参数（$k$ 为傅里叶阶数，通常取 3--4）替代 $O(T \times C_{SH})$ 参数，兼顾时间建模能力和存储效率。

\textbf{性能。} 在 KITTI 和 Waymo Open 数据集上持续超越当时所有方法，以 135 FPS（$1066\times1600$ 分辨率）实现渲染，训练仅需半小时。

\paragraph{渐进式复合高斯图与深度分层}

Street Gaussians 虽然奠定了动态街景显式高斯表示的基础，但其一次性建模策略在远距离场景容易产生尺度混淆。DrivingGaussian（2023.12）~\cite{2312_07920} 提出渐进式静态背景建模——按时间顺序逐步累加静态 3D 高斯——并引入复合动态高斯图处理多个运动物体，解决了这一局限。

\textbf{渐进式静态建模。} 直觉上，这如同画家从近到远逐层绘制风景——先画近处确定的物体建立锚点，再逐步扩展到远处模糊的背景。静态背景根据 LiDAR 深度范围均匀划分为 $N$ 个深度区间（bin），按时间顺序依次构建。第一个区间使用 LiDAR 先验初始化高斯参数（位置 $P(x,y,z)$、协方差 $\Sigma$、球谐颜色 $C(r,g,b)$、不透明度 $o$），后续区间以上一个区间的高斯位置为先验，基于重叠区域对齐。这种渐进策略避免了远距离场景一次性引入导致的尺度混淆和模糊伪影。

\textbf{复合动态高斯图。} 动态物体通过复合图 $H = \langle O, G_d, M, P, A, T \rangle$ 建模，其中每个节点存储一个实例物体 $o \in O$、对应动态高斯 $g_i \in G_d$、变换矩阵 $m_o \in M$、边界框中心和朝向 $(p_o, a_o)$。多个动态物体之间的遮挡通过基于距离的透明度调整处理：距离相机更近的物体获得更高的透明度，这符合光线传播原理。

\textbf{LiDAR 先验与训练。} 使用 LiDAR 点云替代 SfM 初始化，投影到多相机图像获取颜色，再通过密集束调整（Dense Bundle Adjustment, DBA）扩展到多相机设置来优化点位置。训练损失包含 Tile Structural Similarity（TSSIM）、鲁棒损失和 LiDAR 几何监督三项。在 Waymo 上实现高保真环绕视角合成。

\subsubsection{前馈与可泛化重建}

本路线介绍前馈（feed-forward）范式下从图像直接预测高斯参数实现秒级驾驶场景重建的方法，涵盖自监督深度驱动、柱面投影和基础模型蒸馏等技术。

逐场景优化在实际部署中不够可扩展——每遇到一个新场景就需要 30 分钟以上的训练，无法满足大规模自动驾驶测试中快速生成数百个场景的需求。为此，大量工作探索前馈（feed-forward）范式：训练一个通用模型，推理时直接从图像预测高斯参数，实现秒级场景重建。驾驶场景前馈重建的独特难点在于：车载相机重叠区域小，传统极线匹配几乎不可行；不同车辆搭载不同数量和类型的相机，模型需要灵活适配。

\paragraph{自监督深度驱动的环绕视角前馈重建}

DrivingForward（2024.09）~\cite{2409_12753} 是面向驾驶场景的前馈 3DGS 模型，核心突破在于实现了无需深度真值和相机外参的自监督高斯基元预测。

\textbf{Baseline 与问题。} 前馈 3DGS 方法（如 pixelSplat、MVSplat）依赖大重叠的配对视角通过极线几何预测高斯位置。但驾驶场景中车载相机间重叠极小、视角稀疏，极线匹配几乎不可行。

\textbf{核心改动——尺度感知定位。} DrivingForward 放弃多视角匹配范式，转而从单帧图像预测深度，通过三种光度损失学习尺度信息：（1）时间上下文——同一相机相邻帧间的光度重投影；（2）空间上下文——同帧相邻相机在重叠区域的光度约束；（3）时空上下文——相邻相机在不同帧间的联合约束。位姿网络 $P$ 预测车辆运动，深度网络 $D$ 估计深度图，高斯网络 $G$ 从融合特征预测其余参数（缩放 $s$、旋转 $r$、不透明度 $o$、颜色 $c$）。关键设计在于：深度编码器提取几何特征后与深度网络共享的图像特征融合，使高斯参数预测继承尺度感知能力。

\textbf{设计动机。} 利用固定相机到车体变换矩阵 $E_i$（而非时变的世界坐标外参），结合位姿网络预测的运动，训练时无需世界坐标外参。推理时每帧独立预测，因此可以支持灵活的单帧或多帧环绕视角输入。

\textbf{实现原理。} 深度网络 $D$ 从图像 $I_t$ 预测逐像素深度 $d \in \mathbb{R}^{H \times W}$，结合位姿网络 $P$ 预测的自车运动 $\Delta T_t$，通过相机到车体的固定变换 $E_i$ 将像素反投影为 3D 高斯中心：

\begin{equation*}
\boldsymbol{\mu} = K^{-1} [u, v, 1]^T \cdot d(u,v) \cdot R_i^{-1} + t_i
\end{equation*}

其中 $K$ 为相机内参，$R_i, t_i$ 为第 $i$ 个相机到车体的变换。高斯网络 $G$ 从融合的深度特征和图像特征预测其余参数（缩放 $s \in \mathbb{R}^3_+$、旋转四元数 $r \in \mathbb{R}^4$、不透明度 $o \in [0,1]$、球谐颜色 $c$）。尺度信息通过三类光度损失联合学习：时间光度重投影 $\mathcal{L}_{temp} = \|I_t - \hat{I}_{t \to t+1}\|$、空间重叠约束 $\mathcal{L}_{spatial}$ 和时空联合约束 $\mathcal{L}_{st}$，总损失为 $\mathcal{L} = \lambda_1 \mathcal{L}_{render} + \lambda_2 \mathcal{L}_{temp} + \lambda_3 \mathcal{L}_{spatial} + \lambda_4 \mathcal{L}_{st}$。

\textbf{性能。} 在 nuScenes 上超越当时的逐场景优化方法，支持不同数量环绕视角帧的一次性前馈重建。

\paragraph{统一柱面投影与体素级高斯生成}

XYZCylinder（2025.10）~\cite{2510_07856} 在 DrivingForward 的前馈框架基础上，针对其两个结构性瓶颈提出改进。

\textbf{锚点回顾与差异。} DrivingForward 通过隐式网络学习视角变换，产生逐像素深度的 ``2.5D'' 场景表示。XYZCylinder 放弃这一隐式范式，转而采用确定性的统一柱面相机建模（Unified Cylinder Camera Modeling, UCCM）策略：以中心点 $c_u^x$、半径 $R_u = Z_u / (2 \cdot \tan(f_{\min}/2))$、高度 $Z_u$ 五个参数定义柱面坐标系，将 $N$ 个相机视角显式投影到统一柱面平面上。这一设计无需学习视角对应关系，仅通过调整训练无关参数即可适配不同相机数量和外参内参，从而实现零样本跨相机配置泛化。

\textbf{柱面平面特征组（CPFG）。} DrivingForward 的逐像素深度表示在面对 360\textdegree{} 稀疏视角时会产生变形、空洞和畸变。XYZCylinder 为此设计了柱面平面特征组（Cylinder Plane Feature Group, CPFG），通过三个专用模块将 2D 图像特征提升到 3D 空间：（1）占用感知模块（YNet$_\text{occ}$）区分占用体素与空体素，剪枝空区域以节省显存；（2）体积感知模块（XNet$_\text{vol}$）为占用体素生成 $G_v$ 个 3D 高斯，以几何特征 $F^\text{geo}$ 和外观特征 $F^\text{app}$ 双分支解耦预测高斯位置偏移 $\Delta x_{i^g}$ 与其余参数（球谐系数 $c_{i^g}$、协方差 $\Sigma_{i^g}$、不透明度 $o_{i^g}$）；（3）像素感知模块（YNet$_\text{pix}$）以 $4\times$ 上采样分辨率生成额外高斯，精修纹理细节。

\textbf{前背景解耦。} DrivingForward 未区分前景与背景，3D 高斯可能被误置于远景区域形成浮动伪影。XYZCylinder 将前景（车辆、道路、建筑）用 3DGS 稀疏重建，背景（天空、云层）由 ZNet$_\text{bg}$ 在 2D 空间以 $4H_t \times 4W_t$ 高分辨率生成后下采样，最终通过 $\alpha$-混合 融合：$I = I_\text{fg} + (1 - \hat{\alpha}_\text{fg}) \cdot I_\text{bg}'$。

\textbf{设计动机。} DrivingForward 的隐式视角变换对特定相机配置过拟合，换车即失效；XYZCylinder 用确定性映射替代可学习的视角对应，在 Waymo、Argoverse、Pandaset、ONCE 等未见过数据集上均可直接推理。此外，体素级高斯生成替代逐像素深度预测，避免了 ``2.5D'' 表示在 360\textdegree{} 全景中因尺度歧义导致的变形和空洞。

\textbf{性能。} 在 nuScenes 和 Carla-Centric 两个基准上，XYZCylinder 在 PSNR、SSIM、LPIPS 和深度 PCC 四项指标上均取得最优。相比 DrivingForward 等像素方法在 nuScenes 上 PSNR 相当但几何精度（PCC）显著领先，且无 DrivingForward 渲染图像中常见的白色空洞伪影。仅在 nuScenes 上训练的模型可直接在 Waymo（不同相机数量）和 Argoverse 上零样本重建。

\paragraph{几何基础模型蒸馏的前馈重建}

VGD（2025.10）~\cite{2510_19578} 与 DrivingForward 一样面向前馈环绕视角驾驶场景重建，但核心区别在于几何信息的获取方式：DrivingForward 通过自监督光度损失隐式学习深度，而 VGD 改为从预训练视觉几何基础模型 VGGT（Visual Geometry Grounded Transformer）中蒸馏几何先验，建立显式几何监督。

\textbf{动机。} 稀疏环绕视角下自监督深度学习的几何一致性难以保证，新视角常出现重叠伪影和几何失真；预训练基础模型的几何先验比光度约束更可靠，因此需要将几何信息显式学习而非隐式自监督。

\textbf{核心改动——VGGT 蒸馏与双分支架构。} DrivingForward 的深度网络 $D$ 与高斯网络 $G$ 串联预测；VGD 将其重构为 DPT-Depth $G(\cdot)$ 与 DPT-GS $\hat{G}(\cdot)$ 双分支并行架构。具体地，（1）\textit{几何蒸馏}：设计轻量 VGGT 变体，参数量仅为原始 VGGT 的 5\%，通过软蒸馏损失 $L_{distill}$（Smooth L1 + 仿射不变损失）从预训练 VGGT 教师网络迁移深度先验 $O_{geo}$，替代 DrivingForward 的三类光度损失自监督范式。（2）\textit{几何引导的高斯预测}：DPT-GS 分支不再独立预测，而是显式融合 DPT-Depth 分支提取的多尺度特征 $f_d = \{f_d^{(k)}\}_{k=1}^4$，使高斯参数（旋转 $R$、不透明度 $O$、缩放 $S$、球谐 $SH$）的预测受几何约束。（3）\textit{语义精化模型（SRM）}：DrivingForward 直接输出渲染结果；VGD 新增轻量 U-Net 精化模块（参数量不足 VGGT 的 0.2\%），融合双分支的多尺度特征 $f_d$ 与 $f_{gs}$，通过残差学习 $I_{novel} = I_{render} + E(I_{render}, f_d \| f_{gs})$ 精化渲染质量。整个框架以联合损失 $L = \lambda_r L_{render} + \lambda_d L_{distill} + \lambda_l L_{loc}$（$\lambda_r{=}0.01$, $\lambda_d{=}0.0005$, $\lambda_l{=}1$）端到端训练。

\textbf{设计动机。} DrivingForward 的自监督深度在相机间重叠极小（$<10\%$）的环绕视角配置下，几何一致性难以保证，导致新视角出现重叠伪影和几何失真。VGD 认为：在稀疏环绕视角场景中几何信息必须被显式学习而非隐式自监督，预训练基础模型的几何先验比光度约束更可靠。

\textbf{性能。} 在 nuScenes 上，VGD 单帧模式分别以高分辨率（$352{\times}640$）和低分辨率（$114{\times}228$）相比 DrivingForward 取得 $+1.69$ dB 和 $+2.11$ dB 的 PSNR 提升，多帧模式下 PSNR 提升 $+1.01$ dB、推理速度约 $1.6\times$ 快于 DrivingForward，在极端夜间场景中几何对齐和光照恢复也显著优于后者。

\paragraph{3D 查询点时空融合的前馈重建}

PointForward（2026.05）~\cite{2605_11594} 摒弃 DrivingForward 的逐像素高斯预测范式，转而在世界空间中初始化稀疏 3D 查询点（每个查询是 12 维向量：3D 坐标 + RGB 颜色 + 6 维 Pl\"ucker 射线坐标），通过时空融合聚合多视角信息，在单次前馈中强制跨视图一致性。

\textbf{与 DrivingForward 的差异。} DrivingForward 从单帧图像预测逐像素深度，本质是 ``2.5D'' 表示；PointForward 利用预训练 EfficientNetV2 提取图像特征和冻结的 InfiniDepth 预测密集深度，将像素直接提升到 3D 空间形成稀疏查询集 $Q = Q_s \cup Q_d$（$V \times N_s$ 个静态查询 + $O \times N_d$ 个动态查询），避免了像素级预测中的视角冲突和冗余计算。

\textbf{场景图驱动的动态建模。} 每个 3D 查询点通过可变形注意力机制从多视角图像特征中采样信息，预测对应的高斯属性（位置 $\mu$、旋转 $r$、缩放 $s$、不透明度 $o$、高维颜色特征 $c$）。利用场景图显式组织运动实例——将落入 3D 边界框的查询标记为动态并变换到实例规范空间（以首个边界框为参考），其余标记为静态，实现实例级运动传播和时间一致动态表示。渲染时动态查询从规范空间变换回目标时间戳的世界坐标，高维颜色特征经轻量 UNet 解码为 RGB 图像。不透明度通过时间衰减调整：$o_{t'} = o \cdot e^{-12(t'-t+1)^2}$，确保高斯仅在对应时间窗口内有效。

\textbf{时空融合的跨视图一致性。} 对每个查询，将其 3D 坐标投影到所有输入图像采样特征 $f_i$，同时计算深度差异 $d_i$ 反映几何一致性。融合权重综合考虑深度可靠性和时间邻近性：$w_i = \sigma(W_d \cdot \text{MLP}([f_{s,i}; d_i]))$，其中动态查询在静态区域采样的特征被置零以避免不一致的运动特征。最终融合特征 $f = \sum_{i=1}^{V} w_i f_{t,i}$ 聚合所有视角信息，实现全局一致的 3D 表示。

\paragraph{VGGT 驱动的混合预测头 4D 重建}

ReconDrive（2026.03）~\cite{2603_07552} 在 DrivingForward 的前馈框架基础上引入 3D 基础模型 VGGT 作为几何主干，实现快速高保真 4DGS 生成。

\textbf{与 DrivingForward 的差异。} DrivingForward 使用自监督深度网络预测高斯位置；ReconDrive 直接利用预训练 VGGT（冻结权重 + LoRA 适配，rank $r{=}8$、scaling $\alpha{=}32$）的交替注意力机制提取 24 层 Transformer 特征，通过 DPT 上采样恢复到原始图像分辨率后预测像素级深度图，再利用相机内外参将深度反投影为高斯 3D 中心——这一流程将高斯中心精确定位在自车标定坐标系中，避免了 VGGT 直接用于驾驶场景时的空间对齐偏差。

\textbf{混合高斯预测头。} 引入混合预测头解耦空间坐标和外观属性回归：高斯中心预测头（GCPH）从 Transformer 特征预测深度图并反投影为 3D 中心；高斯参数预测头（GPPH）在上采样特征与原始图像的拼接 $FP_t^o = [F_t^o; I_t^o]$ 上通过卷积层预测旋转、缩放、不透明度和球谐系数。关键设计是快捷连接将原始图像直接拼接到上采样特征中，捕获 Transformer 下采样过程中丢失的高频纹理和颜色细节，避免外观属性的模糊退化。这种解耦使得空间头专注于几何一致性、外观头专注于光度保真度，避免多任务梯度干扰。

\textbf{静动态 4D 组合。} 利用 SAM2（Segment Anything Model 2，通用视频分割模型）提取实例级动态物体掩码，结合 3D 边界框标注估计每个动态物体的速度向量 $\mathbf{v} \in \mathbb{R}^3$（假设短时间内匀速刚体运动）。动态高斯的中心随时间线性运动 $\mu_i(t) = \mu_i(T_s) + \mathbf{v}(i)(t - T_s)$，静态高斯的中心保持不变。两帧上下文的高斯通过时空对齐融合（将 $T_{s+1}$ 帧的高斯变换到 $T_s$ 的自车坐标系并按速度场补偿动态物体位置），拼接后形成最终的 4DGS 表示。

\textbf{性能。} 在 nuScenes 上性能媲美逐场景优化方法（如 Street Gaussians、OmniRe），但速度快数个数量级——前馈推理仅需数秒而非数十分钟训练。

\paragraph{因果动态速度场的无标注前馈重建}

StreetForward（2026.03）~\cite{2603_19552} 与 ReconDrive 同样基于 VGGT 主干，但核心区别在于动态建模方式：ReconDrive 依赖外部检测器（SAM2 + 3D 边界框）估计速度，StreetForward 则通过因果动态模块直接从图像序列学习逐像素速度场，无需任何物体标注或追踪器。

\textbf{因果动态模块。} 在 VGGT 交替注意力层输出的 token 上附加正弦时间嵌入 $\tau_f \in \mathbb{R}^{d_t}$，然后通过帧结构化因果掩码的多头注意力层传递时间信息——掩码 $M[b,h,i,j]$ 仅允许 $frame(i)$ 对 $frame(j)$ 进行前向或后向相邻帧的注意力（$f_t = f_s \pm 1$），严格禁止跨帧信息泄漏。这一设计使网络隐式提取运动场：通过关联当前帧与相邻帧的特征，DPT 风格的运动解码器预测逐像素的前向和后向 3D 速度 $\mathbf{v}_f^+,\mathbf{v}_f^- \in \mathbb{R}^3$ 以及动态概率 $\sigma_{f,u}$。关键创新是丢弃了 PVG 等方法中常用的逐高斯寿命属性——静态 GS 基元在整段序列中持续存在，通过反不透明度塌缩正则化防止优化器在新视角合成中抑制这些基元。

\textbf{跨帧时序融合渲染。} 渲染帧 $f$ 时，故意排除当前帧自身的动态高斯，强制动态物体必须由相邻帧的动态高斯通过速度传播来解释：$G^f = G^{static} \cup_{t \neq f} G^{dynamic}_f(t)$。前向-后向对称损失 $L_{fb} = \sum_u \|\mathbf{v}_u^{f \to f+1} + \mathbf{v}_u^{f \to f-1}\|_1$ 惩罚速度场的不一致性，局部刚体正则化在 2D 邻域窗口和 3D KNN 邻域内鼓励邻近高斯共享相似速度。在 Waymo 上训练后零样本迁移到 CARLA 等数据集。

\paragraph{遮挡演化的多车协作前馈重建}

FRUC（2026.05）~\cite{2605_29997} 面向多车协作场景，与 StreetForward 的单车前馈重建不同，FRUC 解决的是分布式多车网络中视角极端不对称且无需精确标定的挑战——基础设施相机固定在高处俯瞰道路，车辆相机处于低位前向视角，两者之间几乎不存在视角重叠，且无需硬件级时间同步。

\textbf{自车中心因果遮挡场。} 核心洞察是：自车的盲区主要由前方动态物体造成，且这些盲区随物体运动持续演化——遮挡新区域的同时暴露先前隐藏区域。FRUC 首先通过 DPT 风格动态头估计基础动态概率图 $S_{dyn} \in [0,1]^{H \times W}$ 识别当前移动物体的空间足迹。然后，在聚合的多智能体 token 上施加智能体级因果掩码注意力——严格禁止不同智能体或未来帧之间的信息泄漏——提取运动演化 token $Z_{kin}$ 作为隐式运动追踪器。轻量遮挡解码器从基础足迹出发预测时空扩展残差 $\Delta S_{kin}$，通过残差扩张学习遮挡的时间演化，最终遮挡不确定性先验 $M_{occ} = \text{clip}(S_{dyn} + \Delta S_{kin}, 0, 1)$ 为后续跨智能体融合提供结构化空间先验。

\textbf{零初始化残差去噪。} 跨智能体特征对齐被建模为约束残差去噪过程——以自车独占特征 $F_e$ 为参考，双分支卷积架构的条件分支编码遮挡先验和参考特征，基础分支处理混合 token。关键设计是所有注入层使用零初始化 $1 \times 1$ 卷积（$Z_{pri}$ 和 $Z_{out}$），使得训练初期跨智能体更新被自动抑制，网络从纯自车重建起步逐步学习有针对性的残差校正。在遮挡区域外（$M_{occ} \approx 0$），零卷积自然压制协作更新保护自车可靠观测；在动态盲区内（$M_{occ} > 0$），残差更新动态聚合结构上下文修复缺失背景。两阶段渐进训练课程（单智能体预训练 + 跨智能体适配）确保高度非凸的多智能体优化转化为有界稳定的残差学习。在 V2XReal 和 UrbanIng-V2X 数据集上达到新的最优性能。

\subsubsection{外观解耦与多轨迹重建}

本路线介绍将场景的内在几何/材质属性与外在光照/传感器效果解耦的方法，解决多轨迹数据中的跨视角光度不一致问题。

自动驾驶数据通常在不同时间、不同光照条件下采集，同一道路在同一天的早晨和傍晚呈现截然不同的外观。这种跨视角的光度不一致性严重影响了重建质量：优化器在尝试同时拟合不同光照条件下的颜色时，会引入几何伪影（如路面上的"浮游体"）或外观模糊。将场景的内在几何/材质属性与外在光照/传感器效果解耦，是解决这一问题的关键思路。

\paragraph{多尺度双边网格}

Wang et al.（2025.06）~\cite{2506_05280} 提出统一外观码（appearance code）和双边网格的多尺度双边网格方法，系统解决了自动驾驶数据中跨视角光度不一致的问题。

\textbf{问题分析。} 自动驾驶重建中的光度不一致分为两类：一致的内在场景颜色 $C_{c,t}$（物体真实颜色、恒定的传感器调整如归一化曝光）和不一致的非本质变化 $N_{c,t}$（光照变化、不同相机设置、多样化 ISP 效果）。传统方法使用单一外观码（对整张图像建模能力有限）或标准双边网格（空间分辨率不足），无法有效处理从全局光照变化到局部阴影的多尺度光度变化。

\textbf{核心改动——三级双边网格金字塔。} 提出层次化三级双边网格架构：（1）粗粒度层（$2 \times 2 \times 112$ 网格）捕获全局场景结构和近似全局外观码；（2）中间层（$4 \times 4 \times 212$ 网格）表示区域特征和中等范围光度变化；（3）细粒度层（$8 \times 8 \times 412$ 网格）编码局部细节和像素级双边网格。每级网格存储展平的 $3 \times 4$ 仿射颜色变换矩阵（12 通道），通过亮度引导图进行三线性插值获取逐像素变换系数。各级变换从粗到细顺序组合：$I^e = T^{(L-1)} \circ T^{(L-2)} \circ \cdots \circ T^{(0)}(I^r)$，实现尺度自适应的光度校正。

\textbf{设计动机。} 将光度校正与几何优化联合建模——$C_{c,t}^r$ 的调整同时受到几何约束 $D_{c,t}^r$ 的引导，避免将路面阴影误识别为几何变化。降采样引导图在粗粒度层的应用大幅提升计算效率，同时通过空间聚合增强光度变换的鲁棒性。

\textbf{性能。} 在 Waymo、nuScenes、Argoverse、PandaSet 四个数据集上显著提升几何精度，有效抑制由光度不一致导致的浮游体。

\paragraph{材质-光照双层分解的多轨迹重建}

多轨迹数据指同一地理区域在不同时间段被多次采集的数据——例如同一段道路在晴天和阴天各驶过一次。核心观察是：静态几何跨轨迹高度一致（路面、建筑不变），但外观因光照变化而剧烈不同。如果直接将所有轨迹的数据丢给标准 3DGS 一起优化，优化器会试图让同一高斯基元同时拟合不同光照下的颜色，结果要么产生几何伪影（路面"浮游体"），要么外观模糊——这正是 Street Gaussians 和 DrivingForward 等方法在多轨迹场景下遇到的关键瓶颈。

ADM-GS（2026.04）~\cite{2604_05908} 修改了 3DGS 的外观表示环节——用``材质--光照显式分解''替代单一球谐颜色建模——几何骨架和高斯泼溅渲染管线保持不变。

\textbf{Baseline 做法。} 标准 3DGS 为每个高斯基元学习一组球谐系数直接预测视角相关颜色。这种方法不区分材质和光照：同一基元在不同轨迹中因光照不同呈现不同颜色，但球谐系数只有一个，优化只能取折中，导致重建质量下降。

\textbf{核心改动——材质/光照双层分解。} ADM-GS 将渲染颜色拆分为两个因子的 Hadamard 积（逐元素相乘）：

\begin{equation}
\hat{C}(x, v, m) = M(x) \odot L(x, n, v, r, e_m)
\end{equation}

直觉上，可以将其理解为``底色''$M$（材质，跨轨迹不变）乘以``光照滤镜''$L$（随轨迹和视角变化）。类比摄影：同一面墙在不同天气下拍摄，墙的漆色 $M$ 不变，但光照 $L$ 不同，最终照片颜色是两者的综合效果。等式左侧的 $\hat{C}$ 是最终渲染像素颜色，$x$ 是空间位置，$v$ 是视角方向，$m$ 是轨迹索引。

\textbf{材质场——轨迹不变的底色。} 材质映射只依赖空间几何特征，刻意不接触任何视角或轨迹信息，确保分解的一致性：

\begin{equation}
M(x) = \phi_{\text{material}}(f_{\text{geo}}(x)), \quad M(x) \in [0,1]^3
\end{equation}

其中 $f_{\text{geo}}(x)$ 是从高斯场提取的空间几何特征。输出限制在 $[0,1]^3$ 以保证材质值为合法颜色。为处理轨迹间的瞬态几何不一致（如某轨迹中有临时广告牌），引入几何门控机制：独立 MLP 预测不透明度衰减项，抑制瞬态物体区域，使材质优化仅作用于几何一致的区域。

\textbf{神经光照场——轨迹相关的照明。} 光照场同时依赖表面法向 $n$、视角 $v$、反射方向 $r$ 和轨迹嵌入 $e_m$：

\begin{equation}
L(x, n, v, r, e_m) = \phi_{\text{light}}(n, v, r, f_{\text{geo}}(x), e_m), \quad L \in \mathbb{R}_+^3
\end{equation}

输出约束为正数，以覆盖户外照明的大动态范围。轨迹嵌入 $e_m$ 是每条轨迹的可学习向量，使同一位置的照明效果随轨迹不同而变化。

\textbf{设计动机——频率分离混合编码。} 漫反射和镜面反射在频域中呈现截然不同的特征：漫反射主要由低频表面朝向决定（粗略地说，``面朝光源多亮，背朝光源多暗''），而镜面反射包含极高频的视角相关细节（金属和玻璃表面的锐利高光）。若将两者混在一起用同一编码频率建模，要么漫反射被高频噪声干扰，要么镜面高光因编码分辨率不足而模糊。

为此，ADM-GS 设计了频率分离策略：对表面法向 $n$ 仅使用 1 阶球谐编码（低频），捕获漫反射成分；对反射向量 $r$ 使用 4 阶球谐编码（高频），捕获镜面反射的锐利高光；对视角方向 $v$ 使用 3 阶球谐编码作为补充。融合后的特征向量输入光照 MLP：

\begin{equation}
f_{\text{in}}(x, v, m) = [\Phi^{\text{deg}=1}(n),\; \Phi^{\text{deg}=4}(r),\; \Phi^{\text{deg}=3}(v),\; f_{\text{geo}}(x),\; e_m]
\end{equation}

即低频编码看法向（漫反射），高频编码看反射方向（镜面反射），两者拼接后由 MLP 统一融合。反射向量 $r$ 通过 $r = 2(n \cdot v)n - v$ 计算。这一重参数化将输入域对齐到镜面反射瓣的峰值方向，简化了网络的学习任务（从学习空间变化的旋转函数退化为拟合更平滑的环境光照映射），这是 Ref-NeRF 验证过的有效策略。

\textbf{几何正则化——法向质量保障。} 原始 3DGS 缺乏显式几何约束，高斯基元往往是蓬松的体积块而非贴合表面的扁平椭球，导致法向噪声大、光照估计不准。ADM-GS 引入各向异性平坦度约束：

\begin{equation}
\mathcal{L}_{\text{scale}} = \frac{1}{|G|}\sum_k \max(0, -((s_k^{\min}) - (s_k^{\max}))), \quad \delta = 1.0
\end{equation}

其中 $s_k^{\min}$ 和 $s_k^{\max}$ 是高斯基元 $k$ 缩放向量的最小值和最大值，$\delta=1.0$ 对应惩罚 $s_k^{\min}/s_k^{\max} > e^{-1} \approx 0.37$ 的不够平坦的高斯。最短轴方向 $u_k^{\min}$ 直接作为几何法向 $n_k$，无需额外的法向预测网络。

\textbf{传感器对齐。} 不同轨迹的相机自动曝光和白平衡设置不同，为补偿这种光度差异，对光栅化后的图像施加可学习的逐轨迹仿射变换：$\hat{I}_{m,c,t} = s_m \cdot I^{\text{render}}_{m,c,t} + b_m$，其中 $s_m \in \mathbb{R}^3$ 和 $b_m \in \mathbb{R}^3$ 是轨迹特定的缩放和偏移参数，将渲染输出映射到轨迹 $T_m$ 的图像空间。

\textbf{性能。} 在 Argoverse 2 和 Waymo 上比潜在编码基线提升 +0.98 dB PSNR。

\paragraph{外观嵌入引导的前馈重光照}

SpectralSplat（2026.04）~\cite{2604_03462} 在 ADM-GS 将外观解耦引入逐场景优化之后，首次在前馈框架中实现外观与几何的解耦。

\textbf{与 ADM-GS 的差异。} ADM-GS 为每个高斯学习独立的材质 MLP 和光照 MLP，需要逐场景优化；SpectralSplat 基于 UniSplat 前馈骨干，将颜色预测重构为两次共享颜色 MLP 的评估——外观无关基础流（以零嵌入为输入预测固有颜色 $c^{base} = \sigma(f_{rgb}([z^{vox}; z^{img}; \mathbf{0}]))$）和外观条件适应流（以全局外观嵌入为输入预测适应颜色 $c^{ada} = \sigma(f_{rgb}([z^{vox}; z^{img}; a]))$），其中 $\sigma$ 为 sigmoid，$a \in \mathbb{R}^{64}$ 是从 DINOv2 全局平均池化 patch token 经轻量编码器生成的全局外观嵌入。

\textbf{设计动机。} 驾驶场景的外观变化（色调、色偏、光照氛围）主要是全局性的——同一场景在不同天气下整体色调不同但局部纹理不变。单一共享外观嵌入直接施加这一先验，防止逐点外观漂移。外观嵌入仅从当前时间戳的相机计算，未来帧渲染复用当前嵌入——迫使模型生成跨视角泛化的几何，而外观码仅捕获当前时间的全局光照状态。

\textbf{混合重光照训练。} 核心挑战是缺乏配对的不同外观观测。SpectralSplat 设计了物理-生成混合重光照管线：先用 MVInverse 估计逐像素反照率和表面法向，计算 Lambertian 重渲染 $\hat{I}^{mv}$ 保证多视角 3D 一致性但缺乏高频细节；再用 IC-Light 扩散重光照模型精化，通过频谱解耦策略在去噪过程中注入低频物理先验——在每个 DDIM 步替换估计清洁潜变量的低频分量为物理参考的低频分量，权重 $w(t)$ 线性衰减（$\lambda \in [0,1]$），早期步锁定光照方向和色调、后期步放松约束保留扩散模型的高频细节生成能力。

\textbf{外观可适应时间历史。} 修改 UniSplat 的时间累积机制——缓存的静态点高斯保存其 FPN 图像特征 $z^{img}_{cached}$，在后续时间戳回忆时用当前外观嵌入重新计算颜色 $c^{hist} = \sigma(f_{rgb}([z^{vox}_{fused}; z^{img}_{cached}; a^{now}]))$，使历史高斯和当前高斯在相同目标外观下渲染，消除跨时间戳的颜色不一致。

\paragraph{场景图驱动的多轨迹高斯泼溅}

MTGS（2025.03）~\cite{2503_12552} 与 ADM-GS 同样处理多轨迹外观变化，但采用场景图而非物理分解的策略。

\textbf{动机。} 多轨迹数据中静态几何跨轨迹一致但外观因光照剧变，直接联合优化导致几何伪影或外观模糊；需要一个轻量级的数据驱动策略隐式学习外观校正，同时支持不同轨迹动态节点的灵活组合。

\textbf{与 ADM-GS 的差异。} ADM-GS 通过材质/光照双层分解建模外观，需要为每个高斯维护两个 MLP；MTGS 采用更轻量的多轨迹动态场景图策略——共享静态节点存储跨轨迹复用的背景高斯，轨迹专属动态节点为每条轨迹独立建模动态物体高斯，颜色校正节点通过可学习球谐系数残差补偿轨迹间的光照和色调差异（$\Delta c = c^{traj} - c^{shared}$）。这种方法无需显式的物理光照模型，通过数据驱动的方式隐式学习外观校正。

\textbf{跨轨迹场景组合。} 场景图在推理时支持选择性地组合不同轨迹的动态节点——例如将轨迹 A 的车辆插入轨迹 B 的背景中——生成原数据集中不存在的交通场景组合。共享静态几何保证背景一致性，颜色校正节点确保插入物体的外观与目标轨迹匹配。

\textbf{性能。} 在 nuPlan 多轨迹数据上，相比单轨迹基线 LPIPS 提升 23.5\%、几何精度提升 46.3\%，验证了多轨迹数据对重建质量的显著增益。

\paragraph{HDR 辐射不变性的物理光度校正}

P2GS（2026.05）~\cite{2605_16925} 与 ADM-GS 和多尺度双边网格同样处理光度不一致，但采用完全物理驱动的 HDR 域分解——从仅 LDR 图像联合分解三个物理分量：视角不变的线性 HDR 辐射场（场景的真实辐射亮度 $\hat{I}_{linear}$）、逐视图的曝光尺度 $e_i \in \mathbb{R}^+$ 和色调映射函数 $T_i(x) = \text{clamp}(x^{1/\gamma_i}, 0, 1)$。

\textbf{与多尺度双边网格的差异。} 多尺度双边网格学习仿射颜色变换矩阵，是纯数据驱动的光度校正；P2GS 的核心创新是辐射不变性原理（Principle of Invariant Radiance, PIR）——在线性 HDR 域中，不同视角观测到的辐射比值仅取决于曝光：$\hat{I}^{(i)}_{linear} / \hat{I}^{(j)}_{linear} = e_i / e_j$。这一不变性直接约束 HDR 辐射和曝光参数，通过相对曝光一致性损失 $L_{exp} = \frac{1}{M}\sum_{(i,j) \in P} \| (e_j/e_i) \hat{I}^i_{linear} - \hat{I}^j_{linear} \|_1$ 嵌入优化，在无需 HDR 监督的情况下实现物理正确的辐射-曝光解耦。

\textbf{HDR 域优化。} 球谐系数被优化为线性 HDR 辐射而非直接拟合 LDR 像素值，曝光参数初始化为 $e_i \sim \mathcal{N}(1.0, 0.05^2)$，色调参数 $\gamma_i$ 初始化接近 sRGB 先验（$\gamma \sim 2.2$）并通过正则化 $L_{reg}$ 约束——曝光尺度锚定 $E[(e_i-1)^2]$、曝光方差惩罚 $\text{Var}(e_i)$ 和色调先验 $E[(\gamma_i - \gamma_{prior})^2]$ 联合稳定优化，避免全局尺度歧义和退化的色调曲线。推理时取所有训练视角曝光和色调参数的均值，消除视频渲染中的时间闪烁。

\subsubsection{轨迹外推与车道变换}

本路线介绍突破训练视角限制、使重建场景在偏离训练轨迹的视角下仍保持高保真渲染的方法，包括协方差约束、扩散蒸馏和虚拟车道训练等技术。

标准 3DGS 重建在训练轨迹附近渲染效果良好，但偏离到新车道或方向时质量急剧下降。这是因为高斯基元的协方差和密度在仅有前向相机监督时会过拟合到训练视角分布，当从侧面观察时，原本覆盖正面视角的"扁平"高斯暴露了其几何缺陷。轨迹外推的目标正是突破这一限制，使重建场景在偏离训练轨迹的视角下仍保持高保真渲染。

\paragraph{协方差约束与扩散蒸馏的外推视角合成}

VEGS（2024.07）~\cite{2407_02945} 首次系统研究城市场景中的外推视角合成（EVS）问题，揭示了前向相机优化导致高斯协方差"懒惰优化"的关键瓶颈。

\textbf{Baseline 与问题。} 标准 3DGS 在训练轨迹附近渲染质量良好，但偏离到新车道或方向时质量急剧下降：高斯基元的协方差和密度在仅有前向相机监督时会过拟合到训练视角分布，从侧面观察时暴露几何缺陷。此前缺乏系统研究这一外推视角合成（EVS）问题的工作。

\textbf{核心发现——懒惰协方差优化。} VEGS 发现标准 3DGS 在仅有前向相机时，高斯协方差的形状和方向会过拟合到特定视角——协方差被训练为以最小优化代价覆盖训练像素的视锥。当从非训练方向观察时，这些过拟合的协方差在表面下方产生空洞（cavity），导致外推视角渲染质量急剧下降。

\textbf{三项技术贡献。} （1）密集 LiDAR 地图初始化——将多帧 LiDAR 扫描聚合为密集静态地图和逐实例动态点云，替代稀疏 SfM 初始化，为高斯位置提供更可靠的几何先验；（2）表面法向引导协方差优化——设计协方差轴损失 $L_{axis}$（通过球面线性插值渲染协方差方向图，与法向估计器预测的表面法向对齐）和协方差缩放损失 $L_{scale}$（最小化与法向对齐的轴的缩放值，迫使协方差变平贴合表面）；（3）大规模扩散模型知识蒸馏——通过 LoRA 微调 Stable Diffusion 获得场景特定的视觉先验，利用去噪得分匹配对外推视角渲染提供直接监督，平衡场景特定知识与泛化能力。

\textbf{关键公式。} 协方差轴损失和缩放损失联合约束高斯形状贴合表面：

\begin{equation*}
\mathcal{L}_{axis} = \|Slerp(R_k \mathbf{s}_{min}^k, \hat{\mathbf{n}}_k)\|, \quad \mathcal{L}_{scale} = \sum_k s_{min}^k
\end{equation*}

其中 $\mathbf{s}_{min}^k$ 是高斯基元 $k$ 缩放向量的最小分量方向（经旋转 $R_k$ 变换到世界坐标后与估计法向 $\hat{\mathbf{n}}_k$ 对齐），$s_{min}^k$ 为最小缩放值，最小化其迫使高斯在法向方向变薄。直觉上，$\mathcal{L}_{axis}$ 确保高斯最薄方向对齐表面法向，$\mathcal{L}_{scale}$ 迫使该方向厚度进一步压缩——两者协同将过拟合的胖椭球压扁为紧贴表面的薄片。扩散监督通过去噪得分匹配对外推视角渲染施加先验：$\mathcal{L}_{SDS} = \mathbb{E}_{t,\epsilon}[\omega(t) \| \epsilon_\theta(z_t, t) - \epsilon \|_2^2]$，其中 $z_t$ 为外推视角渲染图像在时间步 $t$ 的噪声潜变量。

\textbf{性能。} 在外推城市场景视角合成基准上显著超越标准 3DGS 和其他基线方法。

\paragraph{逆视角变形与深度引导的离轨渲染}

FlexDrive（2025.02）~\cite{2502_21093} 提出逆视角变形（Inverse View Warping, IVW）技术为离轨视角创建高质量监督图像，配合深度引导策略（Depth Bootstrapping）在优化过程中获取密集深度图。

\textbf{逆视角变形（IVW）。} 核心思路是利用训练轨迹上的真实图像为离轨虚拟视角提供监督。具体流程分三步：（1）变形射线图生成——将训练视角 $V_{in}$ 的像素通过渲染深度反投影到 3D 空间，再投影到随机采样的离轨视角 $V_{out}$，生成变形射线图；（2）遮挡感知光栅化——仅对深度大于反投影点的 3D 高斯进行 $\alpha$-混合，避免遮挡区域引入错误颜色监督；（3）像素重排——将 $V_{out}$ 中的渲染像素恢复为 $V_{in}$ 中的相对空间顺序，使其可与 $V_{in}$ 的真值图像直接比较。

\textbf{深度引导（DB）。} IVW 依赖密集深度图，但 LiDAR 深度稀疏且有噪声。FlexDrive 提出两步策略：稀疏深度初始化（多帧 LiDAR 聚合，去除遮挡穿透点和重复点）和密集深度校正（利用渲染深度与稀疏深度之间的强线性关系，通过最小二乘法求解线性变换参数 $D_r^* = a D_r + b$ 校正渲染深度）。随着训练进行，密集深度越来越准确，形成正反馈循环。

\textbf{约束动态物体建模。} 将每个动态物体约束在边界框内，用 sigmoid 函数将逻辑坐标映射到 Euclidean 坐标，有效消除动态物体周围的拖尾浮游体。在 Waymo 上显著提升离轨渲染性能。

\paragraph{虚拟车道闭环训练的跨车道泛化}

GGS（2024.09）~\cite{2409_02382} 修改了可泛化 3DGS（MVSplat）的新视角合成环节——用``虚拟车道生成 + 扩散自监督''替代仅依赖输入视角邻域的渲染策略——3D 高斯基元参数化保持不变。其核心问题是：自动驾驶数据几乎总是从单车道采集的，当渲染视角偏离到相邻车道时，模型从未见过该视角下的场景内容，导致严重的伪影和信息缺失。与 VEGS 通过扩散先验正则化协方差、FlexDrive 通过逆视角变形创建离轨监督不同，GGS 的策略是主动构造虚拟车道数据，让模型在训练阶段就``见过''相邻车道。

\textbf{Baseline 做法。} MVSplat 是一种可泛化 3DGS 方法：给定 $N$ 张稀疏输入图像，通过跨视角注意力构建代价体（cost volume），经 U-Net 预测逐像素深度和高斯参数（位置 $x$、缩放 $s$、旋转 $q$、不透明度 $\alpha$、颜色 $c$），无需逐场景优化即可渲染新视角。然而，MVSplat 只能渲染与输入图像视角接近的新视角——当视角偏移量大到需要``换一条车道''时，模型缺乏该方向的任何信息，渲染质量急剧下降。

\textbf{改动一——多视角深度精化。} MVSplat 的深度估计在遮挡区域不准，而准确深度是后续虚拟车道生成的基础。GGS 引入两个改进。

第一，邻域特征融合：将第 $i$ 帧前后各 $k$ 帧的颜色特征拼接到深度特征中，通过 U-Net 融合：

\begin{equation}
F_{\text{neighbor}}^{(i)} = \bigcup_{m \in [i-k,\, i+k]} F_m, \quad \text{dep}_{\text{ref}} = \mathcal{U}(F_{\text{neighbor}}^{(i)},\; \text{dep}_i)
\end{equation}

直觉上，当某帧的某个像素被遮挡时，前后帧可能从略微不同的角度``看到''该位置，融合邻域信息相当于为遮挡区域补充了间接观测。

第二，置信度引导深度校正：利用重建点云（通过 Agisoft Metashape 生成）作为补充深度先验。当预测高斯的不透明度低于阈值 $\tau$（即模型对该像素的深度估计不自信）时，用重建深度 $D_i$ 替换预测深度；否则保留原始预测。这样可以有选择地校正低置信区域，而不破坏高置信区域的准确预测。

\textbf{改动二——虚拟车道生成模块。} 这是 GGS 的核心创新。关键洞察是：即使没有真实的多车道数据，也可以通过``换出去再换回来''的闭环训练让模型学会车道变换。

训练分两个阶段。第一阶段，输入 $N$ 张单车道图像 $ISet_1 = \{I_1, I_2, \ldots, I_N\}$，渲染当前车道视角：

\begin{equation}
\hat{I}^1 = \mathcal{G}(ISet_1)
\end{equation}

然后通过虚拟车道转换器将渲染图像水平平移 $\delta$ 并施加透视变换，生成虚拟相邻车道的图像集合 $ISet_2 = \mathcal{V}(\{I_k^1, \delta\})$，其中 $\delta$ 是平移系数，每帧的切换角度按周期变化。第二阶段，将虚拟车道图像作为新输入，让模型``换回''真实车道：

\begin{equation}
\hat{I}^2 = \mathcal{G}(ISet_2)
\end{equation}

这样形成闭环：$\text{真实车道} \to \text{虚拟车道} \to \text{真实车道}$。由于第二阶段的渲染结果可以与第一阶段的真值直接比较，因此无需任何多车道标注就能构建车道变换的训练信号。设计动机在于，自动驾驶场景中车道线平行且路面平坦，水平平移在几何上是近似合理的车道变换模拟——模型通过反复练习``换出再换回''，逐步学会处理大视角偏移。

\textbf{改动三——多车道扩散损失。} 虚拟车道图像存在一个天然问题：当车道变换幅度大时，前方车辆等障碍物会遮挡新车道视野，导致虚拟车道图像中出现大片空白区域。GGS 利用扩散模型（Stable Diffusion）的先验知识来``想象''这些缺失内容。

具体流程：将多车道图像（左车道、中间车道、右车道）通过变分自编码器编码为潜表示，执行若干步去噪（以``autonomous driving''文本提示为条件），再解码为图像作为虚拟车道渲染的监督信号。扩散损失不是直接用扩散模型生成完整图像（这会导致帧间不一致），而是计算去噪前后图像之间的感知距离：

\begin{equation}
\mathcal{L}_{\text{diffusion}} = \mathbb{E}_{\pi, t}\left[\sigma(t) \cdot \mathcal{L}_{\text{LPIPS}}(y,\, \hat{y}_\pi)\right]
\end{equation}

其中 $\pi$ 是选取视角的相机位姿，$y$ 是多车道图像，$\hat{y}_\pi$ 是去噪模型输出，$\sigma(t)$ 是与噪声水平相关的权重函数，$\mathcal{L}_{\text{LPIPS}}$ 是感知损失（从人类视觉角度度量图像相似度）。直觉上，扩散损失的作用类似于``图像增强顾问''——它指出虚拟车道渲染中哪些区域看起来不自然，从而引导模型生成更符合真实场景的纹理细节。

\textbf{整体损失函数。} 总损失由四项组成：

\begin{equation}
\mathcal{L} = \mathcal{L}_{\text{recon}} + \mathcal{L}_{\text{depth}} + \mathcal{L}_{\text{switch}} + \mathcal{L}_{\text{diffusion}}
\end{equation}

其中 $\mathcal{L}_{\text{recon}}$ 是渲染图像与真值之间的均方误差，$\mathcal{L}_{\text{depth}}$ 是深度平滑约束（相邻像素的深度梯度应连续，包含一阶和二阶导数惩罚），$\mathcal{L}_{\text{switch}}$ 是虚拟车道闭环训练的重建损失（衡量``换出再换回''后的渲染质量），$\mathcal{L}_{\text{diffusion}}$ 是上述扩散感知损失。

\textbf{设计动机总结。} GGS 的三项改动各有明确的针对性：深度精化模块解决遮挡区域的深度不准问题（为虚拟车道提供可靠的 3D 基础）；虚拟车道生成模块通过闭环``换出--换回''构造自监督信号（解决多车道训练数据缺失的根本困难）；扩散损失弥补虚拟车道图像中大视角偏移导致的信息空白（利用视觉先验``填补''模型未见过的区域）。三者协同，使可泛化 3DGS 在仅使用单车道数据训练的情况下，仍能实现高质量的车道变换渲染。

\paragraph{奖励网络引导的两轮扩散增强重建}

RGE-GS（2025.06）~\cite{2506_22800} 与 VEGS 同样利用扩散先验增强外推视角渲染，但采用两轮重建策略——第一轮从原始视角真值图像重建初始 GS 场 $G_o$，第二轮利用扩散生成的外推视角伪图像扩展场景覆盖。

\textbf{与 VEGS 的差异。} VEGS 将扩散先验作为去噪得分匹配监督直接引导高斯协方差优化；RGE-GS 引入轻量奖励网络 $F_c$（3 层 U-Net）生成逐像素置信度图 $C_e \in [0,1]^{H \times W}$，以对抗验证方式动态调节扩散生成图像的监督贡献。奖励网络通过重投影误差损失 $L_{reproj} = \frac{1}{HW}\|C_e \odot (\text{Proj}(P) - I_e)\|_2$（将着色 LiDAR 点云投影到外推视角提供粗监督）和二值化正则化 $L_{reg} = \sum C_e(1-C_e)/(HW)$（鼓励置信度集中在 0 或 1 避免模糊预测）联合训练。

\textbf{差异化训练策略。} 将高斯基元分为三组按收敛状态施加不同策略：（1）缺失高斯——外推视角可见但原始视角被遮挡的区域，有 LiDAR 覆盖时用 LiDAR 点初始化，否则用 DepthAnything-V2 深度估计（以 LiDAR 测量为尺度校准参考）初始化；（2）未成熟高斯——视角空间位置梯度超过阈值 $\theta = 5 \times 10^{-4}$ 的基元，继续用外推视角图像监督优化直至收敛；（3）成熟高斯——固定属性不受生成图像影响，确保原始区域重建稳定性。奖励图加权重建损失 $L_{I_e} = \lambda_{RGB} L_{RGB}(I_d, I_e, C_e) + (1-\lambda_{RGB}) L_{SSIM}(I_d, I_e, C_e) + \lambda_{I_e} L_{LPIPS}(I_d, I_e)$，使扩散先验仅在可信区域提供监督。

\paragraph{3D 几何先验注入的扩散后处理增强}

SetDiff（2026.01）~\cite{2601_07540} 与 VEGS 同样利用扩散模型增强外推视角渲染，但 VEGS 将扩散先验嵌入 3DGS 优化过程，而 SetDiff 作为后处理增强——以 3DGS 渲染的降质图像为输入，通过扩散去噪恢复高频细节。

\textbf{动机。} 将扩散先验嵌入优化过程需要重新训练整个 3DGS 场景，代价高昂且不灵活；后处理增强方案可以利用任意预训练扩散模型，即插即用地提升任意已重建场景的外推渲染质量。

\textbf{显式 3D 几何先验注入。} 核心创新是将像素对齐坐标图（每个像素编码其对应的 3D 空间位置）和姿态感知 Pl\"ucker 射线嵌入（编码相机 6DoF 位姿的 6 维表示）作为显式 3D 先验融入集混合器（Set Mixer）。与标准扩散模型仅在潜空间操作不同，SetDiff 在全局标记级注意力中允许同一视角内的像素 token 和跨视角的像素 token 在同一注意力层中交互——参考视角和目标视角的所有像素 token 被展平为统一序列后执行全自注意力，使模型能从参考视角的可靠信息中推断目标视角缺失内容的合理外观。

\textbf{可变视角数量处理。} 集合框架天然支持可变数量的参考和目标视角——不同数量的输入图像无需修改网络结构即可处理，通过潜空间监督和选择性解码保持计算效率。在 EUVS、Para-Lane、nuScenes 和 DL3DV 上显著提升感知保真度（LPIPS）和结构相似度（SSIM），尤其在严重外推条件下鲁棒性明显优于无 3D 先验的扩散增强方法。

\paragraph{外推视角合成基准与多车道数据集}

前述 GGS 等方法虽然在单车道训练下实现了车道变换渲染，但缺乏标准化的外推评估基准——现有评测几乎都基于视角内插（训练与测试视角接近），无法量化模型在大视角偏移下的真实泛化能力。EUVS（2024.12）~\cite{2412_05256} 和 Para-Lane（2025.02）~\cite{2502_15635} 分别从``多源数据整合''和``专用多车道采集''两个互补方向填补了这一空白。

\textbf{EUVS：外推城市场景视角合成基准。} EUVS 整合了三个公开自动驾驶数据集——nuPlan（1,200 小时驾驶数据）、Argoverse 2（1,000 个标注场景）和 MARS（多智能体场景）——经 300+ 小时人工轨迹筛选和 800+ 小时 COLMAP 计算后，构建了包含 90,810 帧、345 个视频的评估数据集。其核心贡献是定义了三种外推评估设置：纯平移（模拟车道变换，训练与测试来自不同车道）、纯旋转（模拟摄像头切换，训练前方/后方相机、测试侧方相机）、平移+旋转（模拟路口场景，训练与测试从不同方向接近同一路口）。实验结果揭示了严重的过拟合问题：从内插到外推，PSNR 下降 23--29\%，SSIM 下降约 22\%，LPIPS 恶化约 50\%。此外，EUVS 还通过消融实验分析了平面基元（2DGS/GSPro）与椭球基元（3DGS/3DGM）在道路和纹理物体上的互补性、扩散先验（VEGS）对缺失区域的补全效果、以及多次遍历数据对背景重建的逐步增益。

\textbf{Para-Lane：多车道并行扫描数据集。} 与 EUVS 整合已有数据不同，Para-Lane 从零构建了首个真实世界跨车道 NVS 数据集。其采集方案是让自动驾驶车辆在同一路段的三条车道上分别行驶（多轮采集），每组场景包含三条共享起终点的关联序列，共 25 组场景、75 个序列、约 64,000 张环绕视角图像。硬件配置包括 1 个前视相机（$90^\circ$ FOV（Field of View，视场角））+ 4 个环视鱼眼相机（$190^\circ$ FOV）+ 3 个 32 通道 LiDAR，所有传感器在硬件层同步。数据集的关键技术难点是跨车道序列的精确位姿对齐——Para-Lane 提出两阶段优化：第一阶段通过 LiDAR 建图（LOAM 特征 + 稠密扫描配准 + 环路闭合）构建参考地图，第二阶段通过因子图优化将图像配准到 LiDAR 地图（融合 SfM 稀疏因子、LiDAR 点先验因子和跨模态光度因子）。评估设计了 5 种轨道（单车道回归、相邻车道预测、隔一车道预测、双车道训练预测、夹层预测），关键发现是：当训练视角均匀分布在测试视角两侧时（夹层轨道），渲染质量显著优于训练视角仅在一侧的情况（双车道轨道），且 Scaffold-GS 在跨车道任务上表现最佳——其锚点-MLP 机制建立了视角位姿与渲染结果之间的关联，天然适合外推场景。

\subsubsection{特殊场景与领域方法}

本路线介绍面向自动泊车、行车记录仪、夜间驾驶、无 LiDAR 重建等特定驾驶场景和传感器配置的街景重建方法。

除通用街景重建外，若干工作聚焦于特定驾驶场景和传感器配置的重建需求，涵盖自动泊车、行车记录仪、夜间驾驶、无 LiDAR 重建等 12 条技术路线。

\paragraph{无迹变换鱼眼投影的泊车场景重建}

ParkGaussian（2026.01）~\cite{2601_01386} 首次将 3DGS 引入自动泊车场景重建，创建 ParkRecon3D 基准数据集（约 $220\times110$ m$^2$ 地下车库、430+ 车位、四路环视鱼眼相机 $1280\times960$@10Hz、超过 40K 同步帧和 60K 手动标注车位角点）。技术核心包括两方面：其一，采用基于无迹变换（Unscented Transform）的投影替代传统 EWA 泼溅的雅可比线性化，直接在鱼眼图像上训练 3DGS 而无需逐相机推导雅可比，显著提升强畸变下的几何稳定性；其二，设计可微分逆透视映射（IPM）模块将渲染鱼眼图转换为鸟瞰图，并引入车位感知重建策略——利用预训练车位检测器 DMPR-PS 和 GCN-Parking 的冻结特征图构建车位角点与边缘的混合权重图（温度 $T=0.5$、混合系数 $\alpha=0.8$），反向传播监督信号至 3DGS 优化，使重建结果在车位区域的结构保真度显著提升。

\paragraph{透射-反射分解的挡风玻璃去遮挡}

DC-Gaussian（2024.05）~\cite{2405_17705} 专门处理行车记录仪视频中的挡风玻璃反射和遮挡问题。自适应图像分解模块将观测图像统一分解为场景透射分量（车外真实场景）和反射分量（挡风玻璃上的镜面反射），两者均由独立的 3D 高斯场建模。光照感知障碍物建模为挡风玻璃上的污渍、水滴等遮挡物附加光照条件属性，在不同光照条件下自动调整其不透明度和外观，处理反射随光照变化的问题。

\paragraph{基于物理渲染的夜间驾驶场景重建}

Kim et al.（2026.02）~\cite{2602_13549} 将基于物理的渲染（PBR）引入 3DGS 增强夜间驾驶场景重建。夜间场景的挑战在于光照极度不均匀（路灯、车灯产生局部高光区域）和低信噪比（暗区域细节丢失）。方法联合优化 BRDF 材质属性（反照率、粗糙度、金属度），通过全局光照模块显式建模漫反射分量（使用球谐函数近似环境光），各向异性球面高斯（anisotropic spherical Gaussians）建模镜面反射（车灯在湿润路面上的高光反射），在 nuScenes 和 Waymo 夜间场景上超越当时方法。

\paragraph{2D 边界框替代 3D 标注的精简管线}

S3R-GS（2025.03）~\cite{2503_08217} 与 Street Gaussians 的动静分离依赖 3D 边界框不同，S3R-GS 重新审视大规模街景重建管线，通过系统简化消除三个效率瓶颈。核心改动是用 2D 边界框替代 3D 边界框分离动静组件——在图像空间而非 3D 空间执行动静分割，避免了将 2D 检测结果提升到 3D 的额外变换和投影步骤；同时优化远距离内容渲染策略，减少不必要的投影计算。在 Argoverse2 上将重建时间降至竞争对手的 50\% 甚至 20\%。

\paragraph{鸟瞰图单帧实时道路表面重建}

BEV-GS（2025.04）~\cite{2504_13207} 提出基于鸟瞰图（Bird's Eye View, BEV）范式的单帧实时道路表面重建，分离几何和纹理网络直接从单帧估计参数。几何网络从 BEV 特征预测道路高程图（每个网格位置的高度值），纹理网络预测对应位置的外观参数；使用网格高斯表示道路表面——每个网格单元分配一个扁平高斯基元，高程值决定高斯的垂直位置。道路高程误差仅 1.73 cm，预测 FPS 达 26、渲染 FPS 达 2061。

\paragraph{平面高斯与 LOD 八叉树的可扩展街面重建}

StreetSurfGS（2024.10）~\cite{2410_04354} 首次将高斯泼溅专门应用于可扩展城市街面表面重建。其核心表示将 3D 高斯展平为 2D 平面高斯以更精确地捕捉城市场景的表面几何，并基于 SfM 稀疏点构建八叉树结构，将锚点高斯动态分配至不同细节层次（LOD）。针对街景中物体遮挡导致的深度不准问题，方法提出引导平滑正则化策略——通过离群点过滤（法向量夹角阈值）和边缘过滤（利用 SAM 分割掩码剔除边界点）构建二值掩码，确保法线一致性约束仅在平坦区域生效。为应对稀疏视角和多尺度挑战，双步匹配策略分别利用相邻相机和远距离长期信息评估光度与几何一致性，结合分段训练降低大规模街景的内存开销，最终通过 TSDF 融合提取高质量网格。

\paragraph{深度基础模型驱动的无 LiDAR 城市重建}

D$^2$GS（2025.10）~\cite{2510_25173} 提出无需 LiDAR 的城市重建框架，动机在于实际场景中 LiDAR 与相机的精确时空标定困难且安装位置差异引入重投影误差。方法首先利用预训练多视角深度网络预测各视角度量深度图并反投影为密集点云（约一亿点），通过渐进剪枝策略——基于不透明度阈值逐轮淘汰不重要高斯并重置不透明度——将点云压缩至可管理规模，保持全局一致性。优化过程中设计多视角深度增强器（Multi-view Depth Enhancer），利用 Marigold 深度基础模型的扩散先验迭代细化渲染深度图：通过置信度图识别高置信区域，结合参考损失、多视角翘曲（warping）损失和边缘感知平滑损失联合引导扩散步骤，每 $N$ 轮迭代更新深度监督池。此外引入道路节点对路面高斯施加位置约束（限制 $z$ 坐标接近预计算平均高度）和平面约束（法向对齐地面、最小缩放通道趋近零），在 Waymo 上即使不使用 LiDAR 也超越使用真值 LiDAR 的方法。

\paragraph{道路扁平与天空无穷远的几何约束}

AutoSplat（2024.07）~\cite{2407_02598} 对道路和天空区域的高斯施加几何约束实现多视角一致的车道变换模拟。道路高斯被约束为扁平的各向异性椭球（垂直缩放接近零），天空高斯被推到无穷远并通过环境贴图渲染。利用 3D 模板引入反射高斯一致性约束监督前景物体可见和不可见面——可见面的高斯通过光度损失优化，不可见面的高斯通过模板的对称性约束间接正则化。并估计每个前景高斯的残差球谐以建模动态外观。

\paragraph{道路-环境解耦的深度有序混合渲染}

DHGS（2024.07）~\cite{2407_16600} 提出解耦混合像素级混合器，将道路和非道路层用独立高斯模型分别建模（道路层采用 2DGS，环境层采用 Scaffold-GS），而非使用统一可微分渲染处理整个场景。道路几何通过预训练 SDF 隐式表示提供精细监督：利用 LiDAR 道路点云拟合符号距离函数，通过距离损失、法线损失（基于 SVD 估计邻域法向）和 Eikonal 正则化约束道路高斯的中心位置和朝向。深度有序混合渲染策略是该方法的核心创新——根据道路与环境模型渲染深度的差值通过 sigmoid 函数生成连续融合权重（替代硬切换），确保两层的叠加在边界处保持一致性和连续性，配合透射率损失（约束累积透射率与语义分割掩码一致）和深度一致性损失（惩罚边界区域的深度不匹配），在 Waymo 上实现几乎不可见边界的缝合效果。

\paragraph{孪生编码器多模态联合学习的零样本泛化}

ADGaussian（2025.04）~\cite{2504_00437} 强调图像和深度特征联合优化对准确高斯预测的重要性。动机来自对 DepthSplat 等方法的观察：其图像和深度特征在并行分支中独立处理，缺乏信息共享和同步优化，导致跨数据集精度不一致。方法引入稀疏 LiDAR 深度作为额外输入模态，设计孪生 ViT 编码器以权重共享方式联合处理 RGB 图像和深度图，通过交叉注意力 Transformer 解码器实现多模态特征的双向增强；并提出深度引导位置嵌入（DPE），将空间坐标与对应深度值拼接形成 $xy$-$z$ 联合位置编码，增强 3D 空间感知。解码端采用双头架构——高斯头（Gaussian Head）预测外观参数、几何头（Geometry Head）预测几何参数——均在 DPT 解码器各尺度注入多尺度深度编码，在 Waymo 和 KITTI 上展现优越的零样本泛化能力。

\paragraph{仿射变换序列的统一鱼眼渲染}

UniGaussian（2024.11）~\cite{2411_15355} 与 ParkGaussian 针对鱼眼相机使用无迹变换（Unscented Transform）不同，UniGaussian 提出基于仿射变换序列的可微分鱼眼渲染方法，统一建模多种传感器（针孔/鱼眼相机）和模态（深度、语义、法向、LiDAR 点云）。

\textbf{鱼眼相机的高斯变换。} 核心观察是鱼眼相机模型（Kannala-Brandt 和 MEI）的数学本质是将相机-点向量向光轴旋转后再投影——这与 3DGS 的渲染流程结构相似，只需将 3D 点替换为 3D 高斯并引入变形。具体实现分三步：（1）位置和姿态调整——将高斯中心沿旋转轴 $\mathbf{r}_{rot}$ 旋转角度 $\Delta\theta = \theta_d - \theta$（$\theta_d$ 由鱼眼模型计算），四元数同步更新 $q_g' = q \otimes q_g$；（2）极向和切向压缩——由于鱼眼压缩视场角，高斯需要相应压缩以避免泼溅平面上过度重叠，极向压缩比 $k_\theta = d\theta_d / d\theta$ 通过泰勒展开从鱼眼模型导数计算，切向压缩比 $k_\phi = \theta_d / \theta$；（3）缩放因子和旋转更新——压缩后的协方差 $\Sigma' = S_\theta \Sigma S_\phi^T S_\phi \Sigma S_\theta^T$ 通过特征值分解提取新的缩放因子和四元数。整个过程完全可微分，且不需额外 GPU 显存，适合处理驾驶场景重建中的百万级高斯。

\textbf{多模态统一渲染。} 每个高斯附加语义 logits、法向等属性，通过 $\alpha$-混合 在统一瓦片光栅化器中渲染为深度图、语义图和法向图。LiDAR 仿真通过从渲染深度图中按真实 LiDAR 参数和扫描模式提取点云实现。MCMC 采样和重定位策略替代标准自适应密度控制，鼓励移除冗余高斯并产生更紧凑的场景表示。

\paragraph{其他街景重建方法}

多个工作从不同角度丰富了街景重建工具箱。MagicRoad（2025.07）~\cite{2507_23340} 通过语义引导颜色增强和障碍物修复实现鲁棒道路表面重建。GS-RoadPatching（2025.09）~\cite{2509_19937} 通过 3D 空间结构匹配实现驾驶场景的替代式修复。MOGS（2025.09）~\cite{2509_06685} 用单目深度替代 LiDAR 降低传感器依赖。LT-Gaussian（2025.08）~\cite{2508_01704} 提出高斯地图长期更新方法。LS-Gaussian（2025.07）~\cite{2507_21572} 实现轻量级流式渲染的算法-硬件协同设计。CE-NPBG（2025.04）~\cite{2504_19557} 通过连接关系图改善大规模驾驶场景中点云渲染的可扩展性。MuDG（2025.03）~\cite{2503_10604} 将多模态扩散模型与高斯泼溅结合增强极端视角变化下的渲染鲁棒性。Zheng et al.（2025.03）~\cite{2503_12001} 提出自适应透明度机制消除运动物体。ReconDreamer（2024.11）~\cite{2411_19548} 通过在线修复和渐进数据更新实现大偏移渲染。

MagicDrive3D（2024.05）~\cite{2405_14475} 首先训练多视角视频生成模型再构建 3DGS，支持多条件控制。SGD（2024.03）~\cite{2403_20079} 利用扩散先验和 LiDAR 深度增强 3DGS 的街景外推能力。StreetUnveiler（2024.05）~\cite{2405_18416} 从拥挤观测中恢复空街道。FaithFusion（2025.11）~\cite{2511_21113} 通过像素级期望信息增益融合 3DGS 和扩散模型。Drone-assisted Road GS（2024.08）~\cite{2408_15242} 引入跨视图不确定性使无人机数据辅助地面视角合成。LR-SGS（2026.03）~\cite{2603_12647} 利用 LiDAR 反射率引导显著高斯表示增强复杂光照场景。Omran et al.（2025.10）~\cite{2510_12308} 提出高斯泼溅与扩散模型混合的两阶段新视角合成方案。Hybrid Gaussian（2025.03）~\cite{2503_08217} 和 Hybrid Rendering（2025.03）~\cite{2503_09464} 分别探索了混合表示和混合渲染策略。

\begin{takeaway}
街景重建与渲染的核心突破在于鱼眼相机高斯变换和多模态统一渲染，使 3DGS 能同时建模针孔相机、鱼眼相机和 LiDAR 等多种传感器，为自动驾驶仿真提供完整的场景表示基础设施。
\end{takeaway}

%=====================================================================

\subsection{近期方法脉络：街景重建与仿真压缩、街景训练增强}

这一组近期工作延续 3DGS，主要把改动集中在压缩、增量改进等环节。下面只保留相对既有路线的变化，便于和前文方法直接对照。

\paragraph{SparseStreet}
SparseStreet（2026.06）~\cite{2606_03909} 延续 3DGS 在“街景重建与仿真压缩”上的基本设定，主要改变是：SparseStreet：节点级可学习剪枝保留关键动态区域，并在场景稳定后进行背景压缩。高斯基元仍承担可渲染动态场景表示。
其设计动机是：提出针对街景结构的通用压缩框架，利用动态对象和静态背景的不同重要性分配高斯预算，在保持时序一致性的同时降低存储与渲染成本。
具体而言，方法先通过node-based learnable pruning删除低贡献高斯并保护车辆、行人等视觉关键区域；待表示稳定后再对背景进行进一步压缩，减少静态区域冗余。
适用边界是：依赖动态/静态区域划分；极端复杂街景和长序列中压缩误差可能累积。

\paragraph{Leveraging NeRF-Rendered Images for 3D Gaussian Splatting}
Leveraging NeRF-Rendered Images for 3D Gaussian Splatting（2026.06）~\cite{2606_09034} 延续 3DGS 在“街景训练增强”上的基本设定，主要改变是：用预训练街景NeRF生成额外训练图像和BEV视角，再结合扩散增强训练目标3DGS。高斯基元仍承担可渲染动态场景表示。
其设计动机是：提出NeRF-to-3DGS训练增强策略，将NeRF较高质量渲染和3DGS实时性结合，尤其针对街景中的瞬态对象和缺失视角。
具体而言，方法使用街景专用NeRF渲染目标3DGS训练图像，移除输入视图中的瞬态对象，并生成鸟瞰补充视角；扩散增强进一步改善额外视图质量。
适用边界是：依赖预训练街景NeRF；额外视图错误会传递给3DGS；增加预处理成本。
\section{仿真与数据生成}
\label{sec:18_2}

自动驾驶仿真系统的核心价值在于能够生成海量、多样、安全关键的训练和测试数据。3DGS 以其显式场景表示和实时渲染能力，为构建高保真自动驾驶仿真器提供了新的技术基础。本节从闭环仿真、场景编辑与操控、多模态数据合成和预训练范式四个维度梳理相关工作。

\subsection{闭环仿真与场景编辑}
\label{ss:18_2}

本节讨论2条技术路线：（1）闭环仿真平台；（2）场景编辑与操控。

\subsubsection{闭环仿真平台}

本路线介绍基于 3DGS 的完整闭环驾驶仿真平台，实现感知-规划-控制循环测试的实时照片级渲染。

闭环仿真要求仿真环境能够根据自车的控制指令（转向角、加速度、制动力）动态更新场景中所有参与者的状态，并实时渲染新的传感器观测，实现完整的感知-规划-控制循环测试。相比开环评测（仅在固定视角下评估渲染质量），闭环仿真对渲染速度（需实时）、场景动态性（需支持自车运动和其他车辆响应）和物理合理性（需一致的几何和光照）提出了更高要求。

\paragraph{高斯场景图驱动的实时闭环仿真器}

HUGSIM（2024.12）~\cite{2412_01718} 是首个基于 3DGS 的完整闭环驾驶仿真平台。其核心突破在于将 2D RGB 图像重建为可交互的 3D 高斯场景，并动态更新自车和其他参与者的状态与观测。直觉上，这相当于把 2D 视频游戏升级为 3D 可交互世界——不仅能看到场景，还能在其中自由驾驶并观察所有参与者的实时响应。

\textbf{Baseline 与问题。} 闭环仿真要求仿真器根据自动驾驶系统的控制指令（转向、加速、制动）实时更新场景状态，生成新的传感器观测，形成完整的感知-规划-控制循环。传统 NeRF 仿真器渲染速度太慢（秒级/帧），无法支持实时闭环测试。CARLA 等游戏引擎仿真器虽然实时，但网格资产的建模成本高昂，渲染保真度与真实传感器数据存在显著 sim-to-real 差距。亟需一种兼具照片级保真度和实时渲染速度的场景表示。

\textbf{改了什么。} HUGSIM 将场景表示从 NeRF 的隐式辐射场或游戏引擎的网格资产，替换为 3DGS 的显式高斯场景图——静态背景、动态物体和天空分别建模，统一泼溅渲染。

\textbf{场景表示与渲染。} HUGSIM 将城市场景建模为包含静态背景、动态物体和天空的 3D 高斯场景图。静态背景用标准 3D 高斯表示，动态车辆在物体局部坐标系中建模后通过位姿变换到世界坐标，天空用环境贴图表示。利用 3DGS 的实时渲染能力实现闭环仿真中的连续帧生成，同时解决视角外推（自车偏离训练轨迹时）和 360 度车辆渲染（环绕视角观测动态车辆）两大挑战。

\textbf{设计动机。} 3DGS 的显式点云表示使场景编辑天然可行——动态物体可独立编辑（移动、删除、插入）而不影响背景，且泼溅渲染管线对场景修改实时响应。NeRF 的隐式表示无法实现这一点：修改场景需要重新训练整个辐射场。同时，高斯基元的世界坐标与物理引擎兼容，可直接从高斯位置提取碰撞检测所需的几何信息。

\textbf{性能与基准。} HUGSIM 在 KITTI-360、Waymo、nuScenes、PandaSet 的 70+ 序列和 400+ 变化场景上建立综合基准，支持照片级真实感闭环设置下的算法微调，验证了 3DGS 仿真环境对下游感知和规划算法的训练效果。

\paragraph{Unity 引擎内嵌高斯资产的混合仿真器}

GSAVS（2024.12）~\cite{2412_18816} 修改了经典自动驾驶仿真器的场景表示环节——用 3D 高斯泼溅资产替代传统网格（mesh）资产——物理引擎和碰撞检测机制保持不变。与 HUGSIM 从零构建闭环渲染管线不同，GSAVS 的设计思路是在现有 3D 引擎（Unity）内部嵌入全高斯资产，直接复用引擎的物理引擎、碰撞检测和强化学习接口，同时用 3DGS 的照片级渲染替代网格渲染。

\textbf{Baseline 做法。} 经典自动驾驶仿真器（如 CARLA）将环境、车辆、行人分别建模为网格资产，使用标准 3D 文件格式（如 .fbx）表示场景。网格方式的优势在于定制灵活——可以自由修改光照、交通规则、交叉口参数——但代价是制作照片级数字孪生耗时耗力，且渲染保真度难以匹配真实传感器数据，导致 sim-to-real 差距显著。

\textbf{改了什么。} GSAVS 将场景中的全部资产（环境、自车、其他车辆、行人）替换为 3D 高斯泼溅资产，渲染仍由 Unity 引擎实时完成。具体而言，环境高斯资产通过标准 3DGS 训练流程从 nuScenes 多视角数据生成，再通过 UnityGaussianSplatting 包导入 Unity；车辆高斯资产在封闭场景中采集并训练，经 SuperSplat 手工清理后导入。关键转变在于：3DGS 的输入数据采集遵循标准流程（多视角 RGB 图像 + SfM 点云），而非手工建模。

\textbf{实现原理。} 由于 3DGS 资产本身不携带物理属性（碰撞体、摩擦力等），GSAVS 采用"高斯渲染 + 网格碰撞"的混合架构。自车高斯资产外包一个大尺寸盒状碰撞体（box collider）用于检测与墙壁和其他车辆的碰撞，四个独立的车轮碰撞体（wheel collider）放置在与车轮匹配的位置，接收扭矩向量和转向角输入。物理交互通过隐形道路轨道（Road Spline Track）实现：利用 3DGS 训练过程中产生的相机外参，在相机位置 $(X, Y, Z)$ 序列上构造样条曲线；沿样条线以频率 $f$ 等间距放置预制的道路块资产（RoadBlockAsset），每个道路块宽度约为自车宽度的 2 倍、墙壁高度约为自车高度的 2 倍，方向对齐相机外参的上方向量和前进方向，所有道路块向下偏移自车高度的一半以确保前相机与样条线对齐。网格渲染器被关闭，仅保留网格碰撞体，形成一条隐形轨道。

\textbf{设计动机。} 这一混合架构的动机来自 3DGS 在大范围无界场景中的固有局限：高斯泼溅无法直接生成物理精确的几何，因此不能像网格资产那样直接用于碰撞检测和物理模拟。样条轨道的约束区域恰好对应原始数据采集车辆的行驶路径——在此区域内，高斯泼溅的新视角合成质量最高，因为输入数据提供了充分的视角覆盖。偏离此区域越远，渲染质量越差。因此，将自车约束在样条轨道附近既是物理交互的需要，也是渲染质量的保证。

\textbf{实验结果。} 在 nuScenes 数据上构建仿真环境后，使用 Unity ml-agents 的 PPO 策略训练 250,000 个 episode，评估三项任务：小场景直行（无动态障碍）准确率超过 80\%，大场景转弯准确率约 70\%，小场景带动态车辆避障准确率同样超过 80\%。资源利用方面，场景面积翻倍时 GPU 利用率和显存增长不到两倍，添加动态车辆后额外开销不超过 3\%，平均帧率基本保持不变——这得益于 3DGS 点云资产比同等网格资产占用更少空间。

\paragraph{非对称混合强化学习的端到端泊车}

REAP（2026.05）~\cite{2605_08713} 提出基于强化学习的端到端自动泊车方法，使用 3DGS 构建 Real2Sim2Real 仿真器。与 GSAVS 在 Unity 引擎内部嵌入高斯资产不同，REAP 的 3DGS 仿真器聚焦泊车场景的独特需求——通过手持 LiDAR 扫描仪采集地下车库数据并用精确点云初始化 3DGS 加速训练和提升几何精度，结合 2.5D 高程图和双层车辆模型支持坡道、减速带和机械车位等复杂地形仿真。Real2Sim 步骤用 3DGS 将真实场景转化为数字场景，Sim2Real 步骤将端到端模型部署到物理车辆。核心设计是非对称混合强化学习策略——actor 直接从环视图像学习策略，critic 使用仿真器提供的无伪影 BEV 特权信息评估价值函数，动态混合策略在探索过程中实时查询规则规划器（Reed-Shepp 曲线）生成在线专家监督而非依赖预收集的离线数据集，配合渐进式引导丢弃策略逐步从专家监督过渡到自主决策。在极端狭窄机械车位中验证了端到端强化学习泊车的可行性。

\paragraph{纠错回放缓冲区的样本高效泊车学习}

ParkingWorld（2026.05）~\cite{2605_25029} 提出纠错回路样本高效强化学习（CIL-SERL）框架，在 3DGS 泊车仿真器中训练。与 HUGSIM 面向开放道路闭环仿真不同，ParkingWorld 聚焦极端狭窄泊车场景——特别是机械车位——通过纠错机制大幅提升样本效率。多级回放缓冲区（"错误笔记本"）分层组织标准 RL 回滚、人工纠正干预、失败探索轨迹和基于回滚的纠正片段：每个纠错区域显式配对"策略失败片段"与"人工纠正片段"，使得 critic 和 actor 在每个 mini-batch 中同时看到错误行为及其修正，而非简单混合所有数据。3DGS 仿真器通过 XGRIDS 手持设备采集地下车库数据重建高保真数字孪生，支持碰撞检测、高程感知状态更新和状态回滚。在 3DGS 仿真和物理车辆上均取得显著提升。

\paragraph{预生成合成观测的伪仿真范式}

SceneCrafter（2025.03）~\cite{2503_18108} 提出"伪仿真"（pseudo-simulation）范式，使用 3DGS 在真实数据集上预生成多样化的合成观测（位置、航向、速度变化），通过基于邻近度的加权方案近似自动驾驶系统可能遇到的未来状态，与闭环仿真的相关性（$R^2=0.8$）优于最佳开环方法（$R^2=0.7$）。Pseudo-Simulation（2025.06）~\cite{2506_04218} 同样利用 3DGS 预生成合成观测建立公共排行榜，系统验证了伪仿真与闭环仿真之间的高相关性。

\subsubsection{场景编辑与操控}

本路线介绍基于 3DGS 的场景编辑能力，支持插入/删除/移动车辆、物理仿真碰撞和多样化交通场景生成。

场景编辑能力使仿真器能够生成超越原始采集数据的多样化场景，包括插入/删除/移动车辆、修改轨迹、改变天气等。这是闭环仿真中最具应用价值的环节：安全关键场景（如突然切入、行人横穿）在真实路测中出现概率极低，通过编辑已有场景可以批量生成这些罕见但致命的测试案例。核心挑战在于编辑后的场景需要保持时空一致性（插入的车辆不应穿透其他物体）和视觉真实感（阴影和光照需要与场景匹配）。

\paragraph{高斯基元即物质点的物理仿真编辑}

Real2Sim（2026.05）~\cite{2605_13591} 将 4DGS 与可微分物质点方法（Material Point Method, MPM）求解器结合，是首个实现物理感知驾驶场景编辑的工作。

\textbf{核心创新——高斯基元即物质点。} Real2Sim 的关键洞察是将每个 3D 高斯基元直接视为 MPM 求解器中的物质粒子，携带质量 $m_p$、初始体积 $V_p^0$、位置 $x_p$、速度 $v_p$ 和变形梯度 $F_p$。在每个时间步中，粒子量通过 $C^1$ B 样条核函数传递到欧拉网格，更新网格速度后传回粒子，实现变形仿真。高斯的均值和协方差随物理仿真自然演化：$\mathbf{x}_p(t) = \phi(X_p, t)$，$\mathbf{A}_p(t) = F_p(t) \cdot \mathbf{A}_p \cdot F_p(t)^T$，无需网格提取或重网格化。

\textbf{场景构建与编辑。} 采用 Street Gaussians 方法论分别建模静态背景和动态物体，背景用 LiDAR + SfM 初始化，动态物体用专用高斯簇和 4D 傅里叶 SH 建模。用户可以自由重新定位、旋转车辆或插入外部 3D 资产来多样化交通场景，随后通过 MPM 求解器自动仿真碰撞和碰撞后轨迹。

\textbf{设计动机。} 统一 3DGS 的渲染管线和 MPM 的物理引擎——变形后的高斯基元直接用原始泼溅管线渲染，避免了传统方法中物理仿真到渲染的格式转换开销。不透明度和颜色系数在仿真过程中保持不变，确保物理变形与外观一致。

\textbf{应用。} 生成包括碰撞在内的具有挑战性的极端场景，为下游感知和控制任务提供物理真实的测试数据。

\paragraph{优化无关的单图像道路几何插入}

RoVES（2026.05）~\cite{2605_25373} 与 Real2Sim 同样追求物理感知的编辑，但 Real2Sim 通过 MPM 物理仿真处理碰撞变形，RoVES 聚焦于道路-车辆交互——支持单图像驱动的道路几何插入（减速带、下沉路面），并耦合车辆动力学模型实现物理一致的姿态响应。

\textbf{几何-姿态解耦框架。} 与传统编辑管线需要为每次修改迭代优化高斯参数不同，RoVES 采用优化无关策略——基于 Street Gaussians 重建冻结所有高斯参数，仅保留逐帧物体位姿独立可编辑。编辑操作解耦为静态分支（一次性前馈完成道路几何插入，无需梯度回传）和动态分支（通过更新位姿序列实现平移、旋转、删除等操作，渲染时将物体局部坐标系的高斯按新位姿变换到世界坐标）。

\textbf{单图像道路几何插入。} 使用 Grounding DINO + SAM 级联从文本提示（如"减速带""下沉路面"）定位并分割目标物体，Depth Anything 预测深度后反投影为局部 3D 点云，根据预设真实物理尺寸（$L_x, L_y, L_z$）归一化缩放。关键细节处理：引入全局缩放因子 $\beta = 0.01$ 修正单目深度点云过密导致的高斯重叠和颜色混合；统计颜色迁移在 Lab 空间匹配均值和方差消除插入物体与周围路面的色差；高斯基元直接合并到静态背景集合中完成插入，整个过程无需重训练。

\textbf{4 自由度半车动力学耦合。} 道路几何编辑后，构建高斯密度高度场——将编辑后的道路高斯投影到地面网格，对凸起结构取最大残差、凹陷结构取最小残差。经典 4-DOF 半车动力学模型（簧载质量垂直位移 $z_s$ 和俯仰角 $\theta$、前后非簧载质量垂直位移 $z_{uf}$ 和 $z_{ur}$）以前后接触点高度 $z_{rf}, z_{rr}$ 为激励输入，通过 RK4 数值求解微分方程。解出的车身垂直位移和俯仰角直接叠加到 3DGS 车辆位姿上——车辆经过新插入的减速带时自动产生合理的颠簸响应。一次性优化无关管线在 1.84 秒内完成道路元素插入，全管线 6.24 秒。

\paragraph{语言驱动车辆插入与视频扩散精化}

HorizonForge（2026.02）~\cite{2602_21333} 与 RoVES 的单图像道路编辑不同，聚焦于语言驱动的车辆插入和多智能体轨迹控制。核心思路是将场景重建为高斯泼溅和网格混合表示，通过视频扩散模型精化渲染结果。

\textbf{3D 资产采集与对比。} 系统对比了三种 3D 表示作为编辑条件：彩色 LiDAR 点云（StreetCrafter）、3D 边界框和高斯泼溅。实验表明高斯泼溅在保持精细外观细节方面显著优于稀疏 3D 基元。为支持新车辆插入，额外从视频帧中采集网格资产——对每个物体选择最佳观测帧（包含最多 LiDAR 点），经 SAM 分割和 Pix2Gestalt 非模态修复后，用 Hunyuan3D 生成粗网格，再通过深度差异和 IoU 对齐项联合优化尺度参数 $s^* = \arg\min_s \|D(M,s) - D_{gt}\| - \lambda \text{IoU}(B(M,s), B_{gt})$ 将网格与真值 3D 边界框对齐。推理时，GPT 从文本描述生成参考图像再转为网格资产。

\textbf{噪声感知视频扩散精化。} 将编辑后的 3D 场景渲染为 2D 视频后，利用视频扩散模型（基于 TrajectoryCrafter + CogVideoX 主干）精化渲染伪影。训练数据通过循环重建策略构造——对原始轨迹 $T$ 施加扰动 $T' = T + \Delta T$ 生成退化渲染，再从退化渲染训练新的高斯场并渲染回原始轨迹形成干净-噪声配对。关键发现是视频扩散模型比图像扩散模型在时序一致性和伪影消除方面有显著优势。单次前馈无需逐轨迹优化，用户偏好率比第二名方法提升 83.4\%。

\paragraph{外观-运动-位置三层语言对齐编辑}

SIMSplat（2025.10）~\cite{2510_02469} 与 HorizonForge 的轨迹驱动编辑不同，SIMSplat 实现了自然语言直接操控高斯场景的完整管线——从查询、插入到轨迹修改均可通过语言提示完成。

\textbf{三层语言-高斯对齐。} 外观对齐：对每个物体提取 CLIP（Contrastive Language-Image Pre-training，对比语言-图像预训练）特征后通过轻量自编码器压缩为潜在表示 $h_t(o) = E(f_t^{CLIP}(o)) \in \mathbb{R}^d$，作为每个高斯外观特征 $c_i^{app}$ 的自监督目标。时间对齐：现有方法（如 LangSplat）仅嵌入静态语义特征，无法理解交通场景中特有的运动描述（如"从右向左穿越路口的车辆""在自车左侧横穿马路的行人"）。SIMSplat 设计了轨迹编码器 $E_{traj}$ 将物体轨迹 $X = \{(x_t, y_t)\}_{t=1}^T$ 编码为潜在表示，通过运动码本和位置码本（分别存储典型运动模式如"左转""从右到左"和位置描述如"在自车前方""在自车左侧"）的凸组合生成时间语言特征 $c_o^{temp} = (f^{motion}, f^{location})$。高斯对象查询时，先通过外观相似度缩小候选范围，再结合运动和位置特征选择最匹配的对象。

\textbf{多智能体路径精化。} LLM agent 解析用户提示后生成轨迹编辑，但 LLM 缺乏路径规划能力。SIMSplat 训练 SMART-1B 运动预测模型——以编辑后目标物体的轨迹和所有智能体的历史观测为输入，预测所有智能体的未来轨迹——使编辑后的场景中周围车辆会自然响应（如前方车辆急刹车时后车绕行或停止）。最终通过多轮对话和扩散修复模型精化编辑区域的视觉效果。

\paragraph{高斯渲染条件驱动的对象级姿态编辑}

$\text{G}^2$Editor（2025.08）~\cite{2508_20471} 与 SIMSplat 的语言驱动编辑不同，$\text{G}^2$Editor 聚焦于对象级精确姿态控制和外观保持——核心问题是如何在编辑车辆位置/朝向后维持原始外观细节。

\textbf{3DGS 渲染作为姿态条件。} 现有方法直接将参考图像和 3D 边界框输入扩散模型，要求模型理解 3D 场景几何和视角变换——这对扩散模型是极其困难的。$\text{G}^2$Editor 的关键洞察是将编辑对象的 3DGS 资产按目标位姿渲染为高斯视频 $V_g$（通过 $\alpha$-rendering 在目标视角下渲染每个高斯的投影图像），以此作为条件输入替代原始参考图像。扩散模型的任务从"理解 3D 几何并变换视角"简化为"从略微模糊的高斯渲染恢复清晰图像"——类似高斯图像修复任务，大幅降低训练难度。高斯渲染本身通过可微分泼溅显式保证 3D 空间一致性。

\textbf{场景级布局与外观维护。} 使用场景级 3D 边界框（而非仅目标物体框）保持非编辑物体的布局不变，通过深度感知框编码（投影角点深度并插值生成深度图）和边缘掩码注入扩散 UNet。外观保持采用 ReferenceNet（与原始 SD 共享架构）提取参考图像特征，通过空间注意力层选择性地注入相关信息。关键技巧是参考图像的随机水平翻转——防止扩散模型从 ReferenceNet 泄露姿态信息，确保姿态控制仅来自高斯渲染条件。两阶段训练策略：第一阶段学习单图像的对象绘制，第二阶段通过时间注意力层增强视频时序一致性。统一支持对象重定位、插入和删除。

\paragraph{图像-LiDAR 跨模态联合编辑}

MultiEditor（2025.07）~\cite{2507_21872} 与 $\text{G}^2$Editor 仅编辑图像不同，MultiEditor 首次实现图像和 LiDAR 点云的联合编辑——确保插入或移动的车辆在两种模态中保持 3D 一致性。

\textbf{双分支扩散架构。} 图像分支和 LiDAR 分支各自拥有独立的去噪模型，但共享 3DGS 作为目标物体的结构和外观先验。图像分支采用像素级粘贴（将 3DGS 渲染的目标车辆图像直接粘贴到 ROI 区域作为硬约束，经 VAE 编码后拼接到潜变量）、语义级引导（CLIP 图像编码器提取全局语义特征——CLIP 将图像和文本映射到共享嵌入空间——通过交叉注意力注入）和双分支精化（重建分支学习语义对齐，精化分支在粘贴图像上增强细粒度细节）三层控制机制。LiDAR 分支将点云转换为距离图表示，利用 3DGS 的深度渲染能力合成密集深度图作为空间先验（训练时从稀疏 LiDAR 插值，推理时直接用 3DGS 渲染），替代难以获取的目标物体多姿态 LiDAR 数据。

\textbf{深度引导跨模态条件模块。} 核心挑战是相机深度估计误差导致图像-点云的空间对应关系不准确。MultiEditor 设计了深度引导可变形条件模块——先用几何投影建立粗模态对齐（距离图像素 $(u,v)$ 通过相机内参和 LiDAR-相机外参变换计算），再用可变形交叉注意力从互补模态自适应采样局部特征。Camera-to-LiDAR 方向：在距离图像位置 $(u,v)$ 采样图像潜变量特征，通过可变形注意力生成 LiDAR 分支的控制信号 $h_r^c$；LiDAR-to-camera 方向反之。两个方向均使用零初始化门控（$\alpha_r, \alpha_c$），训练初期跨模态影响被抑制，逐步学习有效的跨模态一致性。

\paragraph{车路协同视角的协作重建与编辑}

CRUISE（2025.07）~\cite{2507_18473} 与 MultiEditor 的单车单模态编辑不同，CRUISE 面向车路协同（V2X）场景——同时从自车和基础设施视角重建、编辑和渲染，生成大规模 V2X 数据增强训练集。

\textbf{V2X 专用重建适配。} 直接将 Street Gaussians 应用于 V2X 数据会遇到独特问题：基础设施视角的标注包含自车边界框，在自车视角训练时该框覆盖大部分可见场景，导致静态背景被错误重建为动态高斯。CRUISE 移除自车框并引入自车掩码遮蔽自车区域，实现动静正确分离。多帧 LiDAR 点云融合（自车 + 基础设施）提供稳健初始化，跟踪框聚合多帧点云增强物体高斯密度。训练损失包含 8 项——颜色、深度、法向、天空、语义、尺度、比率比和浮动正则化——联合约束几何和外观。

\textbf{GPT 驱动的交通流编辑。} 收集互联网多视角车辆图像和尺寸信息，通过 TRELLIS 生成高质量 3D 高斯车辆资产。编辑时简化向量地图（BFS 扩展交叉口相关车道），将车道中心线的起止位置和方向提供给 GPT-4o，结合自车轨迹采样生成合理的多车插入轨迹。生成的高斯车辆资产按轨迹定位后渲染来自两个视角的同步图像和 3D 标注，构建完全标注的合成 V2X 数据集。在 3D 检测和协同 3D 跟踪任务上验证了生成数据对下游任务的有效性。

\subsection{数据合成与增强}
\label{ss:18_3}

本节讨论2条技术路线：（1）多模态数据合成与生成；（2）数据增强与预训练。

\subsubsection{多模态数据合成与生成}

本路线介绍利用视频扩散模型和 3DGS 生成多模态传感器数据（相机图像、LiDAR 点云等）的方法。

自动驾驶仿真需要同时生成相机图像、LiDAR 点云、雷达数据等多模态传感器输出。单一模态的仿真不足以验证完整的自动驾驶系统：视觉感知依赖相机图像，激光感知依赖 LiDAR 点云，传感器融合算法需要两种数据的精确时空对齐。视频扩散模型的引入使得从有限的真实数据生成大量多样化的多视角驾驶视频成为可能；3DGS 则提供将这些视频提升为完整 4D 场景的几何框架。

\paragraph{单步确定性扩散的在线渲染增强}

DiffusionHarmonizer（2026.02）~\cite{2602_24096} 是首个将预训练图像扩散模型转化为确定性单步在线增强器的工作，用于在仿真循环中实时修复 3DGS 重建渲染的缺陷。

\textbf{核心创新——单步确定性增强。} 标准扩散模型需要多步去噪（通常 20--50 步），无法满足在线仿真的实时需求。DiffusionHarmonizer 将预训练扩散模型改造为确定性单步增强器：直接将清洁潜变量 $E_\theta(I_t)$ 输入网络（不注入噪声），时间步和文本条件标记固定为常量"空"值，产生从输入潜变量到增强潜变量的稳定确定性映射。为保持时间一致性，骨干网络 $F_\theta$ 扩展了时间注意力层，以当前帧和 $K=4$ 帧历史增强帧作为上下文输入。

\textbf{定制数据管线。} 由于缺乏公开的合成-真实配对数据，设计了五个互补的数据合成组件：（1）新视角伪影校正——通过稀疏重建、循环重建和模型欠拟合生成降质渲染；（2）ISP 修改——随机化色调映射、曝光、白平衡等参数模拟前景-背景外观不匹配；（3）重光照——用扩散重光照模型在随机光照条件下重新生成选定区域；（4）基于物理的阴影仿真——用物理渲染器在可控光照下合成精确阴影；（5）资产重新插入——将动态前景重新插入 3DGUT 重建的背景中但不生成正确阴影。

\textbf{训练策略。} 引入多尺度感知损失（随机采样 $[128, 512]$ 像素的方形块计算 VGG 特征距离）抑制单步训练导致的棋盘格伪影；光流扭曲损失保持时间一致性。混合时间/非时间批次训练防止过度依赖历史帧。可在单 GPU 上在线运行。

\paragraph{高斯泼溅解码器的跨视图视频扩散}

CVD-STORM（2025.10）~\cite{2510_07944} 提出跨视图视频扩散模型，核心创新在于设计了 STORM-VAE——在预训练图像 VAE（SD3.5）基础上增加并行的 Gaussian Splatting 解码器分支，使 VAE 潜变量同时服务于图像重建和高斯泼溅 3D 重建，从而在潜在空间中隐式编码场景几何。在此基础上，训练了一个基于 DiT 的扩散模型，采用三种互补的 transformer block 分别沿空间维度（MM-DiT）、时间维度和跨视图维度执行注意力，以单阶段联合训练替代多阶段渐进训练，在生成多视角驾驶视频的同时直接输出动态 3DGS 表示（含高斯速度场），实现几何与外观的同步生成。

\paragraph{几何不确定性引导的立体强制 4D 生成}

PhiGenesis（2025.09）~\cite{2509_20251} 提出统一 4D 场景生成框架，第一阶段利用预训练视频 VAE 配合范围视图（range view）适配器实现前馈 4D 重建——将 LiDAR 距离图的视频编码能力迁移到相机图像的 4D 重建任务中。第二阶段引入几何引导视频扩散模型生成未来视角。核心创新"立体强制"（Stereo Forcing）策略在去噪过程中集成几何不确定性——对于几何不确定的区域（如远处弱纹理区域），降低扩散生成的影响以避免引入不一致的几何。

\paragraph{视频扩散先验与自监督动静分解的 4D 建模}

DreamDrive（2025.01）~\cite{2501_00601} 修改了 3DGS 的场景表示和初始化管线——用混合高斯表示（静态时间无关 + 动态时间相关）替代单一时间相关高斯，高斯泼溅渲染函数 $F_\mathrm{splat}$ 保持不变。核心思路是将视频扩散模型的生成能力与 3DGS 的几何一致性渲染相结合，沿着 2D $\to$ 3D $\to$ 4D 的渐进路径构建驾驶场景。

\textbf{Baseline 与问题。} 标准 3DGS 场景重建依赖精心标注的驾驶日志（标定相机位姿、物体边界框、3D 点云），限制了在无标注野外数据上的扩展性。另一条路线——纯视频扩散模型——虽然可以从任意驾驶视频生成动态视频，但因缺乏 3D 几何约束，跨帧几何一致性差，生成的视频无法可靠用于自动驾驶下游任务。Deformable GS 等方法将全部高斯建模为时间相关，对生成图像中固有的 3D 不一致性过拟合，导致静态背景在新视角下出现"虚假变形"（fake dynamics）。

\textbf{2D $\to$ 3D：视频扩散先验与高斯初始化。} DreamDrive 的第一个关键步骤是用视频扩散模型（如 Vista、MagicDrive）从输入控制信号（单张图像 $I_\mathrm{ctrl}$ 或地图 $M_\mathrm{ctrl}$）生成 2D 参考图像序列 $I_\mathrm{ref}$ 和潜特征 $Z_\mathrm{ref}$。直觉上，视频扩散模型充当"视觉想象力"——它提供了场景在各时间步的大致外观，但缺乏精确的 3D 几何。

接下来需要将这些 2D 参考提升到 3D。传统做法使用 COLMAP 恢复稀疏点云，但稀疏点不足以表征大尺度无界驾驶场景。DreamDrive 改用端到端多视角立体网络 DUSt3R 生成像素级稠密 3D 点云，同时恢复相机位姿 $P_\mathrm{ref}$。相机内参通过 Weiszfeld 算法估计，外参通过跨帧全局对齐点云计算。聚合后的稠密点云用于初始化 3D 高斯参数 $G_\mathrm{init}$，每个高斯还附带像素级对齐的潜特征 $Z_\mathrm{ref}$：

\begin{equation}
G_\mathrm{init}, P_\mathrm{ref} = F_\mathrm{MVS}(I_\mathrm{ref}, Z_\mathrm{ref})
\end{equation}

即多视角立体网络同时输出初始化高斯集合和相机位姿，为后续 4D 建模提供稠密几何基础。

\textbf{自监督静动态分解。} 有了初始化高斯后，下一步是区分场景中的静态背景（道路、建筑）和动态物体（车辆、行人）。传统方法依赖标注的物体边界框，DreamDrive 则提出纯自监督的分解策略。

核心洞察是：\emph{渲染误差图可以作为动静区域的天然指示器}。具体而言，先将所有高斯 $G_\mathrm{init}$ 按静态假设整体优化，再泼溅为静态图像 $I_\mathrm{static}^t = F_\mathrm{splat}(G_\mathrm{init}, P_\mathrm{ref}^t)$。渲染误差图 $I_\mathrm{err}^t = |I_\mathrm{static}^t - I_\mathrm{ref}^t|$ 中高误差的像素区域，正是静态高斯无法解释的区域——即动态物体所在位置。

训练一个评分网络 $F_\mathrm{score}$ 以误差图为监督信号，为每个高斯输出动静分数 $S = F_\mathrm{score}(G_\mathrm{init}, Z_\mathrm{ref})$。分数通过泼溅投影到图像平面，用二元交叉熵损失优化：

\begin{equation}
\mathcal{L}_\mathrm{dec} = \sum_{t=0}^{T} \mathcal{L}_\mathrm{bce}\bigl(F_\mathrm{splat}(S, P_\mathrm{ref}^t),\, I_\mathrm{err}^t\bigr)
\end{equation}

即分解损失在所有时间步上对泼溅后的动静分数与误差图之间的交叉熵求和，端到端优化评分网络。最后，以阈值 $\tau$ 分割：$G_\mathrm{dynamic}^* = \{g \in G_\mathrm{init} \mid S(g) > \tau\}$，$G_\mathrm{static}^* = \{g \in G_\mathrm{init} \mid S(g) \leq \tau\}$。

\textbf{时空聚类纠错。} 由于生成图像固有的 3D 不一致性，仅凭误差图分割可能将静态物体上的局部变形误判为动态。DreamDrive 引入时空聚类分组策略 $F_\mathrm{group}$——核心假设是同一物体的高斯具有相同的动静属性。将高斯按时空邻近性聚类，若簇中大多数高斯为静态，则整个簇标记为静态（反之亦然），有效纠正误判并消除虚假变形。

\textbf{混合高斯 4D 表示与渲染。} 分解完成后，静态高斯 $G_\mathrm{static}$ 的参数 $G(x, r, s, \alpha, c)$ 在所有时间步保持不变。动态高斯则通过变形网络 $F_\mathrm{deform}$ 预测时间偏移量：$(\Delta x, \Delta r, \Delta s) = F_\mathrm{deform}(x, t)$，得到各时间步的动态高斯 $G_\mathrm{dynamic}^t = G(x + \Delta x, r + \Delta r, s + \Delta s, \alpha, c)$。最终渲染时将两者合并泼溅：

\begin{equation}
I_\mathrm{render}^t = F_\mathrm{splat}\bigl(G_\mathrm{static} \cup G_\mathrm{dynamic}^t,\, P_\mathrm{ref}^t\bigr)
\end{equation}

即将静态高斯和第 $t$ 步的变形动态高斯合并后泼溅到对应相机位姿，渲染损失为 $\mathcal{L}_\mathrm{render} = \sum_{t=0}^{T} \bigl(\mathcal{L}_1(I_\mathrm{render}^t, I_\mathrm{ref}^t) + \mathcal{L}_\mathrm{SSIM}(I_\mathrm{render}^t, I_\mathrm{ref}^t)\bigr)$。

\textbf{设计动机。} 分离动静表示的根本动机是消除生成图像 3D 不一致性带来的虚假动态——静态高斯参数固定不变，天然保证背景在所有时间步和视角下的一致性；变形网络仅需学习真实物体运动的偏移量，大幅降低拟合难度。相比需要物体标注的 Street Gaussians 或全部使用时间相关高斯的 S$^3$Gaussian，DreamDrive 的自监督混合表示在无标注场景下视觉质量提升约 30\%（FID 45.59，FVD 374.02），并能将规划任务的碰撞率降低 25\%。

\paragraph{世界模型合成新轨迹增强 4D 重建}

DriveDreamer4D（2024.10）~\cite{2410_13571} 首次利用视频生成模型改善驾驶场景的 4D 重建，将世界模型作为"数据机器"合成新轨迹视频，结构化条件显式控制交通元素的时空一致性。视频生成模型以 BEV 布局、3D 边界框序列和文本描述为条件，生成新轨迹上时空一致的多帧视频；"亲缘数据训练"策略融合真实和合成数据优化 4DGS，使重建场景在训练轨迹外也保持高质量。

\paragraph{三平面扩散与一致性外推的大规模生成}

X-Scene（2025.06）~\cite{2506_13558} 统一生成 3D 语义占用、多视角图像和视频，支持多粒度控制（低级布局 + 高级语义引导），通过一致性感知外推将局部区域扩展为大规模场景，最终提升为高质量 3DGS 表示。与 DreamDrive 从驾驶视频出发不同，X-Scene 采用"3D 先生成、2D 后引导"的层次化生成策略——首先用三平面扩散模型生成 3D 语义占用（占用扩散变换器在紧凑三平面潜空间中合成几何），再将占用体素转换为 3D 高斯基元渲染为语义和深度图引导多视角图像生成，最后通过时间感知扩散扩展为视频。场景外推通过三平面扩展实现：在已知区域和待生成区域的公共掩码约束下同步去噪，保持重叠区域的几何一致性并向外合理扩展。

\paragraph{占用中心化大规模生成}

Li et al.（2025.10）~\cite{2510_22973} 构建 Nuplan-Occ（迄今最大语义占用数据集），开发统一框架联合合成高质量语义占用、多视角视频和 LiDAR 点云。与 X-Scene 的渐进外推策略不同，该框架采用时空解耦生成策略——空间扩展生成器学习自车移动时静态结构的空间关系，时间预测生成器专注建模周围车辆的运动模式——两者共享架构但分别微调。视频生成阶段引入高斯泼溅稀疏点图渲染策略：将占用体素转换为极小尺度 3D 高斯（默认 0.01），通过无迹变换（Unscented Transform）校正传感器标定偏差后泼溅为语义和深度图引导视频扩散。LiDAR 生成阶段设计多 LiDAR 传感器感知嵌入，通过傅里叶编码各 LiDAR 传感器原点生成独立嵌入后 AdaLN-Zero 注入稀疏 UNet，使网络区分不同 LiDAR 的扫描模式并支持灵活的传感器配置仿真。

\subsubsection{数据增强与预训练}

本路线介绍利用 3DGS 重建场景进行数据增强和下游模型预训练的方法，为感知与建图提供低成本高保真数据。

3DGS 重建的场景不仅可以渲染训练图像，还可以直接在 3D 空间中编辑、增删物体并生成对应标注，为下游感知与建图模型提供低成本、高保真的数据来源。本节梳理 12 条相关技术路线。

\paragraph{3DGS 作为 3D 检测数据生成器}

Zanjani et al.（2025.04）~\cite{2504_16740} 探索将 3DGS 用于 3D 目标检测的数据增强——直接在重建的 3D 空间放置物体，确保物理合理性和精确 3D 标注。该方法修改了 3DGS 的场景图操作环节：利用 OmniRe 的场景图将场景分解为静态背景、刚性车辆和非刚性行人节点，物体插入通过在场景图中添加新的刚性节点实现。与 DiffusionHarmonizer 等基于扩散模型的方法不同，该方法直接在 3D 高斯空间操作——从 BEV 道路地图中采样可行驶区域确定放置位置，通过可见性检查（比较代理高斯到相机的距离与已渲染深度图）和碰撞检查过滤不合理的放置，保证插入物体在多相机视角中均可见且不与现有物体重叠。实验表明增加物体放置的几何多样性（位置、朝向的变化）比外观多样性影响更大。

\paragraph{自监督体素速度的视觉中心预训练}

VisionPAD（2024.11）~\cite{2411_14716} 提出面向自动驾驶的自监督预训练范式，利用 3DGS 仅以图像为监督重建多视角表示。自监督体素速度估计通过扭曲体素到相邻帧学习运动线索——将当前帧的 3D 体素特征按预测的体素速度变换到相邻帧后投影渲染，与相邻帧真值比较提供自监督信号；多帧光度一致性通过跨视角重建损失增强几何感知，无需任何 LiDAR 或深度标注。

\paragraph{任意传感器配置的虚拟相机渲染增强}

MapGS（2025.01）~\cite{2501_06660} 利用高斯泼溅重建场景并渲染目标传感器配置的相机图像，解决在线建图模型跨传感器配置泛化问题。

\textbf{Baseline 与问题。} 在线建图（online mapping）模型通常在特定传感器配置（相机数量、位置、内外参）的数据上训练，部署到不同传感器配置的车辆时性能显著下降——例如用 6 相机数据训练的模型直接应用于 4 相机配置会出现大量伪影和漏检。为每种传感器配置单独采集和标注训练数据成本高昂且不现实。核心问题是如何在不实际部署目标传感器的情况下，为任意传感器配置生成训练数据。

\textbf{核心改动。} MapGS 的思路极为简洁：用 3DGS 重建场景后，在重建的 3D 空间中虚拟放置任意数量和位置的相机，通过泼溅渲染为这些虚拟相机生成训练图像和对应的 BEV 标注——本质上将 3DGS 用作"任意传感器配置的图像合成器"。

\textbf{实现原理。} 管线分三步：（1）场景重建——对每段驾驶数据用标准 3DGS 训练，获得场景的显式高斯表示；（2）传感器配置采样——定义目标传感器配置 $\mathcal{C}_{target} = \{K_i, E_i, P_i\}_{i=1}^{N_c}$（相机内参、外参到车体变换、图像分辨率），其中 $N_c$ 可与训练配置不同；（3）泼溅渲染——将重建的 3D 高斯按目标相机参数泼溅渲染为图像 $I_i^{target} = F_{splat}(G_{scene}, P_i)$，并同步生成对应的 BEV 标注（道路边界、车道线等从 3D 高斯直接投影到鸟瞰图）。

\textbf{设计动机。} 3DGS 的显式 3D 表示使得虚拟相机渲染在数学上等价于改变泼溅的投影参数——无需任何额外训练或微调即可适配新配置。相比 NeRF 需要为每个新视角重新运行体渲染，3DGS 的光栅化泼溅在 $O(1)$ 时间内完成单帧渲染，使得大规模数据生成（数千配置 $\times$ 数万帧）在计算上可行。

\textbf{性能。} 在 nuScenes 和 Argoverse 2 上通过数据增强提升 18\% 性能，仅用 25\% 原始训练数据即超越当时最优。跨传感器配置泛化实验表明，MapGS 生成的数据训练的模型在目标配置上的性能接近在真实目标配置数据上训练的模型。

\paragraph{3D 表面重建提升 2D 语义的伪标签建图}

PseudoMapTrainer（2025.08）~\cite{2508_18788} 利用高斯泼溅从多相机图像重建道路表面，结合预训练 2D 分割网络的语义生成伪标签，首次实现无需任何真值高精地图的在线建图模型训练。与 MapGS 用 3DGS 渲染新传感器配置图像的思路不同，PseudoMapTrainer 利用 3DGS 的核心能力在于"精确 3D 表面重建"——从多相机图像重建场景后，沿道路表面均匀采样 3D 点并投影到各相机视角，用预训练 2D 语义分割网络预测各投影点的语义标签，再将多视角语义投票聚合回 3D 空间生成伪高精地图标注。这种范式将 3DGS 作为几何桥梁，将 2D 语义信息提升到 3D 鸟瞰图空间，绕过了高精地图真值标注的成本瓶颈。

\paragraph{CAD 检索优化的可控车辆数字孪生}

UrbanCAD（2024.11）~\cite{2411_19292} 从单张城市场景图像生成高度可控的照片级 3D 车辆数字孪生。采用检索-优化管线：首先从 CAD 模型库中检索与观测车辆形状最相似的模型作为初始化，然后在保留 CAD 模型的专家设计先验（精确的几何结构和拓扑）的同时通过 3DGS 优化适应观测数据中的具体外观细节。支持 360 度渲染、材质迁移和重光照。

\paragraph{大规模记忆库检索的全视角车辆替换}

MADrive（2025.06）~\cite{2506_21520} 从大规模外部记忆库检索视觉相似 3D 资产替换场景中的车辆，发布 MAD-Cars 数据集（$\sim$70K 360 度汽车视频），通过方向对齐和重光照集成到目标场景。与 UrbanCAD 通过 CAD 检索+3DGS 优化的生成式路线不同，MADrive 采用检索式路线——用 SigLIP2 嵌入和颜色估计从记忆库中匹配最相似车辆，再用 2D 高斯泼溅重建可重光照的 3D 资产（基于 Cook-Torrance 着色模型解耦材质和环境光照）。关键设计在于解决驾驶场景中车辆仅有单侧观测的问题：替换为全视角 3D 资产后，即使车辆转弯或出现在新视角也能渲染完整外观。插入时通过 ICP 对齐 LiDAR 点云精确匹配位姿，用 DiffusionLight 从前相机估计环境贴图实现光照一致的阴影和反射，车底添加半透明黑色平面模拟阴影。

\paragraph{一致性蒸馏的一步扩散资产插入}

R3D2（2025.06）~\cite{2506_07826} 是轻量级一步扩散模型，实时生成逼真的渲染效果（阴影、一致光照）实现 3D 资产插入。与多步扩散模型不同，R3D2 通过一致性蒸馏将预训练的多步模型压缩为单步推理，在单次前向传播中同时生成被遮挡区域的补全、环境光照一致的阴影和与新位置匹配的反射。训练数据通过图像条件 3D 生成模型从野外驾驶数据创建 3DGS 资产并合成放置到神经渲染环境中。

\paragraph{统一资产插入与新视角合成的视频扩散}

SCPainter（2025.12）~\cite{2512_22706} 统一 3DGS 车辆资产表示和 3D 场景点云与扩散生成，联合实现逼真的 3D 资产插入和新视角合成。与 R3D2 用一步扩散模型修复插入资产周围区域不同，SCPainter 基于 Stable Video Diffusion 构建统一的图像修复框架——3D 车辆资产用 Amodal3R 从部分观测重建完整高斯模型，场景几何用 VGGT 恢复着色点云。两者在新视角下联合渲染为带掩码的条件图像（资产区域和空白区域分别标记），输入 SVD 模型同时完成资产的环境光照适应（阴影、反射）和稀疏投影区域的细节修复。训练时通过"资产替换配对"策略构造监督：将场景中原有车辆替换为生成的 3DGS 资产，模型学习从生成资产渲染恢复到真实车辆外观的映射。

\paragraph{自适应深度变形的跨视角渲染转换}

I2V-GS（2025.07）~\cite{2507_23683} 首次实现从基础设施视角到车辆视角的高斯泼溅转换，采用自适应深度变形生成密集训练视角，级联策略确保跨视图修复一致性，在 NTA-IoU、NTL-IoU、FID 上分别超越 StreetGaussian 45.7\%、34.2\%、14.9\%。与标准 3DGS 街景重建假设自车中心采集不同，I2V-GS 处理的是视角极端不对称的挑战——基础设施相机固定在高处俯瞰道路，车辆相机处于低位前向视角，两者之间几乎不存在视角重叠。自适应深度变形模块将基础设施视角的 3D 高斯按深度依赖的偏移量变形到车辆视角的几何空间，生成原本不存在的稠密训练视图来弥补视角间隙；级联策略先修复几何再修复外观，逐步提升跨视角渲染的视觉一致性和感知精度。

\paragraph{复合嵌入解耦几何与风格的街景编辑}

StyledStreets（2025.03）~\cite{2503_21104} 实现指令驱动的场景编辑，保证时空一致性——混合嵌入解耦持久场景几何与瞬态风格属性，不确定性感知渲染缓解扩散先验的监督噪声，统一参数模型防止几何漂移。与 DreamDrive 的混合高斯 4D 表示不同，StyledStreets 的核心改动在于为 Street Gaussians 的每个高斯基元引入复合嵌入 $E_i(T) = [e_i^{cam}; e^{style}(T)]$，将相机参数和风格指令（通过 CLIP 编码）分别编码后驱动独立的 MLP 变换网络——位置和协方差网络仅通过渲染损失更新保持几何稳定，外观网络联合渲染损失和风格损失实现视觉变化。这种设计使得单个统一模型可以通过不同文本指令（如"暴风雪夜晚""沙尘暴"）产生多样风格渲染，同时通过几何正则化 $L_{reg}$ 约束高斯参数不偏离初始值，确保编辑后场景的深度误差低于 2\%，维持下游规划任务的几何可靠性。

\paragraph{视觉基础模型驱动的跨域扩散泊车}

Dino-Diffusion Parking（2025.10）~\cite{2510_20335} 将视觉基础模型与扩散规划集成，在 CARLA 常规设置训练后零样本迁移到对抗设置，跨域泊车成功率保持 90\% 以上。与 ParkingWorld 在 3DGS 泊车仿真器中训练 RL 策略不同，Dino-Diffusion Parking 的核心设计在于跨域鲁棒性——使用冻结的 DINOv2 提取通用视觉特征后 Lift-Splat 到 BEV 空间，确保视觉等效场景产生一致特征表示（天气变化或周围车辆类型改变不影响 BEV 表征）；扩散模型在 SE(2) 笛卡尔空间规划轨迹（而非终端控制空间），配合 Stanley 控制器跟踪执行，降低复合误差。事后目标重标注策略通过随机扰动目标位姿并重新生成分割标签增强数据多样性，解决感知模块对训练域成功模式的过拟合问题。在真实世界 3DGS 重建环境中的测试展示了有前景的仿真到现实迁移。

\paragraph{其他仿真方法}

多个工作从工业应用和特殊角度推进仿真技术。Nouri et al.（2026.05）~\cite{2605_01995} 系统分析 3DGS 重建在安全相关场景中的保真度退化。Tan et al.（2026.05）~\cite{2605_21032} 提出正交投影梯度（OPG）解决闭环仿真中空间-时间参数的信用分配困境。3DArticCyclists（2024.10）~\cite{2410_10782} 使用 3DGS 生成可控关节自行车骑行者数据。RenderWorld（2024.09）~\cite{2409_11356} 用自监督高斯 Img2Occ 模块生成 3D 占用标签，实现仅视觉端到端驾驶。Gaussian Splatting as Data Generator（2025.04）~\cite{2504_16740} 证明 3DGS 重建场景可直接用作 3D 检测的高效数据增强源。Zeng et al.（2025.03）~\cite{2503_11731} 提出工业级传感器仿真模块化框架。NeRF2GS（2025.03）~\cite{2503_09464} 利用 NeRF 的泛化能力蒸馏到 3DGS 实现混合渲染。Sensor2Sensor（2026.05）~\cite{2605_22809} 通过 4DGS 重建实现行车记录仪到多模态传感器套件的生成转换。ReCamDriving（2025.12）~\cite{2512_03621} 提出纯视觉相机控制新轨迹视频生成框架，利用 3DGS 渲染提供密集几何引导实现精确相机可控生成，构建 ParaDrive 数据集包含超过 110K 并行轨迹视频对。

\begin{takeaway}
仿真与数据生成是 3DGS 在自动驾驶中的核心应用场景：闭环仿真系统将真实世界驾驶日志转化为可编辑的虚拟环境，视觉基础模型驱动的扩散规划实现跨域迁移，3DGS 正成为自动驾驶数据闭环的关键基础设施。
\end{takeaway}

%=====================================================================

\subsection{近期方法脉络：可编辑仿真与Sim-to-Real、车辆资产生成}

这一组近期工作延续 editable 3DGS driving simulation、feed-forward 3D Gaussians，主要把改动集中在增量改进、应用等环节。下面只保留相对既有路线的变化，便于和前文方法直接对照。

\paragraph{RealityBridge}
RealityBridge（2026.06）~\cite{2606_16278} 延续 editable 3DGS driving simulation 在“可编辑仿真与Sim-to-Real”上的基本设定，主要改变是：RealityBridge：针对3DGS渲染驾驶视频的恢复/生成框架，联合修复伪影、前景退化、光照不一致和闪烁。高斯基元仍承担可渲染动态场景表示。
其设计动机是：面向长尾危险驾驶场景，把3DGS可编辑仿真和视频恢复生成结合起来，处理3DGS特有渲染伪影和时序不稳定。
具体而言，方法从真实驾驶场景重建可编辑3DGS，进行场景编辑后用专门模型桥接到真实视频域。
适用边界是：恢复模型可能改变物理/几何真实性；复杂动态交通场景仍需严谨验证。

\paragraph{3DCarGen}
3DCarGen（2026.06）~\cite{2606_24257} 延续 feed-forward 3D Gaussians 在“车辆资产生成”上的基本设定，主要改变是：3DCarGen：先生成固定视角多视图，再合成任意数量3D一致视图并用于feed-forward 3D Gaussians/mesh生成。高斯基元仍承担可渲染动态场景表示。
其设计动机是：提出可扩展单图车辆生成框架，解决多视图扩散视角数量有限和跨视角几何不一致导致的重建保真下降。
具体而言，方法从单图合成固定视角图像，再扩展为任意数量3D一致多视图，并结合前馈3D Gaussian路线生成车辆3D资产。
适用边界是：单图输入下不可见区域仍依赖生成假设；车辆类别外泛化需验证。
\section{动态驾驶场景}
\label{sec:18_3}

动态驾驶场景重建是 3DGS 在自动驾驶中技术挑战最为集中的方向之一。城市街道中车辆、行人、自行车等动态参与者具有复杂的运动模式，而标准 3DGS 假设场景静态，直接应用于动态场景会产生严重伪影。核心问题在于如何准确分解静态背景和动态物体、如何建模动态物体的运动轨迹、如何在没有昂贵标注的情况下实现高质量重建。本节与第~\ref{sec:18_2}~节的仿真场景编辑共享部分技术基础（如场景图建模和动静分解），但本节聚焦于重建质量本身，而仿真场景更关注编辑可控性和下游任务适配。

\subsection{动静分解与运动建模}
\label{ss:18_4}

本节聚焦利用渲染残差和时空特征等线索实现无标注的动态物体分离。

\subsubsection{动静分解与自监督方法}

本路线介绍自监督和无标注的动静分解方法，利用重建残差作为动态性线索实现静态背景与动态物体的分离。

动态驾驶场景重建的第一步是将静态背景与动态物体分离。理想情况下，如果有了 3D 边界框标注，动静分解是直接的：框内的是动态物体，框外的是静态背景。但标注成本高昂且不可扩展，因此自监督和无标注的动静分解成为研究热点。核心思路是利用``3DGS 在动态区域重建失败''这一特性——将重建残差作为动态性线索，再结合正则化约束逐步分离。

\paragraph{HexPlane 时空分解的自监督动静分离}

$S^3$Gaussian（2024.05）~\cite{2405_20323} 首次提出无需 3D 标注的自监督街景高斯方法，从 4D 一致性分解动静元素。

\textbf{Baseline 与问题。} 此前的动态街景方法（如 Street Gaussians、DrivingGaussian）依赖 3D 边界框标注来分离动态物体和静态背景，获取这些标注成本高昂。如何在完全没有物体边界框、分割掩码和光流标注的情况下实现高质量的动静分解是一个开放问题。

\textbf{核心改动——时空场网络。} $S^3$Gaussian 维护一组规范（canonical）3D 高斯，配备时空场网络预测每帧的变形。网络核心是 Multi-resolution HexPlane 编码器，将 4D 时空网格分解为六个多分辨率可学习特征平面：三个空间平面 $P_{xy}, P_{xz}, P_{yz}$ 和三个时空平面 $P_{xt}, P_{yt}, P_{zt}$。这种分解天然有利于动静分离——动态物体在时空平面上呈现明显特征，静态物体仅在空间平面上可见。多解码头分别预测位移 $\Delta X$、SH 系数残差 $\Delta C$ 和语义特征 $f_s$，变形后的 4D 高斯为 $G' = \{X + \Delta X, C + \Delta C, s, r, o, f_s\}$。

\textbf{自监督优化。} 损失函数包含七项：RGB 重建损失 $L_{rgb}$、深度监督 $L_{depth}$、语义特征损失 $L_{feat}$、SSIM 损失 $L_{ssim}$、全变分正则化 $L_{tv}$，以及两项关键正则化——位移正则化 $L_{reg}^x$ 和颜色正则化 $L_{reg}^c$，它们最小化 $\mathbb{E}[\|\Delta X\|]$ 和 $\mathbb{E}[\|\Delta C\|]$，鼓励网络仅在必要时产生偏移，从而自动区分静态和动态高斯。使用 LiDAR 点云替代 SfM 初始化以提供更好的几何先验。

\textbf{设计动机。} HexPlane 分解的核心洞察是：静态物体在时间维度上不变，因此其特征完全由空间平面 $P_{xy}, P_{xz}, P_{yz}$ 编码，时空平面 $P_{xt}, P_{yt}, P_{zt}$ 的贡献趋近于零；动态物体则必须在时空平面上产生显著特征才能准确重建——这种天然的特征分离使得变形场自动区分动静区域。位移和颜色正则化进一步强化这一先验，最小化运动假设的复杂度，使"无需运动"成为默认选择。

\textbf{关键公式。} 对规范高斯 $\mathbf{g} = (X, C, s, r, o)$ 在时刻 $t$ 的变形由 HexPlane 特征解码：

\begin{equation*}
f(\mathbf{x}, t) = \sum_{r=1}^{R} \left( P_{xy}^r(\mathbf{x}) + P_{xz}^r(\mathbf{x}) + P_{yz}^r(\mathbf{x}) + P_{xt}^r(\mathbf{x}, t) + P_{yt}^r(\mathbf{x}, t) + P_{zt}^r(\mathbf{x}, t) \right)
\end{equation*}

其中 $R$ 为多分辨率层级数，$P^r$ 为第 $r$ 层的特征平面。总损失为 $\mathcal{L} = \mathcal{L}_{rgb} + \lambda_d \mathcal{L}_{depth} + \lambda_f \mathcal{L}_{feat} + \lambda_{ssim} \mathcal{L}_{ssim} + \lambda_{tv} \mathcal{L}_{tv} + \lambda_x \mathcal{L}_{reg}^x + \lambda_c \mathcal{L}_{reg}^c$。

\textbf{性能。} 在 Waymo-Open 上实现无标注情况下的最优动静分解，无需任何物体标注即可高质量重建包含车辆和行人的动态街景。

\paragraph{渲染残差驱动的两阶段自监督动静分解}

DeSiRe-GS（2024.11）~\cite{2411_11921} 修改了 3DGS 的动态场景处理环节——用两阶段自监督分解管线替代 $S^3$Gaussian 的统一变形场方案——静态重建能力本身被用作运动检测器。核心观察极为简洁：原始 3DGS 不包含时间信息，因此能准确重建静态区域（停放的车辆、建筑物），但在动态区域会产生伪影（漂浮的"鬼影"点），导致渲染图像与真值之间出现显著差异。

\textbf{第一阶段：2D 运动掩码提取。} 问题：如何不依赖任何标注来识别动态区域？方案：直接利用 3DGS 的重建失败作为动态性信号。具体而言，用预训练基础模型分别提取渲染图像 $\hat{I}$ 和真值图像 $I$ 的特征 $\hat{F}$ 和 $F$，计算逐像素不相似度：

\begin{equation}
D = \frac{1 - \cos(\hat{F}, F)}{2}
\end{equation}

不相似度 $D$ 在静态区域趋近 0（重建成功），在动态区域趋近 1（重建失败），且该计算不涉及任何可学习参数。但 DeSiRe-GS 并未对 $D$ 直接阈值化，而是引入一个 MLP 解码器预测动态概率 $\hat{\mathcal{D}} \in \mathbb{R}^{H \times W}$，以特征中的丰富语义信息为输入、以 $D$ 为监督信号：

\begin{equation}
L_{\text{dyn}} = \hat{\mathcal{D}} \odot D
\end{equation}

即鼓励解码器在 $D$ 高的区域输出高值。最终二值掩码 $M = \mathbb{I}(\hat{\mathcal{D}} > \tau)$（$\tau$ 为固定阈值）。关键设计：联合优化图像渲染与掩码预测——用掩码遮蔽动态区域后的渲染损失 $L_{\text{masked-render}} = M \odot \|\hat{I} - I\|$ 反过来让动态区域的差异更显著，形成互惠增强。第一阶段训练共 30,000 次迭代，运动解码器从第 6,000 次迭代开始训练。

\textbf{第二阶段：掩码蒸馏至高斯空间。} 问题：第一阶段的掩码仅存在于 2D 图像空间且依赖真值图像，无法用于新视角合成。方案：以可微分方式将 2D 运动信息蒸馏到 3D 高斯空间。DeSiRe-GS 采用 PVG（Periodic Vibration Gaussian）这一时间变化高斯表示——每个高斯附加速度 $v$，渲染时生成速度图 $V \in \mathbb{R}^{H \times W}$。但 PVG 的原始速度图包含大量噪声（道路和建筑物等静态区域的速度应为零却不为零），动静分解不理想。DeSiRe-GS 的核心改进是用第一阶段获得的掩码 $M$ 正则化速度图：

\begin{equation}
L_v = \|V \odot M\|
\end{equation}

该损失惩罚静态区域（掩码 $M$ 中为 0 的像素）中非零速度，有效消除 PVG 的噪声异常值。最小化 $L_v$ 后，每个高斯拥有精确速度 $v$，通过简单阈值即可区分动静高斯。这一过程将运动信息从 2D 局部帧传播到全局高斯空间，实现纯自监督分解。

\textbf{几何正则化。} 驾驶场景中视角稀疏（通常仅 2--3 个前向相机），容易导致高斯过拟合训练视角。DeSiRe-GS 引入三项正则化：(1) \emph{扁平化}：受 2DGS 启发，最小化最小尺度分量 $L_s = \|(s_1, s_2, s_3)\|_1$ 将 3D 椭球压扁为 2D 盘面，更好地贴合物体表面；(2) \emph{法向--尺度耦合}：不额外添加法向量参数，而是直接从尺度向量推导法向 $n = R \cdot \text{argmin}(s_1, s_2, s_3)$（$R$ 为旋转矩阵），梯度可回传至尺度向量；(3) \emph{巨型高斯惩罚}：对尺度超过阈值 $\theta$ 的高斯施加惩罚 $L_g = s_g \cdot \mathbb{I}(s_g > \theta)$（$s_g$ 为最大尺度分量），防止无界场景中出现覆盖大片区域的低透明度巨型椭球。

\textbf{时间跨视图一致性。} 假设静态区域深度在时间上保持一致。对参考帧中的静态像素 $(u_r, v_r)$，利用相机内外参投影到最近邻视角再投影回来，计算重投影误差：

\begin{equation}
L_{uv} = \|(u_r, v_r) - (u_{nr}, v_{nr})\|_2
\end{equation}

该损失聚合不同时间视角的信息，显著缓解稀疏视角下的几何过拟合。

\textbf{与 $S^3$Gaussian 的对比。} $S^3$Gaussian 将所有高斯视为潜在的动态元素，通过 HexPlane 变形场统一建模，依赖正则化间接学习运动模式——即"先假设全部可能运动，再通过约束抑制不必要的运动"。DeSiRe-GS 则采用截然不同的策略："先检测哪里在运动，再将运动信息注入高斯"。两阶段管线使运动信息的学习从间接的损失正则化变为直接的掩码蒸馏，在 Waymo 和 KITTI 上实现了超越 PVG 和 $S^3$Gaussian 的自监督分解质量，甚至在部分指标上媲美使用 3D 边界框标注的方法（HUGS、StreetGS）。渲染速度约 40 FPS，PSNR 提升 1.1 dB 以上。

\paragraph{互惠训练的实例级动态场景重建}

DIAL-GS（2025.11）~\cite{2511_06632} 采用两阶段自监督管线实现实例级动态场景重建。第一阶段建立"瞬时高斯场"：先用所有帧重建一个过滤波的静态场景 $S_{over}$，再为每帧生成动态候选高斯 $C_t$，将二者组合后变形到后续帧进行新视角渲染，通过比较渲染结果与真值之间的外观不一致（像素差异）和位置不一致（边界框中心、面积和边缘差异）计算每个实例的动态分数，经三次方放大后阈值化筛选出动态实例集合。第二阶段采用实例感知 4D 高斯（基于 PVG 扩展，每个高斯附加 8 位 ID 嵌入向量），提出互惠训练机制：身份掩码约束动态变形区域（限制速度和寿命在掩码外为零），而同一实例内高斯的速度方向、幅度和寿命一致性损失反过来精化身份分割——形成身份与动态相互增强的闭环。最终无需任何 3D 标注即可实现实例级场景编辑（移除、移动、修改动态物体外观）。

\paragraph{4D 超点时空聚类的无监督实例分解}

UnIRe（2025.04）~\cite{2504_00763} 提出 4D 超点（4D superpoints）表示——在 4D 时空空间中聚类多帧 LiDAR 点，基于时空关联实现无监督实例分离。具体而言，先对各帧 LiDAR 点独立运行 DBSCAN 获取逐帧聚类，再用自监督场景流建立帧间簇对应关系（匹配场景流对齐后落入同一簇的点数最多的簇对），处理消失、出现和分裂三种变换后得到在时间上保持一致身份的 4D 超点。为解决分裂导致的过分解问题，UnIRe 进一步利用时空相似度矩阵（融合运动方向余弦相似度和空间邻近度）对 4D 超点进行二次聚类合并。最终 4D 超点为规范空间初始化和逐点变形场提供精确的时空锚点，训练动态 3DGS 无需边界框或物体标注。

\paragraph{零样本掩码与刚体变换的实例分解}

IDSplat（2025.11）~\cite{2511_19235} 将动态物体建模为经历 SE(3)（特殊欧氏群，描述三维空间中的刚体变换）刚体变换的连贯实例，每个实例的所有高斯共享同一变换 $T_{z,t} \in \text{SE}(3)$，从规范坐标系变换到世界坐标系。实例分解方面，利用 Grounded-SAM-2（基于语言提示的通用视频分割模型）生成零样本语言接地的视频实例掩码，将 2D 掩码投影到 LiDAR 点云后经腐蚀和 DBSCAN 过滤锚定到 3D。轨迹估计方面，以实例最密集帧的包围盒中心定义规范坐标系，利用 DINOv3（DINO 系列第三代自监督视觉Transformer）特征建立帧间对应关系并通过 RANSAC+Umeyama 估计刚体变换；再通过协调转弯（Coordinated-Turn）平滑优化构建物理一致的运动轨迹，迭代剔除离群测量值提升鲁棒性。该方法在 Waymo 上无需重训练即泛化到多样序列。

\paragraph{对比学习与码本量化的 3D 实例分割}

InstDrive（2025.08）~\cite{2508_12015} 是首个在动态开放世界驾驶场景中实现 3D 实例分割的框架，核心思路是分阶段学习实例特征再量化为离散身份标记。2D 层面使用 SAM 掩码引导对比学习——同一掩码内像素特征拉近（类内一致性）、不同掩码原型推远（类间可分性）；3D 层面引入体素一致性损失（随机偏移体素网格内高斯特征对齐到均值），补偿纯 2D 监督导致的 3D 特征不连贯。轻量级静态码本（codebook）基于 $d$ 维二值向量的笛卡尔积构建 $2^d$ 个均匀分布的码字，无需额外训练参数即可将连续特征量化为离散身份标记；配合基于多数投票的伪监督机制（掩码内最频繁码字作为统一标签）和排斥正则化进一步锐化实例边界，实现推理时无需 SAM 掩码的 3D 实例分割。

\subsection{综合动态场景重建}
\label{ss:18_5}

本节介绍从参数化轨迹表示到统一框架的综合动态街景重建方法。

\subsubsection{运动建模与轨迹表示}

本路线介绍为动态物体建立参数化运动模型的方法，利用贝塞尔曲线、B 样条等轨迹表示约束运动的物理合理性。

分离出动态物体后，需要为其建立准确的运动模型。简单的全局变形场（如 HexPlane 的 4D 特征）对长序列运动建模不够精确，且难以保证物理合理性。参数化轨迹表示（贝塞尔曲线、B 样条等）将运动约束在平滑路径上，天然适合车辆和行人的刚体运动模式。

\paragraph{双尺度可学习运动嵌入的即插即用分解}

EMD（2024.11）~\cite{2411_15582} 通过引入可学习运动嵌入到高斯基元中增强动静分解，作为即插即用模块补偿自监督方法中缺失的运动建模。与 $S^3$Gaussian 用全局 HexPlane 统一建模所有运动不同，EMD 的核心改动是为每个高斯分配可学习的个体嵌入 $z_k \in \mathbb{R}^{32}$，并通过双尺度变形框架分别建模大尺度运动（车辆平移，$D_{coarse}$）和精细变形（行人体态变化，$D_{fine}$）。设计动机是街景中车辆和行人的运动模式差异显著——车辆以刚体平移为主、速度较快，行人则是非刚体运动、速度较慢——全局变形场无法有效捕捉这种异质性。具体差异在于：自适应时间嵌入 $F_{temp}(t)$ 从 $N_{min}=30$ 渐进增加到 $N_{max}=150$ 个时间采样，训练初期学习粗粒度运动模式、后期精细化时间动态；EMD 可无缝集成到 $S^3$Gaussian、DeSiRe-GS 等自监督方法和 Street Gaussians、OmniRe 等有监督方法中，在 Waymo 上稳定提升动态区域 PSNR。

\paragraph{可学习贝塞尔曲线的参数化运动建模}

B\'ezierGS（2025.06）~\cite{2506_22099} 修改了 3DGS 的动态物体运动建模环节——用可学习贝塞尔曲线轨迹替代全局变形场（如 HexPlane）的隐式运动表示——静态背景的标准 3DGS 重建保持不变。

\textbf{Baseline 做法与问题。} 此前动态街景方法对运动轨迹的建模可分为两类。第一类是边界框方法（Street Gaussians、DrivingGaussian），它们依赖手动标注的 3D 边界框位置和朝向来放置动态物体，但标注本身存在误差和遗漏（尤其在 nuPlan 等标注质量较差的数据集上），直接影响重建质量。第二类是自监督方法（$S^3$Gaussian），用 HexPlane 等时空分解网络隐式建模运动轨迹——将 4D 时空网格分解为多分辨率特征平面，通过网络预测每帧变形。这种隐式表示有两个困难：(1) 全局变形场的自由度极高，优化困难，难以保证轨迹的物理合理性；(2) 无法显式约束同一物体上不同高斯基元的运动一致性，容易出现某些高斯脱离物体整体运动的"漂浮"现象。PVG 虽尝试用周期振动拼接长轨迹，但振动模式和实际运动不匹配，且分段拼接无法利用同一物体跨时间的完整一致性。

\textbf{核心改动——可学习贝塞尔曲线轨迹。} B\'ezierGS 的核心思路极其直观：与其让神经网络隐式学习运动，不如直接用一条参数化曲线拟合物体运动轨迹。就像用一根弹性绳穿过若干锚点来勾勒平滑路径一样，贝塞尔曲线用少量控制点就能精确描述车辆从起点到终点的完整运动，曲线的平滑性天然保证运动没有突变和抖动。

具体而言，B\'ezierGS 将场景分为静态背景和动态前景两部分。静态背景直接用标准 3DGS 重建（高斯属性 $\{\mu, q, s, o, c\}$），因为背景在所有帧中保持一致。对于动态前景，每个动态物体用一组独立的高斯基元表示，并为每个高斯引入分组属性 $g \in \mathbb{N}_+$ 标识其所属物体。

\textbf{物体中心轨迹。} 每个物体 $g$ 的中心运动轨迹由一条 $n$ 阶贝塞尔曲线定义，由 $n+1$ 个可学习控制点 $\{p_i^g\}_{i=0}^n$ 控制。直觉上，控制点就像是路线图上的关键途经点，曲线自动在这些途经点之间生成平滑插值路径。给定曲线参数 $t \in [0,1]$（$t=0$ 对应第一个控制点，$t=1$ 对应最后一个），物体中心位置为：

\begin{equation}
\mathbf{c}(t, g) = \sum_{i=0}^{n} b_{i,n}(t) \, p_i^g, \quad t \in [0,1]
\end{equation}

其中 $b_{i,n}(t) = \binom{n}{i} t^i (1-t)^{n-i}$ 是 Bernstein 基多项式，决定了每个控制点在参数 $t$ 处的影响力。这可以理解为对控制点的加权平均，权重随 $t$ 平滑变化。

\textbf{高斯偏移轨迹。} 物体上每个高斯基元相对于物体中心的偏移也用贝塞尔曲线建模（控制点为 $\{o_i\}_{i=0}^n$），偏移在世界坐标系中表示：$\boldsymbol{\delta}(t) = \sum_{i=0}^{n} b_{i,n}(t) \, o_i$。这样，高斯基元在时刻 $\tau$ 的完整运动轨迹为：

\begin{equation}
\mathbf{pos}(\tau, g) = \boldsymbol{\delta}(t) + \mathbf{c}(t, g), \quad t = f(\tau, g)
\end{equation}

即"物体整体运动 $+$ 局部偏移"。关键在于 $f(\tau, g)$ 是时间-参数映射：物体沿曲线的运动速度并不均匀，同一物体在不同时间段可能有加速或减速，因此需要将物理时间戳 $\tau$ 映射到曲线参数 $t$。B\'ezierGS 同样用一条额外的贝塞尔曲线来学习这个映射，而不是假设匀速运动。与边界框方法将偏移定义在物体局部坐标系且保持恒定不同，B\'ezierGS 的偏移定义在世界坐标系且随时间变化，因此既能建模刚体（车辆），也能建模非刚体（行人上半身）。标准实现中采用三次贝塞尔曲线（$n=3$，即 4 个控制点），这是轨迹建模中公认有效的配置。

\textbf{设计动机——与全局变形场的对比。} B\'ezierGS 选择参数化曲线而非全局变形场，有三个核心原因。第一，\textbf{显式轨迹表示}：HexPlane 等全局变形场将运动隐式编码在高维特征中，缺乏直接的几何约束；贝塞尔曲线则显式定义了轨迹的数学形式，控制点与物理位置直接对应，优化目标明确。第二，\textbf{自动姿态校正}：边界框方法完全依赖标注精度，标注偏移会直接导致重建失败；贝塞尔曲线的控制点与高斯参数联合优化，在训练过程中自动修正标注误差，使 B\'ezierGS 在标注质量较差的 nuPlan 上也能实现高质量重建（PSNR 提升 3.04 dB）。第三，\textbf{组级一致性约束}：同一物体上的高斯基元共享物体中心轨迹 $\mathbf{c}(t,g)$，只通过各自的偏移轨迹 $\boldsymbol{\delta}(t)$ 区分，天然保证了组内运动一致性，而全局变形场需要额外的复杂正则化才能实现类似效果。

\textbf{曲线间一致性损失。} 尽管分组机制提供了基本约束，但单个高斯基元仍可能在优化中脱离物体整体。问题在于：高斯基元自由度很高，某个基元可能在某些时间步由不同基元代表物体的不同部分，导致新视角渲染不一致。方案是引入曲线间一致性损失，利用刚体运动的几何特性——刚体（如车辆）各部分的轨迹偏离量大小保持恒定。因为高斯基元在 $t=0$ 和 $t=1$ 时恰好与首尾控制点重合，所以可以约束任意时刻 $t$ 的偏移量范数不超过首尾控制点偏移量范数的平均值：

\begin{equation}
L_{\text{icc}} = \left\| \|\boldsymbol{\delta}(t)\| - \frac{\|o_0\| + \|o_n\|}{2} \right\|_1
\end{equation}

该损失有效抑制局部几何的过度变化，增强物体结构在时间上的整体一致性和稳定性。

\textbf{速度损失。} 为了从多维度约束动态分量重建，B\'ezierGS 还对动态高斯渲染的速度图 $V_{\alpha}^{\text{dyn}}$ 施加约束。速度由轨迹曲线对参数 $t$ 的导数直接计算：$\mathbf{v}(\tau,g) = \frac{d\mathbf{c}}{dt} + \frac{d\boldsymbol{\delta}}{dt}$，并通过 $\alpha$-混合 渲染为速度图。速度损失 $L_v = \|V_{\alpha}^{\text{dyn}} \odot (1 - M^{\text{dyn}})\|_2$ 惩罚动态物体区域之外的非零速度，隐式阻止动态高斯漂移到静态区域。完整损失函数为 $L = (1-\lambda_r)L_1 + \lambda_r L_{\text{ssim}} + \lambda_o^{\text{sky}} L_o^{\text{sky}} + \lambda_{\text{icc}} L_{\text{icc}} + \lambda_{\text{dr}} L_{\text{dr}} + \lambda_v L_v + \lambda_d L_d$，其中 $\lambda_r=0.2$，$\lambda_d=1.0$，$\lambda_{\text{icc}}=0.01$，$\lambda_v=1.0$。

\textbf{性能。} 在 Waymo 上，新视角合成 PSNR 提升 1.87 dB，动态区域 Dyn-PSNR 提升 2.66 dB；在标注质量较差的 nuPlan 上，PSNR 提升 3.04 dB，LPIPS 降低 16.35\%，显著验证了贝塞尔曲线自动校正标注误差的能力。消融实验表明：曲线间一致性损失消除漂浮物并提升动态建模能力，动态渲染损失促进前景-背景分离，速度损失进一步约束动态高斯漂移。

\paragraph{B 样条与三角函数混合的运动建模}

AD-GS（2025.07）~\cite{2507_12137} 与 B\'ezierGS 用贝塞尔曲线建模运动轨迹不同，AD-GS 提出可学习运动模型集成局部感知 B 样条曲线和全局感知三角函数——B 样条捕获物体的短距离平滑运动（如变道），三角函数建模长距离周期运动（如匀速行驶）。

\textbf{与 B\'ezierGS 的差异。} B\'ezierGS 的贝塞尔曲线具有全局支撑性——移动一个控制点会影响整条曲线形状，适合建模完整的单次运动轨迹；AD-GS 选用 B 样条曲线正是因为其局部支撑性——每个控制点仅影响相邻 $k$ 个区间的曲线（$k$ 为 B 样条阶数），局部运动细节可通过仅调整相关控制点来精化，不会扰动已收敛的远端轨迹。运动建模表示为 $\mu' = \mu + p(t) + \sum_{l=1}^{L} a_l \sin(lt\omega) + b_l \cos(lt\omega)$，其中 B 样条项 $p(t)$ 处理局部运动变化，三角函数项 $\sum a_l \sin + b_l \cos$ 捕获全局周期模式。旋转使用 B 样条四元数曲线建模（基于 De Casteljau 算法的球面插值），确保旋转的平滑性和单位四元数约束。

\textbf{简化的伪 2D 分割。} 不依赖复杂的实例分割，而是将 SAM 分割结果简化为二类（物体/背景）——将所有可能运动的类别（车辆、行人等）统一标记为物体类，其余标记为背景。这种简化在自动驾驶场景中足够鲁棒，避免了实例分割中的噪声和漏检。双向时间可见性掩码 $\alpha'(t) = e^{-(t-\hat{t})^2 / (2s^2)}$（$\hat{t}$ 为 LiDAR 点采集时间，非可学习参数）控制物体高斯在对应时间窗口内的可见性，不同尺度的前后向衰减 $s_0, s_1$ 处理物体出现和消失的非对称性。配合光流伪监督（CoTracker3 生成）和逆深度渲染监督，无需综合语义标注即达到有标注方法的竞争性能。

\paragraph{语义-变形-时间三重感知的变形补偿}

CoDa-4DGS（2025.03）~\cite{2503_06744} 利用 2D 语义分割基础模型自监督高斯的 4D 语义特征，追踪每个高斯跨相邻帧的时间变形，聚合编码语义和时间变形特征使高斯具备变形补偿线索。与 $S^3$Gaussian 仅依赖 HexPlane 编码时空特征不同，CoDa-4DGS 在变形场中引入三重感知信号：语义感知（通过 LSeg 提取的 2D 语义特征经可微分泼溅为每个高斯赋予高维语义特征，并在变形过程中与位置、时间同步变形）、时间变形感知（将 HexPlane 已计算的变形量 $\Delta G$ 作为显式输入，而非仅依赖时间编码隐式推断）和时间嵌入（正弦位置编码提供帧级唯一标识）。变形补偿网络（DCN）通过双通道设计精化 HexPlane 的输出——MLP 通道预测变形补偿量，Sigmoid 通道学习空间掩码过滤天空和地面等无关高斯的噪声干扰，有效提升了驾驶场景中有限视角下的动态重建质量。

\paragraph{语义类别先验驱动的动静分离}

UrbanGS（2024.12）~\cite{2412_03473} 利用 2D 语义图区分确定静态和潜在动态元素，对静态区域强制全局一致性并引入 KNN 正则化改善低纹理地面，对潜在动态物体用可学习时间嵌入聚合时间信息建模变形。与 DeSiRe-GS 用渲染残差检测动态区域不同，UrbanGS 采用语义驱动的先验分解——用预训练分割模型将高斯分为确定静态类（建筑、道路、树木）和潜在动态类（车辆、行人、自行车），这种基于语义的分解无需训练额外的运动检测器，直接利用类别先验。静态一致性损失通过语义加权约束静态高斯的位置和不透明度在时间变换后不偏离初始值（权重 $w_i$ 从 softmax 比例导出，完全静态的点 $w_i \to 1$ 保持不变），有效消除了 PVG 等方法中时间变换对静态区域的干扰。低纹理地面的 KNN 正则化约束邻近地面高斯的尺度参数一致，间接保持法向方向的一致性，显著改善了驾驶场景中大面积弱纹理路面的重建质量。

\subsubsection{综合动态场景重建}

本路线介绍将动静分解、运动建模和多视角一致性整合为统一框架的综合动态场景重建方法。

综合动态场景重建将动静分解、运动建模和多视角一致性整合为统一框架，目标是实现包含所有类型动态参与者的完整街景重建。这一方向的关键挑战在于可扩展性——从仅处理车辆到同时建模行人、自行车等非刚性动态物体，从单一数据集到跨数据集泛化。

\paragraph{体积-像素-远景三分支前馈 4D 合成}

EVolSplat4D（2026.01）~\cite{2601_15951} 提出前馈框架统一体积和像素高斯预测——近距离静态区域从 3D 特征体预测一致的 3D 高斯（通过体素化特征网格直接查询），动态角色利用对象中心规范空间和运动调整渲染模块聚合时间特征（在规范空间中建模动态物体的变形，再通过位姿变换到世界坐标），远景由高效逐像素高斯分支覆盖（避免对远处弱纹理区域构建完整 3D 表示的开销）。在 KITTI-360、KITTI、Waymo、PandaSet 上超越所有基线。

\paragraph{体素可见性推理与多视角立体恢复}

VAD-GS（2025.10）~\cite{2510_09364} 针对无界动态城市环境中初始点云覆盖不完整的问题，通过体素化可见性推理识别不可靠几何，多样性感知视角选择确定信息性支撑视角，基于块匹配的多视角立体重建恢复缺失结构。与 PAGS 通过语义优先级简化已有高斯不同，VAD-GS 关注的是初始点云根本不存在时的几何恢复——驾驶场景中，远距离区域、动态物体遮挡区域和视角稀疏区域在 LiDAR 点云中往往存在大面积空洞。体素化可见性推理将场景划分为均匀体素格，统计每个体素被多少训练视角观测到来标记可靠程度；多样性感知视角选择在候选视角中挑选信息增益最大的支撑视角进行多视角立体匹配，通过块匹配恢复稠密深度并新增高斯基元填补空洞，从而在不增加额外传感器数据的情况下显著提升无界场景的重建完整性。

\paragraph{语义优先级驱动的自适应高斯简化}

PAGS（2025.10）~\cite{2510_12282} 将任务感知语义优先级注入 3D 重建管线——语义引导剪枝和正则化策略激进简化非关键场景元素同时保留安全关键物体细节，优先级驱动渲染管线通过深度预通（depth pre-pass）剔除被遮挡基元。语义优先级从 2D 分割结果传播到 3D 高斯，高优先级物体（如车辆、行人）保留更多高斯基元和更高的最大阶数 SH，低优先级区域（如路面、远距建筑）被大幅简化，渲染速度超过 350 FPS。

\paragraph{可学习 LOD 驱动的多尺度高斯分配}

DriveSplat（2025.08）~\cite{2508_15376} 提出场景感知可学习 LOD 建模策略显式处理近/中/远深度范围的自适应多尺度高斯分配（近距离使用密集小高斯捕获细节，远距离使用稀疏大高斯覆盖广阔区域）。与 OmniRe 为每种动态参与者设计独立高斯簇不同，DriveSplat 的核心改动是统一的 LOD 驱动分配机制——可学习的细节层次参数根据深度范围自动调节高斯密度和尺度，无需针对不同物体类型手动设计表示策略。动态角色通过全局刚体变换和两阶段变形处理非刚性动态——第一阶段调整锚点位置捕获物体整体运动，第二阶段更新每个高斯的局部偏移建模非刚性变形（如行人姿态变化）。这种分层变形策略兼顾了车辆等刚体物体的运动约束和行人的非刚性特征，在保持渲染质量的同时大幅降低远距离区域的高斯数量。

\paragraph{分类动态表示的全类型参与者在建}

OmniRe（2024.08）~\cite{2408_16760} 扩展动态建模从车辆到行人、自行车等所有动态前景，在 3DGS 上构建场景图并在规范空间构建多种高斯表示。

\textbf{Baseline 与问题。} 此前动态街景重建方法（如 Street Gaussians）仅处理车辆这一类动态物体，采用刚体运动模型。但城市场景中行人、自行车等非刚性动态参与者的运动模式与车辆截然不同——行人有体态变化（摆臂、迈步），自行车有倾斜转向——刚体模型无法捕捉这些非刚性变形。如何在统一框架中同时建模刚性和非刚性动态物体，是实现完整街景重建的关键挑战。

\textbf{改了什么。} OmniRe 将 Street Gaussians 的"单一刚体动态模型"扩展为"分类动态表示"——根据动态参与者的运动特性将其分为三类：静态背景（不变的高斯基元）、刚性物体（车辆，用刚体 SE(3) 变换建模）和非刚性物体（行人、自行车，在规范空间额外配备变形场）。

\textbf{实现原理。} 场景图为每种动态参与者分配独立的高斯簇和可优化的 SE(3) 位姿序列。刚性物体的建模沿用 Street Gaussians 的局部坐标系方法：车辆高斯定义在物体坐标系中，通过可优化位姿 $(R_t, T_t)$ 变换到世界坐标。非刚性物体的建模引入规范空间表示——行人等可变形物体的所有高斯定义在一个共享的规范坐标系中，通过两个变形阶段映射到世界坐标：

\begin{equation*}
\boldsymbol{\mu}_w = R_t \cdot \text{Deform}(\boldsymbol{\mu}_{cano}) + T_t
\end{equation*}

其中 $\text{Deform}(\cdot)$ 是基于 MLP 的非刚性变形场，以规范空间位置和时间 $t$ 为输入预测逐高斯偏移量。外观建模方面，所有动态物体使用 4D 傅里叶球谐编码时间变化外观（与 Street Gaussians 相同的策略）。背景使用 LiDAR 点云初始化的标准 3D 高斯。

\textbf{设计动机。} 分类动态表示的核心洞察是不同参与者的运动自由度差异显著：车辆为 6-DOF 刚体运动（仅平移+旋转），行人则涉及数十个关节的非刚性变形。对车辆使用刚体模型可避免过拟合（6 个参数 vs MLP 变形场的数百参数），对行人使用变形场则可捕捉体态变化。这种"按需复杂度"设计在保证渲染质量的同时控制了参数量。场景图结构使编辑操作（添加/删除/移动特定物体）天然可行——每个参与者是独立的图节点。

\textbf{性能。} 以 $\sim$60Hz 实现包含行人行为仿真和人车交互的综合仿真。在 Waymo 上首次实现同时重建车辆、行人和自行车的完整动态街景，非刚性物体区域的 PSNR 相比仅使用刚体模型的基线显著提升。

\paragraph{连续运动向量场的无追踪 4D 重建}

SplatFlow（2024.11）~\cite{2411_15482} 设计神经运动流场（Neural Motion Flow Field, NMFF）学习无追踪 3D 边界框的 4D 时空表示。与 B\'ezierGS 用参数化曲线显式建模轨迹不同，SplatFlow 的 NMFF 用连续运动向量场隐式建模运动——每个 MLP 场预测任意 3D 点在相邻帧间的 3D 运动流和运动角度，先在时序 LiDAR 点云上通过双向 Chamfer 距离预训练获得几何一致的运动先验，再与 4D 高斯联合优化。关键设计在于动静分解的几何驱动方式：对每个 LiDAR 点计算 NMFF 预测的运动流大小，通过阈值化自然分离静态和动态点，分别合并后初始化静态 3D 高斯和动态 4D 高斯。动态 4D 高斯在中间时间戳的参考帧中聚合所有观测增加密度和完整性，推理时通过 NMFF 的逐步传播将聚合高斯变形到目标时刻。同时将 SEA-RAFT 光流知识蒸馏到 4D 表示中——利用 NMFF 的 3D 运动流渲染 2D 光流作为监督信号，实现无需人工标注的端到端自监督动态场景重建。

\paragraph{运动感知时间注意力的前馈 4D 重建}

DynamicVGGT（2026.03）~\cite{2603_08254} 将 VGGT 从静态 3D 扩展到动态 4D 重建，联合预测当前帧和未来帧的点图（pointmap），引入运动感知时间注意力模块学习运动连续性。与 DrivingScene 的两阶段静态-动态分离策略不同，DynamicVGGT 的核心改动是在 VGGT 的交替注意力块之后插入运动感知时间注意力（MTA）模块——通过可学习运动标记在时间维度上跨帧传递信息，使网络在特征层面建立帧间运动对应关系。未来点预测头（FPH）从当前帧的增强特征直接预测未来帧的 3D 点图，通过时间一致性正则化约束预测位移与真值位移对齐，在点图层面隐式学习运动。动态 3DGS 头通过场景流监督预测每个高斯的 3D 速度向量——将 $t$ 时刻的高斯按预测速度平移到 $t+1$ 时刻，渲染后与 $t+1$ 真值比较提供显式的运动监督信号，与 FPH 的隐式点图位移监督形成互补。

\paragraph{静态先验冻结的残差流在线 4D 重建}

DrivingScene（2025.10）~\cite{2510_24734} 仅从两帧连续环绕视角图像在线重建 4D 动态场景，轻量级残差流网络预测动态物体的非刚性运动，从粗到细训练范式避免端到端方法的不稳定性。与 EVolSplat4D 的前馈 4D 重建不同，DrivingScene 采用显式的两阶段解耦策略：第一阶段专注训练静态场景先验网络，从大规模数据学习稳健的静态几何基础（建筑、道路等）；第二阶段冻结静态主干，仅训练轻量级残差流网络预测动态物体的非刚性运动——共享骨干学习通用运动先验，各相机头适配不同外参和内参保持一致的尺度预测。这种从静态到动态的渐进学习策略避免了 Driv3R 等端到端方法中动静耦合导致的训练不稳定，仅需两帧输入即可实现实时 4D 重建，全部使用自监督目标训练无需任何标注。

\paragraph{其他动态场景方法}

DENSER（2024.09）~\cite{2409_10041} 引入小波动态估计 SH 基增强远距离动态物体外观，跨多帧致密化点云加速收敛。VDG（2024.06）~\cite{2406_18198} 首次集成自监督视觉里程计实现仅 RGB 输入的无位姿动态高斯。EGSRAL（2024.12）~\cite{2412_15550} 引入自适应器实现自动标注，PSNR 在 nuScenes 达 29.04。ArmGS（2025.07）~\cite{2507_03886} 通过多粒度外观精化（局部高斯级、全局图像级、动态角色级）提升帧间和相机间精细变化建模。2D 先验与 SDF 引导（2025.10）~\cite{2510_13381} 结合 SDF 与 3DGS 实现无 LiDAR 城市场景渲染和场景编辑。SRIF（2024.09）~\cite{2409_11682} 利用扩散图像变形和流估计实现语义形状注册。

\begin{takeaway}
动态驾驶场景重建的核心是静态背景与动态物体的准确分解：静态先验冻结策略避免动静耦合的训练不稳定性，残差流网络以轻量级架构实现动态运动的实时预测，仅需两帧环绕视角即可在线 4D 重建。
\end{takeaway}

%=====================================================================

\subsection{近期方法脉络：未来场景预测、动态驾驶场景替代表示}

这一组近期工作延续 4D Gaussian Splatting、dynamic 3DGS，主要把改动集中在扩展、增量改进等环节。下面只保留相对既有路线的变化，便于和前文方法直接对照。

\paragraph{Envision4D}
Envision4D（2026.06）~\cite{2606_10656} 延续 4D Gaussian Splatting 在“未来场景预测”上的基本设定，主要改变是：Envision4D：自监督feed-forward 4D Gaussian Splatting，包含未来位姿预测、层内时间注意力和条件运动提升。高斯基元仍承担可渲染动态场景表示。
其设计动机是：把前馈3D/4DGS从插值扩展到未来外推，通过未来相机位姿去噪预测和非线性动态建模减少大位移ghosting。
具体而言，Future Pose Prediction模块迭代推断未来相机参数；In-layer Temporal Attention捕获时间关系；Conditioned Motion Lifting把不确定外推转化为关系映射并输出未来4D高斯场。
适用边界是：未来预测受运动不确定性和长时外推影响；复杂交通交互可能产生幻觉。

\paragraph{DrivingVoxels}
DrivingVoxels（2026.06）~\cite{2606_23031} 延续 dynamic 3DGS 在“动态驾驶场景替代表示”上的基本设定，主要改变是：DrivingVoxels：组合式稀疏体素光栅化框架，联合渲染多对象动态驾驶场景。高斯基元仍承担可渲染动态场景表示。
其设计动机是：提出比动态3DGS可能更快、更可控的组合式稀疏体素表示，服务无界驾驶环境和多动态对象。
具体而言，方法将静态背景和多个动态对象组织为稀疏体素并联合rasterize，控制内存增长。
适用边界是：不是3DGS方法；与3DGS的关系主要是替代和对比。
\section{传感器融合与定位}
\label{sec:18_4}

自动驾驶车辆通常搭载多种传感器（相机、LiDAR、雷达），完整的仿真系统需要统一建模和渲染这些多模态传感器数据。3DGS 原生仅支持针孔相机的光栅化渲染，将其扩展到 LiDAR、雷达和鱼眼相机等非标准传感器面临独特挑战。

\subsection{多传感器统一仿真}
\label{ss:18_6}

本节讨论2条技术路线：（1）LiDAR仿真与渲染；（2）多模态融合与新型传感器。

\subsubsection{LiDAR 仿真与渲染}

本路线介绍将 3DGS 从针孔相机渲染扩展到 LiDAR 点云仿真的方法，重新设计投影模型和渲染管线以适配激光射线的独特物理特性。

LiDAR 是自动驾驶的核心传感器之一，其点云数据具有与相机图像截然不同的特性：主动发射激光脉冲测量飞行时间、距离图（range image）表示、光线丢弃（ray-drop）现象、反射强度通道、旋转式 LiDAR 的滚动快门效应等。将 3DGS 从针孔相机渲染扩展到 LiDAR 仿真，需要重新设计投影模型（激光射线的锥形几何 vs 针孔相机的线性投影）和渲染管线（距离图 vs RGB 图像）。

\paragraph{可微分激光束泼溅的实时 LiDAR 仿真}

LiDAR-GS（2024.10）~\cite{2410_05111} 是首个将 3DGS 扩展到 LiDAR 点云仿真的工作，核心创新是可微分激光束泼溅（differentiable laser beam splatting）。直觉上，这将相机的"拍照"模式转换为激光的"测距"模式——原本投影到 2D 照片的操作变为沿激光射线计算精确距离，每个激光脉冲击中高斯椭球后返回的距离值构成距离图。

\textbf{Baseline 与问题。} 3DGS 的渲染管线基于针孔相机的局部仿射投影——将 3D 高斯投影到 2D 图像平面进行 $\alpha$-混合。但 LiDAR 的成像原理完全不同：主动发射激光脉冲，通过飞行时间（ToF）测量距离，生成距离图（range image）而非 RGB 图像。直接将 3DGS 的针孔投影用于 LiDAR 会产生局部仿射近似伪影，因为激光射线的锥形展开与针孔相机的线性投影在数学上不同。

\textbf{核心改动——可微分激光束泼溅。} LiDAR-GS 使用距离图表示，将激光射线投射到 3D 高斯的微截面（micro-facet）上实现精确表面泼溅。每个高斯基元除了标准的位置、协方差、不透明度外，还附加神经高斯表示——一个小型 MLP 集成视角依赖线索来建模 LiDAR 的关键属性：激光入射方向和外部因素（天气、表面材质）对回波的影响。

\textbf{三通道渲染。} LiDAR-GS 同时再仿真 LiDAR 的三个核心通道：（1）深度通道——通过激光束泼溅计算精确距离；（2）强度通道——建模激光反射强度随材质、距离和入射角的变化；（3）光线丢弃（ray-drop）通道——模拟激光被遮挡、吸收或 unfavorable 入射角导致的回波丢失。

\textbf{设计动机。} 3DGS 的 EWA 泼溅假设局部仿射投影（高斯在图像平面上的投影仍为 2D 高斯），这对针孔相机有效。但 LiDAR 射线在球坐标系中展开，射线的锥形几何使得 3D 高斯在距离图上的投影不再是简单高斯，需要沿射线求解射线-高斯交点而非投影到平面。神经高斯表示的引入解决了强度和光线丢弃的视角依赖建模难题：这两个通道无法用固定的 SH 系数表达，需要 MLP 动态预测。

\textbf{关键公式。} 激光射线 $\mathbf{r}(d) = \mathbf{o} + d\mathbf{v}$（原点 $\mathbf{o}$，方向 $\mathbf{v}$）与 3D 高斯 $\mathcal{N}(\boldsymbol{\mu}, \Sigma)$ 的交点通过最小化射线到高斯中心的马氏距离求解：

\begin{equation*}
d^* = \arg\min_d \, (\mathbf{r}(d) - \boldsymbol{\mu})^T \Sigma^{-1} (\mathbf{r}(d) - \boldsymbol{\mu})
\end{equation*}

深度图通过 $\alpha$-混合渲染：$D = \sum_{i=1}^N T_i \alpha_i d_i$，其中 $T_i = \prod_{j=1}^{i-1}(1-\alpha_j)$ 为透射率，$\alpha_i$ 为高斯 $i$ 沿射线的密度权重。强度通道由神经高斯 MLP 预测：$I_i = \text{MLP}(f_i, \theta_i, d_i)$（$f_i$ 为高斯特征，$\theta_i$ 为入射角，$d_i$ 为距离）。光线丢弃概率通过 sigmoid 输出建模。

\textbf{性能。} 消除局部仿射近似伪影，在 KITTI-360 和 nuScenes 上实现实时 LiDAR 点云再仿真，同时保持了与原始 LiDAR 数据一致的深度、强度和光线丢弃分布。

\paragraph{周期振动 2D 高斯的显式射线泼溅}

与 LiDAR-GS 使用 3D 高斯基元不同，GS-LiDAR（2025.01）~\cite{2501_13971} 提出基于 2D 高斯基元的全景 LiDAR 仿真框架，核心创新在于周期振动 2D 高斯表示和显式射线-泼溅求交的统一渲染。直觉上，2D 高斯圆盘如同无数个跳动的圆形薄膜拼接成运动物体，每个薄膜独立振动，接力完成复杂运动轨迹的表示。

\textbf{周期振动 2D 高斯。} 驾驶场景中车辆和行人持续运动，GS-LiDAR 使用 2D 高斯圆盘（而非 3D 椭球）作为基元，每个圆盘定义在局部切平面空间中（中心点、两个主切向量 $\mathbf{t}_u, \mathbf{t}_v$、缩放向量 $(s_u, s_v)$）。关键创新是为每个 2D 高斯引入周期振动属性——振动方向 $\mathbf{d}$、生命峰值 $p$ 和时间衰减率 $\gamma$：

\begin{equation*}
\boldsymbol{\mu}(t) = \boldsymbol{\mu} + \mathbf{d} \sin(2\pi(t - p)/l), \quad o(t) = o \cdot \exp(-\frac{(t - p)^2}{2\gamma^2})
\end{equation*}

每个基元执行简单周期运动，多个基元以接力方式组合即可表示动态元素的复杂运动。

\textbf{显式射线-泼溅求交。} LiDAR 的激光射线在全景坐标系中参数化，射线由两个正交齐次平面定义。射线与 2D 高斯圆盘的交点通过求解线性方程组直接计算 $(u, v)$ 坐标，然后从交点坐标精确计算距离 $r$。这避免了 3DGS 中局部仿射近似在 LiDAR 射线几何上的误差。

\textbf{三通道 LiDAR 渲染。} 每个高斯附加视图依赖的强度值和光线丢弃概率（均用球谐建模），通过 $\alpha$-混合渲染深度图、强度图和光线丢弃概率图。引入可学习的先验 $P_{prior}$ 建模 LiDAR 传感器本身的丢弃特性，并用 U-Net 全局精化光线丢弃掩码。

\textbf{性能。} 配合全景 LiDAR 监督，在 KITTI-360 上实现高保真动静统一的 LiDAR 点云仿真。

\paragraph{滚动快门建模的相机-LiDAR 联合渲染}

SplatAD（2024.11）~\cite{2411_16816} 是首个同时实现动态场景相机和 LiDAR 实时真实感渲染的 3DGS 方法，精确建模滚动快门、LiDAR 强度和光线丢弃等传感器特定现象。方法将场景分解为静态背景、天空环境和动态物体（支持刚性和可变形物体），动态物体通过 4D 傅里叶球谐建模时变外观，滚动快门通过逐行独立的相机外参建模，LiDAR 模拟通过可微泼溅同时渲染深度、强度和光线丢弃。渲染速度比 NeRF 方法快一个数量级，PSNR 分别提升 +2（NVS）和 +3（重建）。

\paragraph{因子化表示与自动分块的多传感器渲染}

SimULi（2025.10）~\cite{2510_12901} 与 Uni-Gaussians 的混合光栅化-光线追踪策略不同，SimULi 基于 3DGUT 的无迹变换（Unscented Transform）统一处理任意相机模型和 LiDAR 的实时渲染——两者共享同一套粒子表示和泼溅管线。

\textbf{Baseline 与问题。} 现有方法要么为不同传感器维护独立表示（Uni-Gaussians 分别用光栅化和光线追踪），要么强行统一到单一渲染模式（SplatAD 统一用泼溅）。前者增加存储和同步开销，后者牺牲 LiDAR 的仿真精度。同时，相机数据（稠密 RGB）和 LiDAR 数据（稀疏深度+强度）的分布特性差异巨大，联合优化时一种模态的质量往往被另一种模态的不一致性拖累。

\textbf{因子化 3D 高斯表示。} 相机属性集 $G_c$（3 阶球谐编码视角相关颜色）和 LiDAR 属性集 $G_l$（3 阶球谐编码视角相关强度和光线丢弃概率）使用独立的粒子集但共享位置 $\mu_i$、不透明度 $o_i$ 和协方差 $\Sigma_i$。这种因子化解耦使每个传感器模态的重建质量不受另一模态的传感器不一致性影响——避免了 UniSplat 等方法中跨传感器数据不一致时必须优先某一模态质量的权衡。最近邻锚定损失 $L_{anchor} = \frac{1}{n}\sum_{i \in G_c} \|\mu_i - \text{NN}(\mu_i, G_l)\|_2$ 鼓励相机高斯靠近 LiDAR 监督的表面几何，但允许相机高斯在跨传感器不一致区域自由调整，将相机和深度误差降低最多 40\%。

\textbf{LiDAR 自动分块与光线剔除。} LiDAR 的测量分布极不均匀——方位角覆盖 $360^\circ$ 但仰角通常仅 $20^\circ$，使用相机的等尺寸瓦片会极其低效。SimULi 根据仰角累积分布函数自动计算均衡化分块边界，使每个瓦片包含大致相同数量的测量点。光线剔除使用更密集的瓦片分辨率（$r_\phi^d = 1600, r_\theta^d = 8$）构建 2D 射线掩码和求和面积表，在常数时间内（4 次内存读取 + 3 次算术运算）判断高斯是否与任何射线相交，大幅减少不必要的粒子评估。比光线追踪方法快 10--20 倍。

\paragraph{光栅化-光线追踪混合的多传感器仿真}

Uni-Gaussians（2025.03）~\cite{2503_08317} 提出混合方法——光栅化渲染图像数据，高斯光线追踪渲染 LiDAR 数据——统一动态驾驶场景的多传感器仿真。与 SplatAD 通过可微泼溅统一相机和 LiDAR 的渲染管线不同，Uni-Gaussians 的核心洞察是相机和 LiDAR 的成像原理本质不同（相机是被动接收光线的光栅化过程，LiDAR 是主动发射激光的光线追踪过程），强行统一到单一渲染模式会损害各自的仿真质量。因此该框架在场景表示层面统一（所有动态参与者用 2D 高斯基元建模，附加 SH 强度和 SH 光线丢弃概率属性），但在渲染层面分流：相机渲染使用标准 2D-to-2D 投影和 $\alpha$-混合光栅化；LiDAR 渲染将 2D 高斯转换为三角面片构建 BVH，通过 OptiX 光线追踪计算射线-高斯交点，逐射线前向后积分深度、强度和光线丢弃概率。场景图建模沿用 OmniRe 的静态/刚性/非刚性分解，同时解决了几何表面重建中 3D 高斯的体积表示与薄表面之间的冲突。

\subsubsection{多模态融合与新型传感器}

本路线介绍融合 LiDAR 点云与相机图像互补信息的多模态重建方法，以及雷达、鱼眼相机等新型传感器的仿真探索。

除了单独处理 LiDAR 或相机，多模态融合方法旨在利用 LiDAR 点云和相机图像的互补信息（LiDAR 提供精确几何但缺乏纹理，相机提供丰富颜色但深度不确定）实现更高质量的重建和仿真。同时，雷达、鱼眼相机等新型传感器的仿真也在逐步探索中。

\paragraph{网格-八叉树双表示的紧耦合 LiDAR-相机}

TCLC-GS（2024.04）~\cite{2404_02410} 设计紧耦合 LiDAR-相机高斯泼溅，利用 LiDAR-相机数据派生的混合显式（着色 3D 网格）和隐式（层次八叉树特征）3D 表示丰富高斯属性。显式着色网格通过投影为每个高斯提供完整的形状和颜色初始化，层次八叉树在多尺度上查询空间特征为高斯赋予更广泛的场景上下文，两种表示通过注意力机制融合后解码高斯属性。在 Waymo 和 nuScenes 上以 90/120 FPS 实现实时 RGB 和深度渲染。

\paragraph{LiDAR 可见性率建模的高速公路重建}

LiHi-GS（2024.12）~\cite{2412_15447} 首次聚焦高速公路动态场景——稀疏传感器视角和单调背景的挑战场景。核心设计是为每个高斯引入与相机不透明度解耦的 LiDAR 可见性率 $\alpha_i^L = \alpha_i \cdot \beta_i$，以建模半透明表面和远距离目标的 LiDAR 感知差异（VLS-128 LiDAR 对 180m 外目标检测率仅 5\%）。LiDAR 渲染将高斯从笛卡尔坐标转换到球坐标后泼溅为距离-角度图，并设计 2D 高斯尺度补偿机制恢复因数值不稳定而丢失的远距离小高斯。深度不确定性渲染基于 Welford 在线算法沿光线增量计算深度方差，过滤边缘悬浮伪影，Camera-LiDAR 解耦位姿优化分别处理两传感器的运动目标位姿偏差。

\paragraph{伪外推点云监督的离轨 LiDAR 仿真}

LiDAR-EVS（2026.03）~\cite{2603_14763} 提出轻量级框架增强沿未驾驶轨迹的 LiDAR 仿真。伪外推视角点云监督通过四步生成：多帧 LiDAR 融合（聚合不同时刻的点云扩展空间覆盖）、视图变换（将点云投影到外推视角的坐标系）、遮挡卷曲（处理外推视角中新暴露区域的缺失点云）、强度调整（根据新视角的激光入射角重新计算反射强度）。空间约束 dropout 正则化在外推区域随机丢弃点云，防止模型过拟合到不完整的几何。

\paragraph{多路径噪声建模的雷达高斯泼溅}

RadarSplat（2025.06）~\cite{2506_01379} 首次将高斯泼溅与雷达噪声建模结合，实现真实雷达数据合成和增强 3D 重建。雷达成像具有独特的噪声特性——多路径反射（信号在多个表面间弹跳产生虚假回波）、稀疏性（每个角度仅返回少量回波点）和角度模糊（水平分辨率低）。RadarSplat 为每个高斯基元附加雷达特定的属性建模，包括雷达散射截面积（RCS）和多路径概率，通过可微分泼溅渲染雷达距离-角度图。雷达在恶劣天气下具有独特优势，该方法比当时方法雷达图像合成提升 +3.4 PSNR / 2.6 倍 SSIM，几何重建误差降低 40\%。

\paragraph{高斯转网格的 GPU 光线投射超快仿真}

FGGS-LiDAR（2025.09）~\cite{2509_17390} 与 LiDAR-GS 和 GS-LiDAR 直接用高斯基元泼溅渲染 LiDAR 数据不同，FGGS-LiDAR 采用间接策略——先将任意预训练 3DGS 模型转换为高保真水密网格，再通过 GPU 光线投射实现超快 LiDAR 仿真。

\textbf{三阶段转换管线。} 第一阶段：BVH 加速体积离散化——对预训练高斯场构建包围盒层次结构（Morton 码排序 + 基数排序实现近线性构建），将 3D 空间划分为均匀体素，利用 BVH 快速查询每个体素中心的高斯密度 $D(v) = \sum_{G_i \in C_{tile}} \exp(-\frac{1}{2}(v - \mu_i)^T \Sigma_i^{-1}(v - \mu_i)) \cdot f(o_i)$，生成二值占用体。表面掩码 $Surf(i,j,k) = V_{i,j,k} \land \neg Int(i,j,k)$ 剔除完全被内部体素包围的体素，仅保留表面附近的高斯。

第二阶段：窄带 TSDF 网格提取——对二值占用体进行高斯核去噪后，从外部填充帧播种泛洪填充确定外部区域，6-邻域传播计算符号距离场并截断到窄带半径 $r$，Marching Cubes 提取等值面后再经结构优先简化和非收缩平滑得到最终水密网格。

第三阶段：GPU 加速光线投射——每条 LiDAR 射线分配一个 GPU 线程，遍历网格的 BVH 结构，利用最佳深度提前终止裁剪不可见节点和三角形。结构化数组内存布局和共享内存缓存 BVH 节点保证合并内存访问，warp 同步遍历减少分支发散。即使对百万三角面片的网格，仿真延迟仍在毫秒级，实现超过 500 FPS 的 LiDAR 回波仿真，且不需要对原始 3DGS 模型做任何修改——即插即用。

\paragraph{稀疏体素局部隐式场的双模式渲染}

SaLF（2025.07）~\cite{2507_18713} 提出支持光栅化和光线追踪的统一多传感器仿真体积表示——稀疏 3D 体素基元，每个体素是局部隐式场。相机渲染使用体素内的局部隐式场预测高斯属性后进行光栅化泼溅，LiDAR 渲染使用光线追踪与体素内的隐式表面求交计算精确距离。这种双模式设计避免了为不同传感器维护独立表示的开销，训练少于 30 分钟，相机渲染 50+ FPS，LiDAR 渲染 600+ FPS。

\paragraph{全向 RGB-LiDAR 转换}

Bae et al.（2026.03）~\cite{2603_06061} 提出全向 RGB-LiDAR 确定性转换管线，将存档的传感器日志转换为 3DGS 初始化资产。关键技术包括：ERP 到立方体图投影解决全景图极点严重非线性畸变导致的 SfM 特征匹配失败；PRISM（Color-Stratified Point Cloud Sampling）下采样策略将分层域从空间覆盖切换到颜色空间 $C$，按颜色分箱设置最大点容量 $k$，在大幅削减视觉均匀区域点云的同时保留纹理丰富区域用于球谐初始化。多模态对齐绕过容易发散的全局配准搜索，用轨迹元数据初始化鲁棒 ICP 局部优化求解刚性变换 $(R^*, t^*)$，桥接尺度模糊的 SfM 稀疏点云与 LiDAR 点云之间的坐标系差异。

\paragraph{高斯泼溅作为视图变换器的 BEV 融合}

GaussianCaR（2026.02）~\cite{2602_08784} 将高斯泼溅重新用作高效通用视图变换器，将图像像素和雷达点映射到公共 BEV 表示。图像分支为每个像素预测深度分布并提升为 3D 高斯，雷达分支直接利用雷达的距离信息初始化高斯位置，两种模态的高斯在 BEV 平面上通过可微分光栅化泼溅融合为统一特征图。在 nuScenes 上 BEV 分割与当时方法持平但推理速度快 3.2 倍。

\begin{takeaway}
传感器融合与定位将 3DGS 从纯视觉扩展到多模态：鱼眼相机高斯变换和 LiDAR 仿真使单一 3DGS 模型统一渲染多种传感器数据，高斯泼溅作为视图变换器实现高效的 BEV 多模态融合。
\end{takeaway}

%=====================================================================
\section{3D占用估计与感知}
\label{sec:18_5}

3D 语义占用预测是自动驾驶感知的核心任务之一——将周围环境细粒度划分为体素级语义标签。传统方法采用密集体素网格，计算和内存开销随分辨率立方增长。3DGS 的稀疏高斯基元提供了一种天然高效的替代表示，催生了将高斯泼溅引入占用预测的新范式。

\subsection{高斯辅助感知范式}
\label{ss:18_7}

本节讨论2条技术路线：（1）高斯基元用于占用预测；（2）高斯辅助感知范式。

\subsubsection{高斯基元用于占用预测}

本路线介绍用 3DGS 的稀疏高斯基元替代密集体素进行 3D 语义占用预测的方法，大幅降低表示成本并保持感知精度。

3D 语义占用预测将周围环境划分为密集体素网格并赋予语义标签，是自动驾驶感知的核心任务。传统方法使用密集体素表示，计算和内存开销随分辨率立方增长。3DGS 的稀疏高斯基元天然适合替代密集体素——大多数空间区域是空的，只需在有物体表面的位置放置高斯，大大降低了表示成本。

\paragraph{泼溅投影尺度感知与体素快速渲染的占用估计}

GaussianOcc（2024.08）~\cite{2408_11447} 系统研究高斯泼溅在全自监督 3D 占用估计中的两种用法，首次将 3DGS 的投影和渲染能力引入占用预测的监督流程。直觉上，这相当于只在场景中"有东西"的位置放置高斯基元，而非填满整个 3D 体素网格——空无一物的格子被丢弃，大幅减少计算冗余。

\textbf{Baseline 与问题。} 全自监督 3D 占用估计面临两个核心挑战：（1）尺度歧义——单目深度预测的尺度不确定，需要跨视图约束获取真实尺度，但现有方法依赖 SfM 提取稀疏深度（耗时且不稳定）或时空上下文（性能有限）；（2）体渲染效率低——基于密集采样的体渲染在 $180 \times 320$ 分辨率下需要采样 1.08 亿个点，但实际优化的 3D 体素顶点仅 216 万个，存在约 50 倍的采样冗余。

\textbf{核心改动——两阶段高斯泼溅。}

\textit{阶段 1：高斯泼溅投影（GSP）用于尺度感知训练。} 在相邻相机的重叠区域，将一侧重叠区域遮蔽后，利用 3DGS 的泼溅投影渲染相邻视角。具体地，深度网络预测每个像素的高斯属性（深度、缩放、旋转），将 3D 高斯泼溅到相邻视角后与原图比较——如果深度准确则渲染结果应与原图一致，从而提供精确的尺度监督信号。关键设计是遮蔽重叠区域的一侧，强制模型通过泼溅投影恢复完整场景，迫使深度预测学习正确的尺度。

\textit{阶段 2：体素空间高斯泼溅（GSV）用于快速渲染。} 将 3D 体素网格的顶点直接视为 3D 高斯的位置，使用固定缩放和旋转，仅优化语义和不透明度。这利用了高斯泼溅的快速光栅化渲染替代体渲染中密集采样，在空体素区域网络自动将不透明度预测为零，不影响最终结果。

\textbf{性能。} GSP 实现无需真值位姿的尺度感知训练；GSV 使训练速度提升 2.7 倍、渲染速度提升 5 倍，在 nuScenes 和 DDAD 上实现全自监督占用估计的最优性能。

\paragraph{动静双查询的层次粗到细占用预测}

ODG（2025.06）~\cite{2506_09417} 与 GaussianOcc 使用统一高斯查询不同，ODG 基于驾驶场景可分解为静态和动态部分的观察，设计双高斯查询更好地建模多样场景物体。

\textbf{与 GaussianOcc 的差异。} GaussianOcc 将体素顶点直接视为高斯位置，使用固定缩放和旋转；ODG 定义两组独立查询——静态查询 $G^s$ 在固定 BEV 网格位置均匀初始化（适合捕获道路、建筑等静态结构），动态查询 $G^d$ 在检测到的 3D 边界框中心初始化（适合捕获车辆、行人的局部细节）。关键扩展是为动态查询增加 3D 检测框属性 $\mathbf{b} = [l, w, h, \theta, v_x, v_y, v_z]$（长宽高、朝向、速度），使动态高斯同时承担占用预测和物体检测的双重任务。

\textbf{层次粗到细预测。} 不同于前述方法使用固定数量的高斯，ODG 通过 $L$ 层变换器逐层增加每查询的高斯数量（$K_1 < K_2 < \cdots < K_L$），每层预测高斯中心、语义类别和（仅动态查询的）检测框属性。动静交叉注意力（DaS Attention）将两组查询的特征拼接后执行全自注意力，实现动静信息的双向增强——静态查询为动态查询提供场景上下文，动态查询为静态查询补充物体级语义。

\textbf{渲染监督。} 利用 3DGS 实时泼溅渲染深度图和语义图，与 LiDAR 投影的真值比较施加 $L_1$ 深度损失和交叉熵语义损失，注入像素级对齐信号。检测框通过匈牙利算法标签分配和 $L_1$ 损失监督。

\paragraph{几何-语义双图注意力的占用预测}

GraphGSOcc（2025.06）~\cite{2506_14825} 与 ODG 使用双查询分别建模动静不同，GraphGSOcc 的核心创新是在高斯之间建立显式图结构增强空间交互——双高斯图注意力（DGGA）从几何和语义两个维度建模高斯间的关系。

\textbf{双高斯图注意力。} 几何图根据高斯中心的欧氏距离构建 KNN 邻域（Top-K 最近邻），自适应计算搜索半径使大尺度高斯聚合更广邻域特征——覆盖更大的空间范围，捕获道路和建筑等大面积结构的全局上下文。语义图根据高斯特征的余弦相似度保留 Top-M 最相关节点，显式编码实例内外语义关系——将同一物体（如同一辆车）的高斯关联起来，即使它们在空间上不相邻。两个图分别计算注意力后通过自适应融合层加权合并，权重由特征驱动的门控机制动态决定。多尺度图注意力使用不同 K/M 配置（[100, 75, 50, 20]）构建层次化图结构，在多尺度上融合几何和语义信息。

\textbf{动静解耦注意力。} 根据语义分数将高斯分为动态和静态两组后，双向交叉注意力实现信息交互：动态交叉注意力（DCA）利用静态场景的先验布局改善动态物体的运动预测；静态交叉注意力（SCA）捕获动态物体与静态环境的交互（如行人与人行道、车辆与道路）。mIoU 达 25.20\% 且 GPU 内存仅 6.8 GB，显著优于未解耦的基线。

\paragraph{高斯-视图双模式空间聚合的占用预测}

ST-GS（2025.09）~\cite{2509_16552} 与 GraphGSOcc 从图结构增强空间交互不同，ST-GS 从注意力机制的参考点设计出发，提出引导知情空间聚合（GISA）策略——通过双模式注意力（高斯引导注意力 GGA + 视图引导注意力 VGA）提供互补的 3D 空间采样偏移。

\textbf{双模式空间聚合。} 高斯引导注意力（GGA）利用每个高斯固有的均值和协方差属性生成采样偏移——在局部高斯坐标系中初始化 $M$ 个结构化 3D 网格点作为基础偏移 $P_L^G$，通过可学习缩放因子 $s^G$ 调节采样半径，再经旋转 $R^G$ 和缩放 $S^G$ 变换到感知坐标系，使偏移方向与高斯椭球的主轴对齐。视图引导注意力（VGA）沿相机视线方向生成偏移——在局部视线坐标系的 y-z 平面初始化 2D 网格偏移，根据高斯中心的方位角 $\phi$ 构建旋转矩阵 $R^V(\phi)$ 将偏移变换到感知坐标系，利用多视角重叠区域的空间-语义连续性增强跨视角信息聚合。门控空间特征聚合（GSFA）通过注意力驱动的门控 $\sigma_S \in [0,1]^{K \times M}$ 动态平衡两种偏移的贡献。

\textbf{几何感知时间融合。} 利用自车运动信息建立帧间几何对应——将当前帧的参考点通过刚体变换 $T^{\hat{t} \to t}$ 对齐到历史帧坐标系，然后门控时间特征融合（GTFF）模块通过轻量 MLP 预测自适应融合门控 $\sigma_T \in [0,1]^{K \times D}$，选择性聚合历史信息抑制遮挡或动态物体引入的不一致特征。在 nuScenes 上达到当时最优的时间一致性。

\paragraph{光线补偿与动态损失重加权的弱监督占用}

GSRender（2024.12）~\cite{2412_14579} 与 GaussianOcc 同样将 3DGS 用于弱监督占用估计，但针对不同核心问题：GaussianOcc 关注尺度歧义和渲染效率，GSRender 聚焦于射线上的重复预测问题。

\textbf{光线补偿模块。} 在体渲染中沿射线累积不透明度时，多个高斯可能重复覆盖同一体素位置，导致该位置的语义标签被多个高斯的不透明度加权预测污染。GSRender 的光线补偿模块引入多视角补偿策略——利用相邻时间步的占用特征 $P_{adj}$ 通过自车运动变换 $O_{adj}^{curr\_ego} = T_{adj\_ego}^{curr\_ego} O_{adj}^{adj\_ego}$ 将相邻帧的特征对齐到当前帧坐标系，双线性采样获取对应位置特征后通过共享高斯头预测高斯属性。这为当前帧中被遮挡的体素提供了来自其他时间步的补充视角，降低重复预测。

\textbf{动态损失重加权。} 在相邻帧的损失计算中，为动态物体类别分配较低权重 $\alpha_{s(i)}$（通过可学习掩码识别动态区域），因为无法准确估计动态物体的运动轨迹会干扰网络优化——相邻帧中动态物体的位置变化不应作为当前帧占用预测的可靠监督信号。最终的动态损失 $L_{adj}^seg = -\sum_{i=1}^n \alpha_{s(i)} \omega_{s(i)} S_{c(i)} \log(\hat{p}_{pix(i)})$ 通过动态权重和类别平衡权重的双重调制，有效抑制动态物体对跨帧占用一致性的干扰。RayIoU 提升 +6.0。

\subsubsection{高斯辅助感知范式}

本路线介绍利用 3DGS 渲染能力为下游感知任务提供辅助监督信号的方法，实现从"重建工具"向"感知基础设施"的角色转变。

除了直接用于占用预测，3DGS 的渲染能力还可以为下游感知任务提供辅助监督信号——通过渲染语义图、深度图或 BEV 特征图，可以在不需要昂贵 3D 标注的情况下训练感知模型。这一方向代表了 3DGS 从"重建工具"向"感知基础设施"的角色转变。

\paragraph{显式高斯正交投影的 BEV 感知}

Splat2BEV（2026.03）~\cite{2603_19193} 与 GaussianLSS 将深度不确定性引入 Lift-Splat-Shoot 范式不同，Splat2BEV 的核心思路是先预训练一个前馈高斯生成器显式重建 3D 场景，再将重建结果正交投影到 BEV 空间，利用显式 3D 几何为 BEV 感知提供几何对齐的特征表示。

\textbf{Baseline 与问题。} 现有 BEV 感知方法（如 BEVFormer、LSS）通过隐式 Lift-Splat 操作将 2D 图像特征提升到 3D 空间，但深度估计不准确导致 BEV 特征在空间上错位——远距离物体的特征被投影到错误位置，严重影响分割精度。GaussianLSS 引入深度不确定性部分缓解了这一问题，但本质仍是隐式反投影，缺乏显式 3D 几何约束。

\textbf{设计动机。} Splat2BEV 的核心洞察是：3DGS 的显式几何表示天然适合 BEV 投影——高斯基元具有精确的 3D 位置和协方差，正交投影到 BEV 平面时无需深度估计的中间步骤，直接产生几何对齐的特征。相比隐式 Lift-Splat 的多假设深度采样，显式高斯投影在几何精度和计算效率上均具优势。三阶段训练策略确保高斯生成器先学到可靠的几何，再服务于下游感知任务，最后端到端微调使两者协同优化。

\textbf{三阶段训练。} 第一阶段预训练高斯生成器 $G_\theta$——从多视角透视图像预测深度图后反投影为 3D 高斯位置，配合 Metric3Dv2 深度基础模型提供参考深度监督（$L_1$ 损失 + SILog 尺度不变损失）。每个高斯附加可学习特征向量 $f_j \in \mathbb{R}^C$，通过冻结 DINO 编码器蒸馏密集语义先验（余弦相似度损失），使高斯特征兼具几何接地性和语义含义。第二阶段冻结高斯生成器，通过正交投影将高斯投影到 BEV 平面（3D 协方差 $\Sigma_{3D}$ 经正交雅可比 $J_{ortho}$ 变换为 2D 协方差 $\Sigma_{2D} = J_{ortho} \Sigma_{3D} J_{ortho}^T$），经 BEV 编码器和分割头训练下游任务。第三阶段端到端微调所有组件，使 BEV 分割损失的梯度回传至高斯生成器，让 3D 重建适应任务级目标。在 nuScenes 和 Argoverse 上验证了显式 3D 重建对 BEV 感知的显著提升。

\paragraph{深度不确定性高斯泼溅的 BEV 感知}

GaussianLSS（2025.04）~\cite{2504_01957} 回顾反投影方法，在 Lift-Splat-Shoot 范式中引入深度不确定性建模——学习软深度均值和方差隐式捕获物体范围，将深度分布转换为 3D 高斯并通过光栅化构建不确定性感知 BEV 特征。具体地，每个像素的深度分布 $\mathcal{N}(\mu_d, \sigma_d^2)$ 被离散化为一组 3D 高斯，沿射线放置在等间距深度位置，以对应深度概率值作为不透明度，然后通过可微分光栅化泼溅到 BEV 平面。这避免了 Lift-Splat-Shoot 中沿深度轴的穷举外积（frustum pooling），比投影方法快 2.5 倍、内存仅 0.3 倍。

\paragraph{自适应邻域与体渲染监督的跨模态占用}

TACOcc（2025.05）~\cite{2505_12693} 与 Splat2BEV 仅利用高斯泼溅的投影能力不同，TACOcc 在图像-LiDAR 跨模态占用预测中引入目标尺度自适应的双向特征融合，并利用 3DGS 体渲染管线施加光度一致性监督。

\textbf{目标尺度自适应检索。} 核心挑战是最优邻域大小 $k$ 的选择——大目标需要大邻域聚合足够的上下文信息（如整辆车的结构），小目标需要小邻域抑制噪声干扰（如远处行人的精细特征）。TACOcc 通过 MLP 将当前查询特征映射为候选 $k$ 值的离散分布，利用 Gumbel-Softmax 重参数化建立连续松弛机制（使离散采样可微分），结合直通估计器（STE）在前向传播使用离散值、反向传播使用连续值，实现端到端可学习的最优 $k$ 选择。基于选定的 $k$，对每个有效查询体素在跨模态对齐特征空间中用 KNN 采样非零目标体素作为键向量，通过全连接层和注意力加权增强原始查询特征，对称架构在另一模态端执行反向增强，实现双向跨模态互补。

\textbf{3DGS 体渲染监督。} 融合特征通过轻量 3D 卷积网络映射为高斯基元参数，构建高斯场时同时使用点云坐标和非零体素中心（由融合特征过滤）作为锚点——点云提供高精度几何但可能稀疏，体素坐标以规则网格弥补缺失区域，两者互补防止覆盖盲区。可微分渲染将 3D 高斯投影到六视角 2D 图像空间，光度一致性损失驱动高斯参数迭代优化。收敛后的参数与初始参数计算参数一致性损失 $L_{pc} = \sum_{i=1}^N \|\theta_i^{final} - \theta_i^{(0)}\|_1$，补偿 2D 输入图像在融合 3D 多模态特征中可能遗漏的表面细节。

\paragraph{高斯泼溅正则化增强占用与场景流}

VoxelSplat（2025.06）~\cite{2506_05563} 利用 3DGS 作为正则化框架增强占用和场景流预测，核心思路是将占用预测网络解码出的稀疏语义 3D 高斯投影到 2D 相机视角，与输入图像的语义分割结果比较提供额外语义监督。同时使用预测的场景流建模高斯运动（将 $t$ 时刻的高斯按场景流偏移到 $t+1$ 时刻，渲染后与 $t+1$ 的观测比较），实现自监督场景流学习。整个 3DGS 正则化模块仅在训练时使用，可无缝集成到现有占用模型不增加推理时间。

\paragraph{占用网格初始化的语义驱动动静分离}

OG-Gaussian（2025.02）~\cite{2502_14235} 用占用预测网络从环视相机图像生成的占用网格（OGs）替代 LiDAR 点云作为 3DGS 初始化。利用 OG 中的语义信息自动分离动态车辆（语义标签为"车辆"的体素）和静态背景（其余体素），无需依赖 3D 检测器的边界框标注。动态物体的轨迹和姿态通过可学习的位姿参数以端到端方式估计，消除人工标注需求，平均 PSNR 达 35.13、渲染速度 143 FPS。

\paragraph{协作语义占用预测}

Chen et al.（2025.08）~\cite{2508_10936} 首次利用稀疏 3D 语义高斯泼溅进行协作 3D 语义占用预测。核心设计利用高斯基元的稀疏性实现高效的跨车辆通信——每辆车将本地预测的高斯基元（位置、语义、不透明度）而非密集体素网格传输给邻居车辆。接收方通过邻域跨智能体融合将来自不同车辆的高斯基元投影到统一的 3D 空间，几何语义联合编码和稀疏对象中心消息传递确保融合后的高斯具有一致的语义属性，mIoU 提升 +8.42。

\begin{takeaway}
3D 占用估计与感知利用 3DGS 的稀疏性替代密集体素网格：占用网格初始化实现语义驱动的动静分离，稀疏高斯跨车辆通信将协作占用预测的传输开销降至密集体素的极小比例。
\end{takeaway}

%=====================================================================

\section{近期研究（2026 年 8--11 月）}
\subsection{街景前馈重建与几何一致性}

环视前馈重建的难点分布在三个位置：输入特征要提供几何先验，预测时要把原语分配到有用的表面，训练时还要把相机外观和几何分开。以下方法分别从这三个位置切入。

\paragraph{VGGD：把基础模型的几何先验放到环视前馈网络前端}
2026.08~\cite{2608_10682}。单时刻环视图像的相机间重叠很少，几何不确定性很难只靠增加解码器容量解决。VGGD 从预训练 VGGT 提取多视图几何特征，再分别处理几何一致性与外观补全，并在训练早期渐进引入尺度约束。预测头结合像素和体积两种参数化，输出用于新视图渲染的高斯。

论文在 nuScenes 单帧协议中评价画质，并以相对于 Depth Anything V2 伪深度的 Pearson 相关评价结构。该相关指标对全局平移和正比例缩放不敏感，因此支持相对深度一致性的比较，无法单独证明米制深度准确。

\paragraph{LGS：用删点和加点干预学习高斯密度控制}
2026.08~\cite{2608_11077}。LiDAR 引导的前馈重建常让每个观测点初始化一个高斯，重建质量却未必随初始化密度均匀分布。LGS 在训练阶段模拟增加或删除基元，用局部的梯度响应作为结构调整的训练信号，使推断时能够按需要改变高斯数量。

论文在 Waymo 与 PandaSet 上评价，并将动态区域和整张图的画质及深度误差分别列出。与逐场景增密相比，该方法将增密决策转化为训练得到的前馈策略，其在不同扫描模式或数据域上的表现需要重新验证。

\paragraph{SG-GS：用语义区域和边界调节城市高斯分配}
2026.10~\cite{2610_16292}。复杂街景里统一的增密阈值可能忽略细小构件和边界。SG-GS 将 SAM 分割得到的语义区域引入高斯密度控制，在不同区域动态分配梯度阈值，并在结构边界加强增密。该方法也用语义约束组织优化损失，使结构细节与原语数量的分配相互关联。

在 Derrick 数据上，相同划分和训练设置下，作者比较原始 3DGS 并报告 PSNR 提升 1.34 dB。语义分区质量决定阈值能否覆盖真实细部，渲染指标还需与几何误差或下游驾驶指标分别验证。

\paragraph{MCSplat：把环视外观差异从几何监督中分离}
2026.10~\cite{2610_25345}。车载相机的曝光和图像处理差异会让同一表面在不同相机里呈现不同颜色。若重建模型直接降低观测图像之间的光度误差，几何可能被迫吸收这部分变化。MCSplat 利用深度引导的跨相机对应构建统一外观，并用多尺度双边网格 Transformer 把统一外观的渲染结果恢复到相机各自的观测外观。这个设计保留从真实图像得到的训练信号，在外观映射中解释相机差异。

Waymo 协议以时刻 t 的图像评价输入视图重建，以 t+1 的图像评价新视角合成。跨视图深度和法向一致性另行评价。作者也在 nuScenes 测试迁移表现，并指出训练集外观有限时，少数场景中恢复的相机域外观仍会偏色。跨车道外推示例提供定性证据。

\subsection{相机观测退化与天气}

恢复受天气或挡风玻璃影响的观测，需要先区分污染附着在哪个坐标系。可以绑定在相机，也可以随场景深度变化；仿真则要控制这些效应何时进入渲染。

\paragraph{GSRAIN：在街景高斯中分别校准雨滴与整体雨效}
2026.08~\cite{2608_02177}。街景渲染的雨效同时包含快速变化的局部雨滴和低频的整体外观。GSRAIN 根据实测降雨数据构建高频雨滴模型，用受几何引导的单步扩散模型生成低频雨效，再把两者融合进同一高斯场景。作者给出的雨强控制区间为 0 至 13 毫米每小时。

论文的总体雨效对比在 Waymo 上用 8.5 毫米每小时的模拟设定，报告 FID 149.09，对照 CycleGAN-Turbo 与 WeatherEdit 分别为 155.71 和 157.94。整体雨强可控范围与这组图像质量比较的具体雨强属于不同条件。检测和闭环结果反映所测场景、雨量与驾驶模型的组合。

\paragraph{DerainSplat：从稀疏雨天观测恢复可渲染的干净场景}
2026.08~\cite{2608_02191}。与在已有高斯场景中施加降雨的任务相比，DerainSplat 从少量受天气污染的图像出发，把天气因素与场景分开估计。天气网络预测阴天光照、随深度变化的雾、雨线和镜头雨滴，由此生成两类支持图。一类调整跨视图代价体的匹配，另一类在深度对齐后融合用于修复外观的证据。低几何支持的位置还会削弱高斯不透明度，以抑制伪结构。

论文对合成雨天输入做定量比较，也在真实雨天 HydroViews 场景上展示无微调结果。真实雨天没有与之对齐的干净目标，因此这部分主要用定性结果评价；它与合成雨天的成对指标应分开解读。

\paragraph{RIDE-Splat：在高斯重建之外解释挡风玻璃反射}
2026.10~\cite{2610_34632}。透过挡风玻璃拍摄时，外部街景随相机运动产生视差，座舱反射却大体跟着相机。RIDE-Splat 先联合估计世界坐标下的三维高斯和一张相机坐标下共享的有符号 RGB 图，让后者解释反射增亮与透射衰减。随后去掉反射成分，再用处理后的观测拟合干净场景。

推断时只渲染去反射的高斯场景。评价区分纯合成、真实场景加合成反射和真实挡风玻璃反射。真实的无反射参考来自另一次采集，时序和环境难以完全对齐；强烈独立运动的反射也超出共享图近似的适用范围。

\subsection{轨迹外推与视频修复}

高斯渲染在日志轨迹附近通常有较多测量支撑，轨迹外推后可见范围改变，修复模块需限制对可靠原图内容的改写。

\paragraph{SPVC：按空间条件修复整段驾驶渲染}
2026.08~\cite{2608_17420}。新轨迹和车辆插入会同时破坏背景结构、前景外观和帧间连续性。SPVC 在视频级处理这些伪影，并用相机位姿、三维框和高精地图约束修复内容。训练样本用欠约束渲染及插入车辆产生的退化构造，让模型同时见到背景扭曲与前景接缝。

作者在 nuScenes、Waymo 和 PandaSet 上检查新视角渲染修复及资产插入伪影。位姿、三维框或地图不准会影响结构约束；修复后图像的观感还需与原始几何和下游任务的一致性分开评价。

\paragraph{PRG-Fusion：按重建证据决定保留、修复或生成}
2026.09~\cite{2609_06948}。沿日志以外的轨迹渲染驾驶视频时，同一图像可能同时包含可信的重建表面、受退化影响但仍有测量支撑的区域，以及完全缺少观测的区域。PRG-Fusion 根据三维高斯渲染和几何证据按像素划分这三类区域。可信部分保留高斯外观，中间部分用 LiDAR 投影约束结构，缺少支撑的部分交给视频先验补全。

作者在 Waymo 验证集选取八个含较多动态交通参与者的片段进行评价。方法依赖预先重建与对齐的 LiDAR；比较轨迹外推时，应分别报告可由重建或测量支持的区域，以及由生成模型补齐的区域。

\subsection{动态实例与闭环仿真}

可独立编辑的交通参与者需要同时具备可见表面、实例布局和可控运动。实例重建、不可见面补全和行为合成承担不同的工作。

\paragraph{DecoupleGS：把动态交通参与者做成可重用资产}
2026.08~\cite{2608_01761}。闭环测试需要按规划器输出调整自车和周边车辆的轨迹。DecoupleGS 将静态街景与规范坐标系下的动态交通参与者分别表示。动态资产通过感知裁剪和向量量化控制渲染成本，跟随地图拓扑进行几何注册，再以代理重光照减小资产和背景的外观差异。

实验以 nuScenes 和 PandaSet 片段作为背景，从 3DRealCar 选取可操纵的车辆资产，分别评估渲染效率、几何与光度一致性、日志回放及闭环测试。代理重光照是效率优先的近似，地图注册依赖语义拓扑，当前实验主要覆盖刚性交通参与者和相机仿真。

\paragraph{USR-Drive：共同去噪动态高斯与实例三维框}
2026.08~\cite{2608_19036}。驾驶重建输出稠密、可渲染的高斯，感知输出可以描述实例的三维框。USR-Drive 在有位姿环视视频的条件下联合生成两类三维状态。高斯向框提供稠密几何线索，框向动态高斯提供实例级布局约束。

实验采用 nuScenes 六路环视相机的八帧片段，并在 VKitti 的 400 个样本上做零样本评价。位姿用于统一位置编码的米制锚定，不能把这一条件描述为任意无标定视频。渲染重建和三维检测仍按不同指标评价。

\paragraph{单侧车辆观测下的米制初始化与对称补全}
2026.09~\cite{2609_08841}。单侧稀疏图像用于车辆资产时，米制尺度和未观测侧外观是两个不同的问题。该方法用视觉基础模型预测米制点云作为高斯中心锚点，通过学习得到的编码和解码模块估计高斯属性，再在高斯空间沿车辆的近似对称面复制可靠原语。

论文在 3DRealCar 用同侧四张图构建资产，并在车辆对侧八个视角评价。与 RECAST 的学习式未见面生成相比，这条路线依靠车辆的双侧对称约束，车门、标识和附件等非对称细节仍需观测。

\paragraph{RECAST：给车辆资产补齐日志中没有看到的视角}
2026.09~\cite{2609_31374}。即使把动态车辆从街景中独立出来，日志重建的车辆仍可能缺少背面或侧面。RECAST 从一张已分割的车辆图像生成可从各方向渲染的高斯，再将资产注册到日志重建的场景。作者构建 RECAR 车辆数据适配单图到三维生成先验，并通过修改车辆和自车轨迹验证闭环渲染。

闭环比较从九个 Waymo 场景选取一个目标车辆并移除其他动态参与者，使用固定的 GTRS-Dense 规划器。结果侧重目标车辆超出日志轨迹时的渲染和规划反应，复杂多参与者交互还需独立评估。

\paragraph{ControlPed：把风险轨迹转化为多视角可渲染行人}
2026.10~\cite{2610_06171}。端到端驾驶评估需要同时控制行人进入道路的位置、时刻和外观。ControlPed 将交通冲突轨迹转成受文本与轨迹约束的人体运动，再驱动高斯化身渲染多相机观测。高斯化身承担三维外观和动画，风险条件来自交通轨迹与姿态生成流程。

评估保留 88 个经检查并渲染的场景，覆盖七个端到端模型；报告的平均 HDScore 从 88.8 降到 47.4。当前设置为开环，因此这一差值应解读为在固定渲染观测下的退化，不能作为交互闭环成绩。


\section*{本章小结}
\addcontentsline{toc}{section}{本章小结}
\label{sec:18_summary}

\begin{futurebox}{本章总结与展望}{ch18_summary}
本章系统梳理了 3D Gaussian Splatting 在自动驾驶场景中的 136 篇论文，覆盖五个核心方向。

\textbf{街景重建与渲染}（48 篇）是最成熟的子方向，已形成从 Street Gaussians 的基础显式表示到 ReconDrive/PointForward 的前馈 4D 重建的完整技术栈。当前热点集中在三个问题：外观解耦（多尺度双边网格、ADM-GS）解决跨轨迹光照变化；轨迹外推（FlexDrive、VEGS）突破训练轨迹限制；多相机兼容（UniGaussian、XYZCylinder）处理鱼眼和异构传感器。

\textbf{仿真与数据生成}（38 篇）正在从离线重建向实时闭环仿真演进。HUGSIM、GSAVS 等平台实现了完整的感知-规划-控制循环。场景编辑能力（RoVES、SIMSplat、HorizonForge）使生成超越原始数据的多样化安全关键场景成为可能。多模态生成（DiffusionHarmonizer、CVD-STORM）和物理仿真（Real2Sim）正在缩小仿真到现实的差距。

\textbf{动态驾驶场景}（25 篇）的核心趋势是从依赖昂贵 3D 标注向自监督和无标注方法演进——$S^3$Gaussian、DeSiRe-GS、DIAL-GS 等工作实现了无需人工标注的高质量动静分解。运动建模从简单的全局变形发展到贝塞尔曲线（B\'ezierGS）和B 样条（AD-GS）等参数化轨迹表示。

\textbf{传感器融合与定位}（13 篇）将 3DGS 从相机渲染扩展到 LiDAR（LiDAR-GS、SplatAD）、雷达（RadarSplat）和鱼眼相机的统一仿真，混合光栅化-光线追踪策略成为多传感器统一渲染的主流方案。

\textbf{3D 占用估计与感知}（12 篇）代表了 3DGS 从重建走向感知的新趋势。高斯基元的稀疏性天然适合替代密集体素进行占用预测（GaussianOcc、ODG），而 3DGS 的实时渲染能力则为弱监督训练提供了新的监督信号来源（GSRender、VoxelSplat）。

展望未来，五个方向值得关注：（1）前馈可泛化重建取代逐场景优化，实现大规模部署；（2）物理感知的闭环仿真融合动力学约束，缩小仿真到现实的领域差距；（3）统一多模态传感器仿真实现完整的数字孪生；（4）高斯表示与下游感知任务的深度耦合，从重建工具演进为感知基础架构；（5）大规模自动化数据流水线利用 3DGS 为自动驾驶算法提供无限训练数据。
\end{futurebox}
